%%%%%%%%%%%%%%%%%%%%%%%%%%%%%%%%%%%%%%%%%%%%%%%%%%%%%%%
% A template for Wiley article submissions.
% Developed by Overleaf. 
%
% Please note that whilst this template provides a 
% preview of the typeset manuscript for submission, it 
% will not necessarily be the final publication layout.
%
% Usage notes:
% The "blind" option will make anonymous all author, affiliation, correspondence and funding information.
% Use "num-refs" option for numerical citation and references style.
% Use "alpha-refs" option for author-year citation and references style.

\documentclass[alpha-refs]{wiley-article}
% \documentclass[blind,alpha-refs]{wiley-article}


%%%%%%%%%%%%%%%%%%%%%%%%%%%%%%%%%%%%%%%%%%%%%%%%%%%%%%%%%
%%%%%%%%%%% JDT:  Annotation Code %%%%%%%%%%%%%%%%%%%%%%%%%%%%%%%%%%
%%%%%%%%%%%%%%%%%%%%%%%%%%%%%%%%%%%%%%%%%%%%%%%%%%%%%%%%%

\usepackage{color}
\usepackage[normalem]{ulem}

%\usepackage{graphicx,amsfonts,amssymb,amsmath,amsthm,mathrsfs}
\usepackage{graphicx,amsfonts,amssymb,amsmath,mathrsfs}

 % Uncomment to display with annotation; comment out otherwise
\newcommand{\add}[1]{\textcolor{blue}{#1}}
\newcommand{\delete}[1]{\textcolor{red}{\sout{#1}}}
\newcommand{\edit}[2]{\textcolor{red}{\sout{#1}} \textcolor{blue}{#2}}
\newcommand{\mnote}[1]{\marginpar{\textcolor{green}{\textbf{#1}}}}

 % Uncomment to display without annotation; comment out otherwise
%\newcommand{\add}[1]{#1}
%\newcommand{\delete}[1]{}
%\newcommand{\edit}[2]{#2}
%\newcommand{\mnote}[1]{}

%%%%%%%%%%%%%%%%%%%%%%%%%%%%%%%%%%%%%%%%%%%%%%%%%%%%%%%%%
%%%%%%%%%%%%%%%%%%%%%%%%%%%%%%%%%%%%%%%%%%%%%%%%%%%%%%%%%
%%%%%%%%%%%%%%%%%%%%%%%%%%%%%%%%%%%%%%%%%%%%%%%%%%%%%%%%%

% Add additional packages here if required

\usepackage{lineno}
\usepackage[breaklinks=true]{hyperref}
\usepackage{amsmath,amssymb}
\usepackage{caption} 
\usepackage{multirow}
\usepackage{xcolor}
\usepackage{breakcites}

\newtheorem{prop}{Proposition}
\newtheorem{thm}{Theorem}
\newtheorem{cor}{Corollary}
\newtheorem{ex}{Example}
\newtheorem{alg}{Algorithm}

% Update article type if known
%\papertype{Original Article}
\papertype{}
% Include section in journal if known, otherwise delete

\title{Cox processes driven by transformed Gaussian processes on linear networks}

% Include full author names and degrees, when required by the journal.
% Use the \authfn to add symbols for additional footnotes and present addresses, if any. Usually start with 1 for notes about author contributions; then continuing with 2 etc if any author has a different present address.
\author[1]{Jesper M\o ller}
\author[1]{Jakob G.\ Rasmussen}

% Include full affiliation details for all authors
\affil[1]{Department of Mathematical Sciences, Aalborg University, Skjernvej 4A, DK-9220~Aalborg~\O, Denmark}

\corraddress{Jesper M\o ller, Department of Mathematical Sciences, Aalborg University, Skjernvej 4A, DK-9220~Aalborg~\O, Denmark}
\corremail{jm@math.aau.dk}

% Include the name of the author that should appear in the running header
\runningauthor{Jesper M\o ller and Jakob G.\ Rasmussen} % and Ninna Vihrs}

\begin{document}

\maketitle

\begin{abstract}
{There is a lack of point process models on linear networks. For an arbitrary linear network, we
 use isotropic covariance functions with respect to the geodesic metric  or the resistance metric 
 to construct new models for isotropic Gaussian processes and hence new models for various Cox processes with isotropic pair correlation functions. In particular we introduce three model classes given by log Gaussian, interrupted, and permanental Cox processes on linear networks,
 and consider for the first time statistical procedures and applications for parametric families of such models.
 Moreover, we construct new simulation algorithms for Gaussian processes  on linear networks and
discuss whether the geodesic metric  or the resistance metric should be used for the kind of Cox processes studied in this paper.}

% Please include a maximum of seven keywords
\textbf{Keywords} -- 
%determinantal point process, generalized shot-noise Cox processes, log Gaussian Cox process, 
estimation,
geodesic distance,
interrupted point process,
isotropic covariance function,
log Gaussian Cox process,
moments, % properties and statistical procedures,
permanental Cox point process, %linear networks,  
resistance metric,
simulation.
\end{abstract}


%\subtitle{Point processes and covariance functions}

%\author[1]{Jesper M\o ller}
%\affil[1]{Department of Mathematical Sciences, Aalborg University}
%%\affil[2]{Department of Statistical Science, Duke University}

%\begin{document}
%\maketitle

%\begin{abstract}
%In this paper, I review certain Cox and other point process models specified by a covariance function...
%\end{abstract}
%
%\keywords{determinantal point process; generalized shot-noise Cox processes; log Gaussian Cox process; permanental point process; linear networks; log Gaussian Cox process; resistance metric.}

\section{Introduction}\label{s:intro}

In Sir David R.\ Cox's highly influential paper `Some Statistical Methods Connected with Series of Events' \citep{Cox55} he invented  
doubly stochastic Poisson processes obtained by a generalization of Poisson processes where the intensity function $\Lambda$ that varies over space or time is a stochastic process. % which are usually not observed. 
These Cox models play nowadays an important role when analysing point patterns in Euclidean spaces or on  spheres \citep[see][and the references therein]{MW04,LawrenceEtAl}.
In particular, Cox processes driven by a transformed Gaussian process (GP) $Y$ or independent copies $Y_1,\ldots,Y_h$ of $Y$
play a major role: A Log Gaussian Cox process (LGCP) has $\Lambda(u)=\exp(Y(u))$ \citep{LGCP,CM}; LGCPs constitute the most widely used subclass of Cox processes. Further, an interrupted Cox process (ICP) is obtained by an independent thinning of a Poisson process, where the selection probability of a point $u$ is given by  $\exp(-\sum_{i=1}^hY_i(u)^2)$ \citep[][]{Stoyan,FredMe}. Moreover, a permanental Cox point process (PCPP) is obtained if $\Lambda(u)=\sum_{i=1}^hY_i(u)^2$ \citep{Macchi:75,MacCM}.

In recent years there has been an increasing interest in analysing point patterns on a linear network $L$, that is, $L$ is a connected set in $\mathbb R^k$ (the real coordinate space of dimension $k$) given by a finite union of bounded, closed, line segments  
which can only overlap at their endpoints, see
\citet{AngEtAl}, \citet{spatstat}, the references therein as well as further references given later in the present paper.
% \citep[see][and the references therein; many further references are given later in the present paper]{AngEtAl,spatstat}. 
Figure~\ref{fig:data} shows two examples of point patterns observed on linear networks, which we will use for illustrative purposes throughout the paper.
% and analyse using the models developed in the paper. 
The first dataset was first analysed in \citet{AngEtAl} and consists of 116 locations of street crimes reported in the period 25 April to 8 May 2002 in an area close to the University of Chicago. 
%This dataset was first analysed in \citet{AngEtAl} who also considered the types of crimes.
% (since the present paper does not include multitype point process models, the types of crimes are ignored and therefore not shown in the figure). 
The second dataset is 
one 
of the six point pattern datasets analysed in \citet{HeidiMe}  and 
consists of 566 locations of spines on a dendrite tree protruding from a neuron \citep[see also][]{RasmussenChristensen}.
 
\begin{figure}
	\begin{center}
		\includegraphics[width=7cm]{"Figs/chicago-data.pdf"}
		\includegraphics[width=7cm]{"Figs/neu-data.pdf"}
	\end{center}
	\caption{Left: The locations of street crimes in a part of Chicago. Right: The locations of spines on a dendrite tree \citep[dataset number five in][]{HeidiMe}. The circle marks the root of the tree, the black lines are a main branch, and the grey lines are side branches.}
	\label{fig:data}
\end{figure}

The contribution of the present paper is the following. We
 use isotropic covariance function models with respect to the geodesic metric $d_{\mathcal G}$ or the resistance metric $d_{\mathcal R}$
 on $L$ as
developed in \citet{AnderesEtAl} as well as new models developed in the present paper in order to construct models for isotropic GPs and hence new models for LGCPs, ICPs, and PCPPs on linear networks with isotropic pair correlation functions (more details follow in the next paragraph). Also we construct new simulation algorithms for GPs and consider for the first time statistical procedures and applications for parametric families of LGCPs, ICPs, and PCPPs on linear networks. Moreover, in continuation of the considerations in
\citet{AnderesEtAl} and \citet{RakshitEtAl}
we discuss which metric $d_{\mathcal G}$ or $d_{\mathcal R}$ we recommend for the kind of Cox processes studied in this paper. Incidentally, we also establish some useful results for the resistance metric.

In brief, the paper consists of two parts, Sections~\ref{s:2}-\ref{s:GP-cov-fct} on our setting for isotropic covariance functions and related GPs, 
and Sections~\ref{s:3}-\ref{s:conclusion} on point processes, in particular Cox processes, including the cases of LGCPs, ICPs, and PCPPs on linear networks and how these models can be used for fitting real data. 
In more detail, the paper is organized as follows. Section~\ref{s:new-setting} considers a general metric space $(S,d)$ in order later to compare covariance functions and point processes for the case $S=L$ with the cases where $S=\mathbb R^k$ or $S=\mathbb S^k$ (the $k$-dimensional unit sphere). When $S=L$ we have in mind that $d$ is either $d_{\mathcal G}$ or $d_{\mathcal R}$, where Section~\ref{s:res-met} gives a summary of results for $d_{\mathcal R}$, including a useful expression for the metric. 
%In order to construct parametric models for GPs on linear networks, 
Furthermore, when $S=L$, Section~\ref{s:iso-cov-fct} studies isotropic covariance functions of the form $c(u,v)=c_0(d(u,v))$ and provides a less technical summary of results from \citet{AnderesEtAl} together with examples of isotropic covariance functions not appearing in that paper. %\citet{AnderesEtAl}. 
Simulation algorithms for GPs on $L$ with an isotropic covariance function are developed in Section~\ref{s:sim-GP}, where the case with $L$ a tree is particularly tractable. Our setting for point processes is given in Section~\ref{s:3.1}, and Section~\ref{s:3.2} introduces first and higher order intensity functions which become useful when we later study Cox process models. In particular, we focus on the pair correlation function and
%, depending on the choice of space, 
the related $K$-function defined in Section~\ref{s:3.2}. As discussed in Section~\ref{s:est-check},  $g$, $K$, and other functional characteristics for point processes become useful for statistical inference. Section~\ref{s:cox-general} treats Cox processes in our general setting, since moment and many other properties of Cox processes driven by transformed Gaussian processes do not depend on which space we consider -- 
the dependence of the metric space $(S,d)$ occurs when we want to specify well-defined parametric models of covariance functions, simulate the underlying GPs, and use statistical inference procedures, which for the case $S=L$ are the topics of Sections~\ref{s:GP-cov-fct} and \ref{s:appl}. Moreover, Section~\ref{s:models} studies the properties of LGCPs, ICPs, and PCPPs models, and in Section~\ref{s:appl} we demonstrate how these models may be fitted to real and simulated data. Finally, Section~\ref{s:conclusion} summarises our findings and discuss some open problems. 

At this point we should stress the importance of considering an isotropic pair correlation function (pcf) $g(u,v)=g_0(d(u,v))$:
\begin{description}
\item[(I)] it is easier to plot $g_0$ than $g$;
\item[(II)] for $S=L$, the $K$-function is only defined when $g$ is isotropic;
\item[(III)] to the best of our knowledge, 
 nonparametric estimators of  $g$ and $K$ have always been derived under the assumption that $g$ is isotropic;
\item[(IV)] moment based estimation procedures such as minimum contrast, composite likelihood, and Palm likelihood methods 
\citep[see][and the references therein]{MW07,MW17} 
become computational convenient if $g$ is isotropic;
 \item[(V)] in particular, for LGCPs, ICPs, and PCPPs, isotropy of $g$ becomes equivalent to isotropy of the covariance function for the underlying GPs;
 \item[(VI)] so far flexible model classes for covariance functions have mainly been developed in the isotropic case 
 \citep[an exception is][cf.\ Section~\ref{s:7.1}]{BSW}.
\end{description}

A substantial part of this work was the development of an R-package \texttt{coxln}, in which the methods developed in the present paper are implemented. This package is available on Github under the author \texttt{gulddahl}. Moreover, we used the R-package \texttt{spatstat} extensively throughout the paper, see \cite{spatstat}. 
 
%Moreover, a linear network may be extended to the more abstract case of a graph with Euclidean edges, i.e., graphs with edges viewed as line or curve segments allowing us to consider points on the graph which are either vertices or points on an edge \citep{AnderesEtAl}, but for ease of presentation this case is avoided in the present paper.   
%
%
%. A widely used subclass consists of log Gaussian Cox processes where the log intensity is a Gaussian process \citep{LGCP}. Hence a log Gaussian Cox process is specified by the mean and covariance functions of the underlying Gaussian process.  
%
%In fact covariance functions are used for constructing many interesting point process models. An example is a determinantal point process (DPP); in brief a DPP is specified by a so-called kernel which is usually a covariance function \citep{Macchi:75,Soshnikov:00,FredEtAl}. While Cox processes are often used for describing clustered point pattern datasets, DPPs become useful for describing repulsiveness. In turn a DPP can be used for constructing a generalized shot-noise Cox process (GSNCP) where the centre process is a DPP with its kernel equal to a covariance function \citep{MV}, thereby obtaining models with repulsion between cluster centres. Moreover, covariance functions can be used for constructing $\alpha$-determinantal point processes \citep{VJ,MacCM}, with e.g.\ $\alpha=-1$ for a DPP and $\alpha=1$ for a permanental point process (PCPP). Interesting, PCPPs are particular cases of Cox processes, cf.\ Section~???. 
%
%The purpose of the present paper is to discuss model constructions for point processes based on covariance functions and defined on a metric space which concretely is either an Euclidean space, a sphere, or a linear network, and the corresponding metric is the usual Euclidean distance, the geodesic metric on the sphere or the linear network, or the resistance metric on the linear network, cf.\ Section~\ref{s:2}. 
%The focus will be on properties of the models which become useful for statistical inference when only a single realization of the point process is observed. In particular, when covariance functions are used for constructing specific point process models, the covariance function is usually assumed to depend only a given metric which depends on the context; I will refer to such covariance functions as isotropic covariance functions, though this terminology may make better sense in cases of point processes defined on the Euclidean space or the sphere than in the case of point processes defined on a linear network. Historically, point process and covariance function models and corresponding statistical methods have first been developed in the Euclidean case, and I discuss to which extent these models and methods do or do not extend to the cases of the sphere and a linear network. In the latter case, I point out that using the resistance metric as compared to the geodesic metric has several advantages. In passing, some useful results for the resistance metric are established. OR A USEFUL RESULT???

%Outline....  The paper is organized as follows.  .... 

%.... Therefore, an important issue is to develop parametric models for isotropic covariance functions as discussed in Section~\ref{s:covfct}. 


%\section{Preliminaries}\label{s:setting}

%This section provides the setting for point processes considered in this paper, specifies characteristics for moment properties of the  processes, and recalls the definition and properties of covariance functions. % Poisson and Cox processes. 

\section{Metric spaces}\label{s:2}

%Section~\ref{new-setting} considers a general metric space $(S,d)$ in order to compare point processes and covariance functions for the case $S=L$ with the cases where $S=\mathbb R^k$, the real coordinate space of dimension $k$, or $S=\mathbb S^k$, the $k$-dimensional unit sphere. When $S=L$ we have in particular in mind that $d$ is either the geodesic metric or the resistance metric. Section~\ref{s:res-met} provides a less technical summary of results for the resistance metric based on \citet{AnderesEtAl} together with a useful expression for the metric.  

\subsection{Setting}\label{s:new-setting} %s:2.1

The state space for the points of the point processes considered in this paper is a metric space $(S,d)$ equipped with a reference measure $\nu$. We focus on the case where $S$ is a linear network and compare with the cases where $S=\mathbb R^k$ is the $k$-dimensional Euclidean space or 
$S=\mathbb S^k=\{u\in \mathbb R^{k+1}\,|\, \|u\|=1\}$ is the $k$-dimensional unit sphere, with $d$ and $\nu$ specified as follows.
\begin{description}
\item[(a)] The case $S=\mathbb R^k$: $d(u,v)=\|u-v\|$ is given by the usual Euclidean distance and $\nu(A)=\int_A\mathrm du$ is Lebesgue measure.
\item[(b)] The case $S=\mathbb S^k$: $d(u,v)=d_\mathcal{G}(u,v)=\arccos(u\cdot v)$ is the geodesic (orthodromic, great-circle, or shortest path) metric, and $\nu=\nu_k$ is $k$-dimensional Hausdorff/surface measure \citep[see e.g.][Chapter~1]{MR3060033}.
%($k$-dimensional Hausdorff on $\mathbb S^k$, see Appendix~\ref{a:A1}).
 Here, $u\cdot v$ is the usual inner product of $u$ and $v$. 
 \item[(c)] The case $S=L=\cup_{i=1}^m L_i$ of a linear network: We assume $m<\infty$, each $L_i\subset\mathbb R^k$ is a closed line segment of  length $l_i\in (0,\infty)$, $L_i\cap L_j$ is either empty or an endpoint of both $L_i$ and $L_j$ whenever $i\not=j$, and $L$ is a path-connected set. Furthermore, $d$ is a `natural' metric and $\nu(A)=\int_A\mathrm d_L(u)$ is arc length measure. We let $|L|=\sum_{i=1}^m l_i=\int_L\mathrm d_L(u)$ denote the length of the linear network.
\end{description} 
The remainder of this section are remarks 
to the case (c). 

For disconnected linear networks, definitions and results may be applied
separately to each connected component of the network if we consider independent Gaussian processes on the connected components and independent point processes on the connected components.

Our definition in (c) may be extended to the more abstract case of a graph with Euclidean edges  \citep{AnderesEtAl} but 
%, i.e., graphs with edges viewed as line or curve segments allowing us to consider points on the graph which are either vertices or points on an edge \citep{AnderesEtAl}. In the present paper 
we avoid this generalization %partly 
for ease of presentation and %partly 
since statistical methods have so far only been developed for the case (c). 

For each line segment $L_i$, there are two possible arc length parametrisations. We assume one is chosen and given by $u_i(t)=(1-t/l_i)a_i+(t/l_i)b_i$, $t\in[0,l_i]$, where $a_i$ and $b_i$ are the endpoints of $L_i$. The definitions and results in this paper will not depend on this choice, including when calculating arc length measure restricted to $L_i$: For Borel sets $A\subseteq L_i$, $\int_A \mathrm d_L(u)=\int_0^{l_i}1(u_i(t)\in A)\,\mathrm dt$ where $1(\cdot)$ denotes the indicator function. 
 
Let $V$ denote the set of endpoints of $L_1,\ldots,L_m$ and consider the graph with vertex set $V$ and edge set $E$ given by $L_1,\ldots,L_m$. Thus, two distinct vertices $u,v\in V$ form an edge if and only if $\{u,v\}=\{a_j,b_j\}$ for some $j\in\{1,\ldots,m\}$, in which case we write $u\sim v$. 

We have %mainly 
two cases of natural metrics in mind, namely when $d$ is the geodesic metric $d_\mathcal{G}$ or the resistance metric $d_\mathcal{R}$ defined as follows.
Let $u,v\in L$. Then $d_\mathcal{G}(u,v)=\min\int_{p_{uv}}\mathrm d_L(w)$ where the minimum is over all paths $p_{uv}\subseteq L$ connecting $u$ and $v$.
%as developed in \citet[][]{AnderesEtAl} when extending the original definition in \citet{KleinRandic} for $V$ to $L$:   
When defining $d_\mathcal{R}(u,v)$, without loss of generality assume that $u,v\in V$, since if e.g.\ $u\in L_j\setminus V$, we may split $L_j$ into the two line segments with endpoints $\{a_j,v\}$ and $\{v,b_j\}$, and then consider a new graph with vertex set $V\cup\{u\}$ and edge set $E\cup\{\{a_j,v\},\{v,b_j\}\}$ (this follows from (B) in Theorem~\ref{t:res-m} below). Viewing the graph $(V,E)$ as an electrical network with resistor $l_i$ at edge $L_i$, $i=1,\ldots,m$, then
$d_\mathcal{R}(u,v)$ is the effective resistance between $u$ and $v$ as obtained by Kirkhoff's laws. 
%Appendix~\ref{a:A2} 
In Section~\ref{s:res-met} below we provide the detailed definition of $d_\mathcal{R}$.

Indeed there are other interesting metric when $S=L$, including the least-cost metric \citep{RakshitEtAl}, %Rakshit, Nair, Baddeley (2017), Spat. Stat.
but to the best of our knowledge parametric models for isotropic covariance functions $c(u,v)=c_0(d(u,v))$ %(which is one topic of our main interest) 
have so far only been developed when $d=d_\mathcal{G}$, $d=d_\mathcal{R}$, or $d$ is given by the usual Euclidean distance. However, Euclidean distance is usually not a natural metric on a linear network.

\subsection{The resistance metric}\label{s:res-met} %s:2.4

This section defines the resistance metric $d_\mathcal{R}$ for a graph with Euclidean edges \citep{AnderesEtAl} in the special case of
 a linear network $L=\cup_{i=1}^m L_i$ as given in case (c) in Section~\ref{s:new-setting}. The section also summarises some properties of $d_\mathcal{R}$ and compares with $d_\mathcal{G}$. 
 
 Consider the graph $G=(V,E)$ and its relation $\sim$ as defined above, and denote $d_V$  the classic (effective) resistance metric $d_V$ on $V$ \citep{KleinRandic}. % with relation $u\sim v$. 
 Since $d_\mathcal{R}$ is an extension of $d_V$ to $L$, we start by recalling the definition of $d_V$ using a notation as follows.
Let $u_0\in V$ be an arbitrarily chosen vertex called the origin.
For any $u,v\in V$, define the so-called conductance function by
\[{\mathrm{con}}(u,v)=\begin{cases*}
                    1/\|u-v\| & \text{if } $u\sim v$  \\
                     0 & \text{otherwise}
                 \end{cases*}
                 \] 
and define a matrix $\Delta$ with rows and columns indexed by $V$ so that its entry $(u,v)$ is given by
\[\Delta(u,v)=\begin{cases*}
                    1+c(u) & \text{if } $u=v=u_0$  \\
                     c(u) & \text{if } $u=v\not=u_0$\\
                     -{\mathrm{con}}(u,v) & \text{otherwise}
                 \end{cases*}
                 \] 
where $c(u)=\sum_{w\in V:\,w\sim u}{\mathrm{con}}(u,w)$ is the sum of the conductances associated to the edges incident to vertex $u$.
In fact $\Delta$ is symmetric and strictly positive definite, and it is similar to %what would be called 
the `Laplacian matrix' obtained 
%from electrical network theory \citep[
when viewing $G$ as an electrical network over the nodes with resistors given by the length of each line segment \citep[see e.g.][]{Kigami,JorgensenPearse} 
except that $\Delta$ has the additional $1$ added at entry $(u_0,u_0)$ (this makes $\Delta$ invertible). 

Now, let $B_0$ be a zero mean Gaussian vector indexed by $V$ and having covariance matrix $\Sigma=\Delta^{-1}$. Then
the resistance metric on $V$ is the variogram
\begin{equation}
\label{e:d-V}
d_V(u,v)=\mathbb{V}{\mathrm{ar}}(B_0(u)-B_0(v))=\Sigma(u,u)+\Sigma(v,v)-2\Sigma(u,v)\quad\mbox{for } u,v\in V.
\end{equation}
Extend $B_0$ by linear interpolation to a zero mean Gaussian process (GP) $Z_0$ on $L$ so that %$Z_0(u)=B_0(u)$ if $u\in V$ and
\[Z_0(u)=\frac{\|u-b_i\|}{l_i}B_0(a_i)+\frac{\|u-a_i\|}{l_i}B_0(b_i)\quad\mbox{for }u\in L_i.\]
For $i=1,\ldots,m$, define a zero mean Brownian bridge $B_i$ on $L_i$ so that  %$B_i(a_i)=B_i(b_i)=0$ and
\[\mathbb{C}{\mathrm{ov}}(B_i(u),B_i(v))=\min\{\|u-a_i\|\|v-b_i\|,\|v-a_i\|\|u-b_i\|\}/l_i\quad\mbox{for $u,v\in L_i$,}
%\begin{cases*}
%\|u-a_i\|\|v-b_i\|/l_i & if $\|u-a_i\|\le\|v-a_i\|$ \\
%\|v-a_i\|\|u-b_i\|/l_i & if $\|v-a_i\|\le\|u-a_i\|$
%\end{cases*}
%%\left[\min\{d_i(u),d_i(v)\}-d_i(u)d_i(v)\right]l_i,\quad u,v\in L_i.
\]
%as follows. 
%First, $Z_0(u)=B_0(u)$ if $u\in V$. Second, 
% $Z_0$ is defined on each line segment $L_i$ by linear interpolation so that 
%% if 
%%\[d_i(u)= 
%%\begin{cases*}
%%                    \|u-a_i\|/l_i & for  $u\in L_i\setminus\{a_i,b_i\}$  \\
%%                     0 & for $u\in\{a_i,b_i\}$
%%                 \end{cases*}
%%                 \]
%%is the distance of $u\in L_i$ as a proportion of the length
%%of $L_i$, then
%%\[Z_0(u)=\left(1-d_i(u)\right)Z_0(a_i)+d_i(u)Z_0(b_i).\]
%\[Z_0(u)=\frac{\|u-b_i\|}{l_i}Z_0(a_i)+\frac{\|u-a_i\|}{l_i}Z_0(b_i)\quad\mbox{if }u\in L_i.\]
and define
\[Z_i(u)=\begin{cases*}
                    B_i(u) & \text{for } $u\in L_i$  \\
                     0 & \text{for } $u\in L\setminus L_i$.
                 \end{cases*}
                 \]
%Finally, % assuming $B_0,B_1,\ldots,B_m$ are independent, 
%define 
%\begin{equation*}
%    Z(u)=\sum_{i=0}^m\!Z_i(u),\quad u\in L.
%\end{equation*}
%Thus $Z$ is a sum of $m+1$ independent, zero mean, GRFs. 
Finally, the resistance metric on $L$ is defined by %is defined by
 \begin{equation}\label{e:def-dR}
 d_{\mathcal{R}}(u,v) %={\mathrm{var}}(Z(u)-Z(v))
 =\sum_{i=0}^m \mathbb{V}{\mathrm{ar}}(Z_i(u)-Z_i(v))\quad\mbox{for } u,v\in L.
 \end{equation}
 
For the following theorem, which follows from \citet[][Propositions~2-4]{AnderesEtAl}, we use a terminology as follows. 
%Consider again the graph defined by $L_1,\ldots,L_m$. 
A closed line segment in $\mathbb R^k$ with endpoints $a$ and $b$ is denoted $[a,b]=\{at+b(1-t)\,|\, 0\le t\le 1\}$.
A path is a subset of $L$ of the form $[u,v_1]\cup [v_1,v_2]\cdots\cup [v_{i-1},v_{i}]\cup[v_{i},v]$ where $u,v\in L$, $v_1,\ldots,v_i\in V$ are vertices, $i\ge0$ is an integer, and we interpret $[v_1,v_2]\cdots\cup [v_{i-1},v_{i}]$ as the empty set if $i=0$. 
% path connecting $u\in L_j$ and $v\in L_i$ is a set of the form   
If all vertices in $G$ are of order two, we say that $L$ is a loop (since $L$ is isomorphic to a circle). If there is no loop, we say that $L$ is a tree.  

\begin{thm}\label{t:res-m} We have the following properties of $d_{\mathcal{G}}$ and $d_{\mathcal{R}}$.
%for the resistance and geodesic metrics on a linear network $L$. 
\begin{description}
\item[(A)] The definition \eqref{e:def-dR} of $d_{\mathcal{R}}$ does not depend on the choice of origin $u_0\in V$.
\item[(B)] Both $d_{\mathcal{G}}$ and $d_{\mathcal{R}}$ are metrics on $L$, and their definitions are invariant to splitting a line segment $L_i$ into two line segments.
\item[(C)] For every $u,v\in L$, $d_\mathcal{G}(u,v)\ge d_\mathcal{R}(u,v)$, with equality if and only if there is only one path connecting $u$ and $v$. In particular,
$d_\mathcal{G}= d_\mathcal{R}$ if and only if $L$ is a tree.
\item[(D)] If $G$ is a loop, then $d_{\mathcal{R}}(u,v)=d_{\mathcal{G}}(u,v)-d_{\mathcal{G}}(u,v)^2/\sum_{i=1}^ml_i$.
\end{description} 
\end{thm}

Regarding the inequality in (C) above, %$d_{\mathcal R}\le d_{\mathcal G}$, 
considering e.g.\ the Chicago street network in Figure~\ref{fig:data}, 
 $\max d_{\mathcal R}\approx 675$  
%the maximal distance with respect to the resistance metric is about 675 
feet is much smaller than $\max d_{\mathcal G}\approx 2031$ feet. It is also smaller than the side length of a square surrounding the network which is a little less than 1000 feet.

%and with respect to the geodesic metric is about 2031 feet (the side length of a square surrounding the network is a little less than 1000 feet).

The following proposition and the remarks below show that $d_{\mathcal{R}}(u,v)$ is nicely behaving.
 
\begin{prop}\label{p:calc}  
For any $u\in L_j$ and $v\in L_i$, %so that $i=j$ if $L_i\cap L_j\not=\emptyset$ 
let 
\[s=\|u-a_j\|,\quad t=\|v-a_i\|,\quad A_i=d_V(a_i,b_i)/l_i^2-1/ l_i.\]
Then $A_i\le0$ with equality if and only if $L_i$ is the only path connecting $a_i$ and $b_i$, and $d_{\mathcal{R}}(u,v)$ satisfies the following. 
\begin{description}
\item[(A)] If $i=j$ then
\begin{equation}\label{e:d-R1}
d_{\mathcal{R}}(u,v) 
=\begin{cases}
A_i(t-s)^2+t-s & \text{if }t\ge s,\\
A_i(s-t)^2+s-t & \text{if }t\le s,
\end{cases}
\end{equation}
so $d_{\mathcal{R}}(u,v)$ considered as a function of $t$ is linear (the case $A_i=0$) or quadratic (the case $A_i<0$) on each of the intervals $[0,s]$ and $[s,l_i]$, continuous on $[0,l_i]$, and differentiable on $[0,l_i]\setminus\{s\}$.
\item[(B)] If $i\not =j$ then
\begin{equation}\label{e:d-R2}
d_{\mathcal{R}}(u,v) 
=A_it^2+B_{ij}(s)t+C_{ij}(s)
\end{equation}
where
\[B_{ij}(s)=1-\frac{2}{l_i}\left[\Sigma(a_i,a_i)-\Sigma(a_i,b_i)-\frac{l_j-s}{l_j}\Sigma(a_j,a_i)+\frac{l_j-s}{l_j}\Sigma(a_j,b_i)-\frac{s}{l_j}\Sigma(b_j,a_i)+\frac{s}{l_j}\Sigma(b_j,b_i)\right]\]
and
\[C_{ij}(s)=\frac{(l_j-s)^2}{l_j^2}\Sigma(a_j,a_j)+\frac{s^2}{l_j^2}\Sigma(b_j,b_j)+2\frac{s(l_j-s)}{l_j^2}\Sigma(a_j,b_j)+\Sigma(a_i,a_i)-2\frac{l_j-s}{l_j}\Sigma(a_j,a_i)-2\frac{s}{l_j}\Sigma(b_j,a_i)+\frac{s(l_j-s)}{l_j},\]
so $d_{\mathcal{R}}(u,v)$ is a linear or quadratic concave function of $t\in[0,l_i]$.
\item[(C)] If $i\not =j$ then
\begin{equation}
d_{\mathcal{R}}(u,v)\ge\min\{d_{\mathcal{R}}(u,a_i),d_{\mathcal{R}}(u,b_i)\}.\label{e:min}
\end{equation}
\end{description}
\end{prop}
\begin{proof} Since $l_i=d_{\mathcal G}(a_i,b_i)=d_V(a_i,b_i)$, Theorem~\ref{t:res-m}(C) gives that $A_i\le0$ with equality if and only if $L_i$ is the only path connecting $a_i$ and $b_i$. 
From \eqref{e:d-V} and \eqref{e:def-dR} we obtain \eqref{e:d-R1} and \eqref{e:d-R2} by a straightforward calculation, and thereby we easily see that $d_{\mathcal{R}}(u,v)$ as a function of $t$ behaves as stated in (A) and (B). 
Finally, since $d_{\mathcal{R}}(u,v)$ is a concave function of $t\in[0,l_i]$ in the case $i\not=j$, 
the inequality \eqref{e:min} follows.
\end{proof}
  
%\begin{align}
%d_{\mathcal{R}}(u,v) %{\mathrm{var}}(Z_0(u)-Z_0(v))
%=\,&
%\frac{\|u-b_j\|^2}{l_j^2}\Sigma(a_j,a_j)+\frac{\|u-a_j\|^2}{l_j^2}\Sigma(b_j,b_j)+2\frac{\|u-b_j\|\|u-a_j\|}{l_j^2}\Sigma(a_j,b_j)\nonumber\\
%&+\frac{\|v-b_i\|^2}{l_i^2}\Sigma(a_i,a_i)+\frac{\|v-a_i\|^2}{l_i^2}\Sigma(b_i,b_i)+2\frac{\|v-b_i\|\|v-a_i\|}{l_i^2}\Sigma(a_i,b_i)\nonumber\\
%&-2\frac{\|u-b_j\|\|v-b_i\|}{l_jl_i}\Sigma(a_j,a_i)
%-2\frac{\|u-b_j\|\|v-a_i\|}{l_jl_i}\Sigma(a_j,b_i)\nonumber\\
%&-2\frac{\|u-a_j\|\|v-b_i\|}{l_jl_i}\Sigma(b_j,a_i)
%-2\frac{\|u-a_j\|\|v-a_i\|}{l_jl_i}\Sigma(b_j,b_i)\nonumber\\
%&+\frac{\|v-a_i\|\|v-b_i\|}{l_i}+\frac{\|u-a_j\|\|u-b_j\|}{l_j}\nonumber\\
%&-2\frac{1(u\in L_i)}{l_i}\min\{\|u-a_i\|\|v-b_i\|,\|v-a_i\|\|u-b_i\|\}\label{e:d-R} \\
%\ge\,&\min\{d_{\mathcal{R}}(u,a_i),d_{\mathcal{R}}(u,b_i)\}\label{e:min}
%\end{align}
%% \citet[][Proposition 4]{AnderesEtAl} states that $d_{\mathcal{R}}\le d_{\mathcal{G}}$
%%with equality if and only if $G$ is a tree. For example, $d_{\mathcal{R}}(u,v)=d_{\mathcal{G}}(u,v)-d_{\mathcal{G}}(u,v)^2/\sum_{i=1}^ml_i$ if $G$ is a cycle.
%\end{prop}
%\begin{proof}
%From \eqref{e:def-dR} we obtain \eqref{e:d-R} by a straightforward calculation. The equality \eqref{e:min} is verified in the proof of Proposition~\ref{p:weight}.
%\end{proof}

A few final remarks are in order. 

It follows from \eqref{e:d-R1} and \eqref{e:d-R2} that once $\Sigma$ has been calculated, $d_{\mathcal{R}}(u,v)$ can be quickly calculated. 
Moreover, if $V$ is not too large or if $\Delta$ is sparse, there are quick methods for obtaining $\Sigma=\Delta^{-1}$ \citep[see e.g.][]{rue:held}. For example, for
the Chicago street network in Figure~\ref{fig:data} there are 338 vertices and $\Delta$ is sparse since it contains around 1.2\% non-zero entries. In our implementation the inversion of $\Delta$ took less than a 0.1 second using standard methods for inversion, so here sparse methods are not necessary, but for larger networks sparse methods may become relevant.

%or example, for the Chicago street network in Figure~\ref{fig:data} there are {\color{red} ??? vertices and ....}. 
%We have $d_{\mathcal{R}}(v,a_i)=A_it^2+t$, cf.\ 
% \eqref{e:d-R1} with $s=0$, and $d_{\mathcal{R}}(u,a_i)=C_{ij}(s)$, cf.\ \eqref{e:d-R2} with $t=0$, so we can rewrite  \eqref{e:d-R2} as
%\[d_{\mathcal{R}}(u,v)=d_{\mathcal{R}}(u,a_i)+d_{\mathcal{R}}(v,a_i)+\left[B_{ij}(s)-1\right]t.\]
%This together with the triangle inequality show that $B_{ij}(s)\le1$.

The inequality \eqref{e:min} becomes useful when searching for point pairs $u,v\in L$ with $d_{\mathcal{R}}(u,v)\le r$ and $v\in L_i$, since we need only to consider the cases where $d_{\mathcal{R}}(u,a_i)\le r$ or $d_{\mathcal{R}}(u,b_i)\le r$.  

\section{Gaussian processes and isotropic covariance functions}\label{s:GP-cov-fct}

Let $(S,d)$ be a metric space as in Section~\ref{s:new-setting} and recall that a function $c:S\times S\mapsto\mathbb R$ is positive definite (sometimes called positive semi-definite) if  $\sum_{j,\ell=1}^n a_j{a_\ell}c(u_j,u_\ell)\ge0$ for all $a_1,\ldots,a_n\in\mathbb R$, all pairwise distinct $u_1,\ldots,u_n\in S$, and all $n=1,2,\ldots$. Let
$Y=\{Y(u)\,|\,u\in S\}$ be a Gaussian process (GP) % on a metric space $(S,d)$ as given in Section~\ref{s:new-setting}, 
where each $Y(u)$ is a real-valued random variable. The distribution of $Y$ is specified by the mean function $\mu(u)=\mathbb EY(u)$ and the covariance function 
\[c(u,v)=\mathbb C{\mathrm{ov}}(Y(u),Y(v))=\mathbb E[Y(u) %\overline
{Y(v)}]-\mu(u) %\overline
{\mu(v)}.\] 
The necessary and sufficient condition for a well-defined GP in terms of such two functions $\mu$ and $c$ is just that $c$ is symmetric %Hermitian 
and positive definite.

%that for $n=1,2,\ldots$ and every $u_1,\ldots,u_n$, $(Y(u_1),\ldots,Y(u_n))$ follows an $n$-dimensional %complex 
%normal distribution given by 
%%Here, for $x,y\in\mathbb R$, we let $\overline{x+iy}=x-iy$ denote complex conjugation where $i^2=-1$. Some further remarks are in order.
%The necessary and sufficient condition for a well-defined GP is the property of the covariance function that it is symmetric %Hermitian 
%and positive definite.

% (sometimes called positive semi-definite). Recall that a function $c:S\times S\mapsto\mathbb C$ is Hermitian if $c(u,v)=\overline{c(v,u)}$ for all $u,v\in S$, and $c$ is positive definite if $\sum_{j,\ell=1}^n z_j\overline{z_\ell}c(u_j,u_\ell)\ge0$ for all $z_1,\ldots,z_n\in\mathbb C$, all pairwise distinct $u_1,\ldots,u_n\in S$, and $n=1,2,\ldots$. We say that $c$ is strictly positive definite if    
%$\sum_{j,\ell=1}^n z_j\overline{z_\ell}c(u_j,u_\ell)=0$ implies that $z_1=\ldots=z_n=0$.
%
%Usually in this paper $Y$ is real. Then $c$ is real and symmetric, and positive definiteness means that $\sum_{j,\ell=1}^n a_j{a_\ell}c(u_j,u_\ell)\ge0$ for all $a_1,\ldots,a_n\in\mathbb R$, all pairwise distinct $u_1,\ldots,u_n\in S$, and $n=1,2,\ldots$. The case where $Y$ is complex becomes only relevant in Section~\ref... (the case of permanental point processes) and in Section~\ref... (the case of determinantal point processes). 

\subsection{Isotropic covariance functions}\label{s:iso-cov-fct}

We are in particular interested in isotropic covariance functions %$c$ meaning that $c$ is of the form 
$c(u,v)=c_0(d(u,v))$ for all $u,v\in S$, where 
 with some abuse of terminology we also call $c_0$ a covariance function. So $c_0$ is required to be positive definite, that is, $\sum_{j,\ell=1}^n a_j{a_\ell}c_0(d(u_j,u_\ell))\ge0$ for all $a_1,\ldots,a_n\in\mathbb R$, all pairwise distinct $u_1,\ldots,u_n\in S$, and all $n=1,2,\ldots$. 
 % the necessary and sufficient condition for a well-defined GP . It is easily seen that $c_0$ is a real function, and hence the necessary and sufficient condition for $c_0$ is that it is positive definite in the sense that $\sum_{j,\ell=1}^n a_j{a_\ell}c_0(d(u_j,u_\ell))\ge0$ for all $a_1,\ldots,a_n\in\mathbb R$, all pairwise distinct $u_1,\ldots,u_n\in S$, and $n=1,2,\ldots$. 
 %In Sections~\ref{s:classical} and \ref{s:L-cov} below 
 Henceforth, we assume that the variance $\sigma^2=c_0(0)$ is strictly positive. Below in Sections~\ref{s:classical} and \ref{s:L-cov} we pay attention to the correlation function $r_0(t)=c_0(t)/\sigma^2$ in connection to the cases (a)--(c) in Section~\ref{s:new-setting}. 
 %Thus, $r_0:[0,\infty)\mapsto[0,1]$ is a positive definite function with $r_0(0)=1$; 

\subsubsection{The classical cases of Euclidean spaces and spheres}\label{s:classical}

Let $S=\mathbb R^k$ and $d(u,v)=\|u-v\|$, cf.\ case (a) in Section~\ref{s:new-setting}. %Then $c$ is said to be stationary if $c(u,v)$ depends only on the lag $u-v$, where $C$ is also called a covariance function. By the Bocher-Khinchin theorem, for a continuous function $C:\mathbb R^k\mapsto\mathbb C$, $c(u,v)=C(u-v)$ is positive definite if and only if $C$ is the Fourier transform of a finite Borel measure on $\mathbb R^k$.   
%Clearly, $c$ is stationary if it is isotropic or just geometric anisotropic, meaning that $c(u,v)=c_0(\|D\mathcal O(u-v)\|)$ where $D$ is a $k\times k$ diagonal matrix with positive diagonal elements and $\mathcal O$ is a $k\times k$ rotation matrix. 
\citet{Schoenberg38} and \citet{Gneiting01} provided detailed studies of continuous isotropic covariance functions for this case. 
Table~\ref{tab:covariances-sphere} shows
examples of popular parametric models for continuous isotropic correlation functions which are well-defined for every dimension $k=1,2,\ldots$. 

Instead let $S=\mathbb S^k$ and $d(u,v)=d_\mathcal{G}(u,v)$, cf.\ case (b) in Section~\ref{s:new-setting}. \citet{Schoenberg} gave a complete characterization of continuous isotropic covariance functions in this case; see also \citet{gneiting2013strictly}. In particular, it was proven that $r_0$ defines a continuous correlation function for any dimension $k=1,2,\ldots$ if and only if $r_0(t)=\sum_{\ell=0}^\infty\beta_\ell\cos^\ell(t)$ where the sequence $\beta_0,\beta_1,\ldots$ is a probability mass function. Examples of such correlation functions are seen in Table~\ref{tab:covariances-sphere}. Compared to the case $S=\mathbb R^k$, the ranges of shape and smoothness parameters are more restrictive.


\begin{table}
    \centering
    \begin{tabular}{l||l|l}
        Model & Correlation function $r_0(t)$ & Range of shape and smoothness parameters\\
        \hline
        \hline
        Powered exponential & $\exp\left(-t^\alpha/\phi\right)$ & $\alpha\in(0,2]$ if $S=\mathbb R^k$; $\alpha\in(0,1]$ if $S\in\{\mathbb S^k,L\}$\\
        Mat{\'e}rn & $\frac{2^{1-\alpha}}{\Gamma(\alpha)}\left( \sqrt{2\alpha}\frac{t}{\phi} \right)^{\alpha} K_{\alpha}\left(\sqrt{2\alpha}\frac{t}{\phi}\right)$ &  $\alpha>0$  if $S=\mathbb R^k$; $0< \alpha\leq \frac{1}{2}$ if $S\in\{\mathbb S^k,L\}$\\
        Generalized Cauchy & $(1 + (\frac{t}{\phi})^{\alpha})^{-\tau/\alpha}$ & $\tau > 0$; $\alpha \in (0,2]$ if $S=\mathbb R^k$; $\alpha \in (0,1]$ if $S\in\{\mathbb S^k,L\}$\\
        Dagum & $1 - ((\frac{t}{\phi})^{\tau}/(1+(\frac{t}{\phi})^{\tau}))^{\frac{\alpha}{\tau}}$ & $\tau \in (0,2]$ and $\alpha \in (0,\tau)$ if $S=\mathbb R^k$; \\
        & & $\tau \in (0,1]$ and $\alpha \in (0,1]$ if $S\in\{\mathbb S^k,L\}$
    \end{tabular}
    \caption{Four parametric models for an isotropic correlation function $r_0(r)$. Here, $\Gamma$ is the gamma function, $K_{\nu}$ is the modified Bessel function of the second kind, $\phi$ is a scale parameter, $\tau$ is a 
shape parameter, and $\alpha$ is a smooth\-ness parameter. The correlation functions are well-defined at all scales $\phi>0$ but the range of shape and smoothness parameters depend on the model and the space $S$. For $S\in\{\mathbb R^k,\mathbb S^k\}$, the correlation functions are well-defined for every dimension $k=1,2,\ldots$. For $S=L$, conditions on $L$ may be needed if distance is not measured by the resistance but the geodesic metric, see Section~\ref{s:L-cov}.
    % and $t_{+} := \max\{t,0\}$ for $t\in\mathbb R$. For the powered exponential, Mat{\'e}rn,  generalized Cauchy, Dagum, multiquadric, and %sine power models, $k\in \{1,2,\ldots\}$, whilst for the spherical, Askey, $C^2$-Wendland, and $C^4$-Wendland models, $k\in\{1,2,3\}$.
    } 
    \label{tab:covariances-sphere}
\end{table}

\subsubsection{The case of linear networks}\label{s:L-cov}

Suppose $S=L$ is a linear network, cf.\ case (c) in Section~\ref{s:new-setting}, where we let either $d=d_{\mathcal G}$ or $d=d_{\mathcal R}$.  
Many of the commonly used isotropic correlation functions, including those in Table~\ref{tab:covariances-sphere}, 
are valid with respect to the resistance metric but not always with respect to the geodesic metric. The reason for this is discussed 
%at the end of this section and is based on various results summarized 
in this section.  

Typically (including all of our examples), $r_0$ will %also 
 be a completely monotone function. Recall that 
a function $f : [0,\infty)\mapsto\mathbb R$ is completely monotonic if it is non-negative and continuous on $[0,\infty)$ and for $j=0,1,\ldots$ and all $u>0$, the $j$-th derivative $f^{(j)}(u)$ exists and satisfies $(-1)^jf^{(j)}(u)\ge0$. By Bernstein's theorem, $f$ is completely monotone if and only if it is the Laplace transform of a non-negative finite measure on $[0,\infty)$, meaning that for every $t\ge0$,
\begin{equation}\label{e:bernstein}
f(t)=f(0)\int\exp(-st)\,\mathrm dF(s)
\end{equation}
where $F$ is a cumulative distribution function with $F(s)=0$ for $s<0$. We refer to $F$ as the Bernstein CDF corresponding to $f$. 

Thus, any non-negative valued distribution with a known Laplace transform can be used to produce a completely monotone function. This fact is used in the following example. 

\begin{ex}\label{ex:0} 
The following functions $f_1,f_2,f_3$ are completely monotone functions with $f_1(0)=f_2(0)=f_3(0)=1$ and they have corresponding Bernstein CDFs $F_1,F_2,F_3$ defined as follows.
For $\tau>0$, $\phi>0$, and $t\ge0$,
\begin{equation}\label{e:f_1}
f_1(t)=(1+t/\phi)^{-\tau},\quad F_1\sim\Gamma(\tau,\phi),
\end{equation}
where $\Gamma(\tau,\phi)$ denotes the gamma distribution with shape parameter $\tau$ and rate parameter $\phi$, and
%\begin{equation}
%f_2(t)=\frac{2}{\Gamma(\tau)}(t/\phi)^{\tau/2}K_\alpha(2\sqrt{t/\phi}),\quad F_2\sim\Gamma^{-1}(\tau,\phi).
%\end{equation}
\begin{equation}\label{e:f_2}
f_2(t)={2(t\phi)^{\tau/2}}K_\tau(2\sqrt{t\phi})/{\Gamma(\tau)},\quad F_2\sim\Gamma^{-1}(\tau,\phi),
\end{equation}
where $\Gamma^{-1}(\tau,\phi)$ denotes the inverse gamma distribution with shape parameter $\tau$ and scale parameter $\phi$.
Moreover, for $\psi>0$, $\chi>0$, $\lambda\in\mathbb R$, and $t\ge0$,
%\begin{equation}\label{e:f_3}
%f_3(t)= (1+2t/\psi)^{\lambda/2}K_\lambda(\sqrt{\psi\chi})/K_\lambda(\sqrt{(2t+\psi)\chi})
%\end{equation}
\begin{equation}\label{e:f_3}
f_3(t)= (1+2t/\psi)^{-\lambda/2}{K_\lambda(\sqrt{(2t+\psi)\chi})}/{K_\lambda(\sqrt{\psi\chi})}
\end{equation}
and $F_3$ is the CDF for a generalized inverse Gaussian distribution with probability density function
%\[
%2\frac{K_\lambda(\sqrt{\psi\chi})}{(\psi/\chi)^{\lambda/2}}  s^{\lambda-1}\exp(-s \psi/2 -\chi/(2s)),\quad s\ge 0.
%\]
\[
\frac{(\psi/\chi)^{\lambda/2}} {2K_\lambda(\sqrt{\psi\chi})} s^{\lambda-1}\exp(-s \psi/2 -\chi/(2s)),\quad s\ge 0.
\]
\end{ex}

In the next theorem, which summarises Theorems 1 and 2 in \citet[][]{AnderesEtAl}, we need the following definition. We say that $L$ is a 1-sum of $\mathcal L_1=L_1\cup\ldots\cup L_j$ and $\mathcal L_2=L_{j+1}\cup\ldots\cup L_m$ if $\mathcal L_1$ and $\mathcal L_2$ are (connected) linear networks where $1\le j<m$, $\mathcal L_1\cap\mathcal L_2=\{u_0\}$ consists of a single point $u_0$, and 
\[d(u,v)=d(u,u_0)+d(v,u_0)\quad\mbox{whenever $u\in\mathcal L_1$ and $v\in\mathcal L_2$.}\]
This property is possible if $d=d_{\mathcal G}$ or $d=d_{\mathcal R}$ but unless $L$ is a straight line segment it is impossible if $d$ is the given by the usual Euclidean distance. Using induction we say for $n=3,4,\ldots$ that $L=\mathcal L_1\cup\ldots\cup\mathcal L_n$ is a 1-sum of  $\mathcal L_1,\ldots,\mathcal L_n$ if $L$ is a 1-sum of $\mathcal L_1\cup\ldots\cup\mathcal L_{n-1}$ and $\mathcal L_n$. 

%say that $L$ is a finite 1-sum of trees and cycles if the set $\{L_1,\ldots,L_m\}$ can be divided into disjoint sets of line segments whose unions are trees or cycles which pairwise share at most one vertex/endpoint. 

\begin{thm}\label{t:1}
Let $f : [0, \infty)\mapsto\mathbb R$ be a completely monotone and non-constant function. Then
 $f(d_{\mathcal R}(u, v))$ is strictly
positive definite over $(u, v)\in L\times L$. Moreover, if $L$ is a 1-sum of trees and loops, then $f(d_{\mathcal G}(u, v))$ is strictly
positive definite over $(u, v)\in L\times L$. However,
if there are three distinct paths between two points on $L$, then there exists a constant $\phi>0$ so that  $\exp(-d_{\mathcal G}(u, v)/\phi)$ is not positive definite over $(u, v)\in L\times L$. 
\end{thm} 

In Table~\ref{tab:covariances-sphere}, for
 each model, $r_0$ is completely monotone for the ranges of the parameters \citep[cf.\ the comments to Theorem~1 in][]{AnderesEtAl}. 
So for each example of $r_0$ in Table~\ref{tab:covariances-sphere} and for $r_0$ given by $f_1$, $f_2$, or $f_3$ in \eqref{e:f_1}--\eqref{e:f_3}, $r_0(d_{\mathcal R}(u, v))$ is a valid correlation function, but by Theorem~\ref{t:1} 
we only know that
$r_0(d_{\mathcal G}(u, v))$ is 
 valid if $L$ is a 1-sum of trees and loops. 
 Note that in Table~\ref{tab:covariances-sphere}, the ranges of the parameters are the same for the two cases $S=\mathbb S^k$ and $S=L$. Finally, \eqref{e:f_1} is the special case of the generalized Cauchy function when $\alpha=1$ in Table~\ref{tab:covariances-sphere}, whilst \eqref{e:f_2} and \eqref{e:f_3} are not covered by Table~\ref{tab:covariances-sphere}. 
 
For example, consider the Chicago street network in the left panel of Figure~\ref{fig:data}. Here $L$ is not a 1-sum of trees and loops and therefore we cannot use the geodesic metric for the cases of covariance functions related to the Chicago street network. See also the counter examples in \citet[][Section~5]{AnderesEtAl}. On the other hand, the dendrite data shown in the right panel of Figure~\ref{fig:data} is observed on a tree, so here we can use the geodesic/resistance metric (by Theorem~\ref{t:1}(C) the two metrics are equal in this case).

\subsection{Simulation of GPs on linear networks}\label{s:sim-GP} 

This section discusses how to simulate a GP $Y=\{Y(u)\,|\,u\in L\}$ on a linear network $L=\cup_{i=1}^m L_i$ as given in case (c) in Section~\ref{s:new-setting} using three different algorithms, which are all available in our package \texttt{coxln}.
We assume without loss of generality that the mean function of $Y$ is zero.

The following is a straightforward algorithm applicable to any metric $d$ and any linear network $L$.

\begin{alg} Select a finite subset $D\subset L$ and make the following steps.
	\begin{itemize}\label{a:1}
		\item Simulate $Y$ restricted to $D$, e.g.\ by using  eigenvalue decomposition of the corresponding covariance matrix $\Sigma_D$.
		\item For $u\in L\setminus D$, approximate $Y(u)$ by the average of those $Y(v)$ where $v\in D$ is closest to $u$  with respect to $d_{\mathcal G}$.
	\end{itemize}
\end{alg}

Specifically, we have chosen a grid $D=V\cup D_1\cup…\cup D_m$ where each $D_i$ is a fine discretization of $L_i$ as described after the proof of Theorem~\ref{t:simalgotree} below.
%first selecting finite subsets $D_j\subset L_j$, $j=1,\ldots,m$; then simulating $Y$ restricted to the set $D$ given by the union of $V$ and $D_1,\ldots,D_m$, e.g.\ by using Choleski decomposition of the corresponding covariance matrix $\Sigma_D$; and finally, 
%for $u$ not in $V$ or $\cup_{j=1}^m D_j$, approximating $Y(u)$ by the average of those $Y(v)$ where $v\in D$ is closest to $u$ (`closest' with respect to $d_{\mathcal G}$). 
The disadvantage of Algorithm~\ref{a:1} is of course that the dimension of $\Sigma_D$ can be large and hence eigenvalue decomposition (as well as other methods) can be slow. Algorithm~\ref{a:2} below is much faster but requires $L$ to be a tree and $c(u,v) = \sigma^2\exp(-s d(u, v))$ to be an exponential covariance function with parameter $s>0$ and $d=d_{\mathcal G}=d_{\mathcal R}$ (the exponential correlation function $r_0(t)=\exp(-ts)$ appears as two special cases in Table~\ref{tab:covariances-sphere} with scale parameter $\phi=1/s$, namely the powered exponential model with $\alpha=1$ and the M{\'a}tern model with $\alpha=\tfrac12$). But we first need to establish a Markov property given in the following theorem, where we denote the shortest path between $u,v\in L$ by $p_{uv}$.

%Let $Y$ be a GP on $L$. The simplest algorithm to describe consists in selecting finite subsets $D_j\subset L_j$, $j=1,\ldots,m$, with e.g.\ $\{a_j,b_j\}\subset D_j$ and the geodesic distance between consecutive points in $D_j$ equal to a fixed number $c_j>0$ so that $D_j$ consists of $1+l_j/c_j$ points. First, we simulate $Y$ restricted to the set $D$ given by the union of $V$ and $D_1,\ldots,D_m$, e.g.\ using Choleski decomposition of the corresponding covariance matrix $\Sigma_D$. Second, for $u$ not in $V$ or $\cup_{j=1}^m D_j$, we approximate $Y(u)$ by the average of those $Y(v)$ where $v\in D$ is closest to $u$ (with respect to $d_{\mathbb G}$). The disadvantage of this algorithm is of course that the dimension of $\Sigma_D$ can be large and hence Choleski decomposition can be slow.   



%Assuming $L$ is a tree and $c(u,v) = \sigma^2\exp(-s d(u, v))$ is an exponential covariance function with parameter $s>0$ and $d=d_{\mathcal G}=d_{\mathcal R}$,  
%the algorithm in Theorem~\ref{t:simalgotree} below is much faster to use (the exponential correlation function $r_0(t)=\exp(-ts)$ appears as two special cases in Table~\ref{tab:covariances-sphere} with scale parameter $\phi=1/s$, namely the powered exponential model with $\alpha=1$ and the M{\'a}tern model with $\alpha=\tfrac12$). But we first need to establish a Markov property given in the following theorem, where we denote the shortest path between $u,v\in L$ by $p_{uv}$.
%% To prove Theorem~\ref{t:simalgotree} we first need to establish a Markov property given in the following theorem, where we denote the shortest path between $u,v\in L$ by $p_{uv}$.

\begin{thm}\label{t:markov}
	Suppose that $Y$ is a GP on a tree $L$ with exponential covariance function $c(u,v) = \sigma^2\exp(-s d(u, v))$ where $\sigma>0$, $s>0$, and $d=d_{\mathcal G}=d_{\mathcal R}$. 
	%For every pairwise distinct points $u,v,w\in L$ with $w\in p_{uv}$, we have that $Y(u)$ and $Y(v)$ are conditionally independent given $Y(w)$. Furthermore, f
	For $n=1,2,...$ and every pairwise distinct points $u,v,w_1,\ldots,w_n\in L$ so that $w_i\in p_{uv}$ for at least one $w_i$, we have that $Y(u)$ and $Y(v)$ are conditionally independent given $Y(w_1),\ldots,Y(w_n)$.
\end{thm}

\begin{proof}
	Let $n=1$ and $w=w_1$.
	Since $w\in p_{uv}$ and $L$ is a tree, $d(u,v) = d(u,w) + d(w,v)$ and therefore $c(u,v) = c(u,w)c(w,v)/\sigma^2$. Thus the covariance matrix for $(u,v,w)$ has the form 
	\begin{equation*}
	\Sigma_{u,v,w} = 
	\begin{pmatrix}
	\sigma^2 & c(u,w)c(w,v)/\sigma^2 & c(u,w) \\
	c(u,w)c(w,v)/\sigma^2 & \sigma^2 & c(w,v) \\
	c(u,w) & c(w,v) & \sigma^2
	\end{pmatrix}.
	\end{equation*}
	Inverting the covariance matrix, we get that the corresponding precision matrix has 0 at entries $(1,2)$ and $(2,1)$, thus implying that $Y(u)$ and $Y(v)$ are conditionally independent given $Y(w)$. 
	
	 Consider the case $n=2$ and e.g.\ $w_1\in p_{uv}$. Since $L$ is a tree, $Y(w_2)$ must be conditionally independent of either $Y(u)$ or $Y(v)$ given $Y(w_1)$. Assume without loss of generality that this is $Y(u)$. Thus, $Y(u)$ is conditionally independent of $(Y(v),Y(w_2))$ given $Y(w_1)$, which implies that $Y(u)$ and $Y(v)$ are conditionally independent given $(Y(w_1),Y(w_2))$. In a similar way we  verify the case with $n\ge3$. 
\end{proof}

%{\color{red}
%
%\begin{thm}\label{t:2}
%Suppose that for some $s>0$, $\exp(-s d_{\mathcal G}(u, v))$ is positive definite over $(u, v)\in L\times L$. Let 
%%$c(u,v)=\sigma^2r(u,v)$ where 
%$\sigma>0$ %and $r(u,v)=\exp(-s d_{\mathcal G}(u, v))$, 
%and consider the following algorithm, where $r(u,a_j)=\exp(-\|u-a_j\|s)$, $r(u,b_j)=\exp(-\|u-b_j\|s)$, $r(v,a_j)=\exp(-\|v-a_j\|s)$, and $r(v,b_j)=\exp(-\|v-b_j\|s)$. 
%\begin{itemize}
%\item Generate $Y_V=\{Y(u)\,|\,u\in V\}$ as a zero mean Gaussian vector with covariances $c(u,v)$, $u,v\in V$.
%\item Conditioned on $Y_V$, generate independent GPs $Y_j=\{Y_j(u)\,|\,u\in L_j\setminus\{a_j,b_j\}\}$, $j=1,\ldots,m$ so that for $u,v\in L_j\setminus\{a_j,b_j\}$,
%the conditional mean of $Y_j(u)$ depends only on $Y(a_j)$ and $Y_j(b_j)$ and is given by
%\begin{align}
%\mathbb E[Y_j(u)\,|\,Y_V] %&=\mathbb E[Y_j(u)\,|\,Y(a_j),Y(b_j)] \nonumbers\\
%&=\frac{r(u,a_j)-r(u,b_j)\exp(-l_js)}{1-\exp(-2l_js)}Y(a_j)+\frac{r(u,b_j)-r(u,a_j)\exp(-l_js)}{1-\exp(-2l_js)}Y(b_j)
%%\frac{\left(r(u,a_j)-r(u,b_j)\exp(-l_js)\right)Y(a_j)+\left(r(u,b_j)-r(u,a_j)\exp(-l_js)\right)Y(b_j)}{1-\exp(-2l_js)}
%\label{e:M-a}
%\end{align}
%and the conditional covariance of $Y(u)$ and $Y(v)$ does not depend on $Y_V$ and is given by
%\begin{align}
%\mathbb C{\mathrm{ov}}[Y_j(u),Y_j(v)\,|\,Y_V] %&=\mathbb C{\mathrm{ov}}(Y_j(u),Y_j(v)) \nonumbers\\
%&=\sigma^2\frac{r(u,v)\left(1+\exp(-2l_js)\right)-r(u,a_j)r(v,a_j)-r(u,b_j)r(v,b_j)}{1-\exp(-2l_js)}.\label{e:C-a}
%\end{align}
%\item Output $Y=Y_V\cup Y_1\cup\ldots\cup Y_m$.
%\end{itemize}
%Then $Y$ is a zero mean GP on $L$ with exponential covariance function $c(u,v)=\sigma^2\exp(-s d_{\mathcal G}(u, v))$.
%\end{thm}
%
%\begin{proof}
%First, recall the following properties for the special case where $L$ is a closed interval divided into $m$ consecutive closed intervals $L_j=[a_j,b_j]$, $j=1,\ldots,m$, with $a_{j+1}=b_j$ if $2\le j<m$. Suppose $Y=\{Y(u)\,|\,a_1\le u\le b_m\}$ is a zero mean GP with covariance function $c(u,v)=\sigma^2\exp(-|u-v|s)$. Let $Y_j=\{Y(t)\,|\,a_j< t< b_{j}\}$, $j=1,\ldots,m$.
%Since $Y$ is Markov, conditioned on $Y_V=\{Y(a_j)\,|\,j=1,\ldots,m\}$, the $Y_j$ are conditionally independent GPs with mean and covariance functions as given in \eqref{e:M-a} and \eqref{e:C-a}. 
%
%Second, consider the algorithm in the theorem for any linear network $L$. Clearly, the output is a zero mean GP and $Y_V$ has been generated in accordance to a GP with the exponential covariance function $c$. For any $u,v\in L\setminus V$, let $L_j$ and $L_i$ be the line segments with $u\in L_j$ and $v\in L_i$, let
%$p_{uv}$ be a shortest path connecting $u$ and $v$, and let $u_1,\ldots,u_h$ be the $h\ge2$ vertices such that $u_1=a_j$, $u_2=b_j$, $u_{h-1}=a_i$, $u_h=b_i$, $u_{r-1}\sim u_{r}$ for $r=2,\ldots,h$, and $u_r\in p_{uv}$ if $2\le r\le h-1$. Since the linear network given by the $h$ consecutive line segments with endpoints $u_0,\ldots,u_h$ can be identified with a linear network on an interval given by $h$ closed line segments/intervals of lengths $\|b_1-a_1\|,\ldots,\|b_h-a_h\|$, it follows that the output in the algorithm is a zero mean GP with covariance function $c(u,v)=\sigma^2\exp(-s d_{\mathcal G}(u, v))$.
%\end{proof}
%}

We use the following notation in Algorithm~\ref{a:2}. Pick an arbitrary origin $u_0\in V$ and set $G_0(u_0)=\{u_0\}$. For $j=1,2,\ldots$, if $u\in V\setminus\cup_{i=0}^{j-1}G_i(u_0)$ and $u\sim v$ for some $v\in G_{j-1}(u_0)$, we call $u$ a child of $j$-th generation to $u_0$ and define $G_j(u_0)\subset V$ as the set of all children of $j$-th generation to $u_0$. 
Moreover, %for the GP $Y=\{Y(w)\,|\,w\in L\}$ constructed in the following theorem,
if $u\in G_{j-1}(u_0)$, $v\in G_j(u_0)$, and $u\sim v$, we define a GP $Y(u,v)=\{Y(w)\,|\,w\in(u,v]\}$ where $(u,v]$ is the half-open
line segment with endpoints $u$ and $v$ so that $u$ is excluded and $v$ is included. 
%Moreover, for the GP $Y=\{Y(u)\,|\,u\in L\}$ constructed in the following theorem, if $u\in G_{j-1}(u_0)$, $v\in G_j(u_0)$, and $u\sim v$, we define $Y(u,v)$ as the collection of all $Y(w)$ where $w$ belongs to the half-open line segment with endpoints $u$ and $v$ so that $u$ is excluded and $v$ is included. 

\begin{alg}\label{a:2}
Suppose that $L$ is a tree. Let $s>0$ and $\sigma>0$ be parameters, and pick an arbitrary vertex $u_0\in V$. Construct random variables $Y(w)$ for all $w\in L$ by the following steps.
	%$c(u,v)=\sigma^2r(u,v)$ where 
	%$\sigma>0$ and %$r(u,v)=\exp(-s d_{\mathcal G}(u, v))$, and 
	%Consider the following iterative construction of random variables $Y(w)$ where . 
	\begin{description}
		\item[(I)] For $w=u_0$, generate $Y(w)$ from $N(0,\sigma^2)$.% and set $Y_0=\{Y(u_0)\}$.
		\item[(II)] For $j=1,2,\ldots$, conditioned on all the $Y(w)$ so far generated,  generate independent GPs $Y(u,v)$ for all $u\in G_{j-1}(u_0)$ and all $v\in G_j(u_0)$ with $u\sim v$, where $Y(u,v)$ depends only on $Y(u)$ and for every $w,w_1,w_2\in(u,v]$ we have
		\begin{align}\label{e:M2}
		\mathbb E[Y(w)\,|\,Y(u)]&=\exp(-\|w-u\|s))Y(u)\\
		%,\quad \mathbb C{\mathrm{ov}}[Y(w_1,w_2)\,|\,Y(u)]=\sigma^2\left(\exp(-\|w_1-w_2\|s)-\exp(-\|w_1-u\|s-\|w_2-u\|s)\right).
		\label{e:M22}
		\mathbb C{\mathrm{ov}}[Y(w_1,w_2)\,|\,Y(u)]&=\sigma^2\left(\exp(-\|w_1-w_2\|s)-\exp(-\|w_1-u\|s-\|w_2-u\|s)\right).
		\end{align} 
		\item[(III)] Output $Y=\{Y(w)\,|\,w=u_0\mbox{ or }w\in (u,v]\mbox{ for some }j\in\mathbb N,\ u\in G_{j-1}(u_0),\ v\in G_j(u_0)\mbox{ with }u\sim v\}$.
	\end{description}
\end{alg}

%Using Theorem~\ref{t:markov}, we can now prove that the simulation algorithm given in Theorem~\ref{t:simalgotree} below generates a GP $Y=\{Y(w)\,|\,w\in L\}$ with an exponential covariance function on a tree. We use the following notation in Theorem~\ref{t:simalgotree}. Pick an arbitrary origin $u_0\in V$ and set $G_0(u_0)=\{u_0\}$. For $j=1,2,\ldots$, if $u\in V\setminus\cup_{i=0}^{j-1}G_i(u_0)$ and $u\sim v$ for some $v\in G_{j-1}(u_0)$, we call $u$ a child of $j$-th generation to $u_0$ and define $G_j(u_0)\subset V$ as the set of all children of $j$-th generation to $u_0$. 
%Moreover, %for the GP $Y=\{Y(w)\,|\,w\in L\}$ constructed in the following theorem,
%if $u\in G_{j-1}(u_0)$, $v\in G_j(u_0)$, and $u\sim v$, we define $Y(u,v)=\{Y(w)\,|\,w\in(u,v]\}$ where $(u,v]$ is the half-open
%line segment with endpoints $u$ and $v$ so that $u$ is excluded and $v$ is included. 
%%Moreover, for the GP $Y=\{Y(u)\,|\,u\in L\}$ constructed in the following theorem, if $u\in G_{j-1}(u_0)$, $v\in G_j(u_0)$, and $u\sim v$, we define $Y(u,v)$ as the collection of all $Y(w)$ where $w$ belongs to the half-open line segment with endpoints $u$ and $v$ so that $u$ is excluded and $v$ is included. 


\begin{thm}\label{t:simalgotree}
	Let the situation be as in Algorithm~\ref{a:2}.
	Then $Y$ is a zero mean GP on $L$ with exponential covariance function $c(u,v)=\sigma^2\exp(-s d(u, v))$. % where $d=d_{\mathcal G}=d_{\mathcal R}$.
\end{thm} 

\begin{proof}
	By considering the GPs associated to the successive generations to $u_0$, we prove by induction that %if $Y$ is a GP as stated in t theorem, then 
	the iterative construction in Algorithm~\ref{a:2} makes $Y$ a GP as stated. 
	
	Clearly the distribution of $Y(u_0)$ is correctly generated, cf.\ (I). This is induction step $j=0$. 
	
	For induction step $j\ge1$, we condition on all the $Y(w)$ so far generated, i.e., $w$ is equal to either $u_0$ or $Y(w)$ is a member of a GP associated to two neighbouring vertices, with one being a child of $i$-th generation $i$ to $u_0$ and the other being a child of $(i-1)$-th generation $i$ to $u_0$ where $1\le i\le j$, cf.\ (I) and (II) in Algorithm~\ref{a:2} 
	(here we interpret $u_0$ as a child of $0$-th generation to $u_0$). 
	By the induction hypothesis, these $Y(w)$ have been correctly generated. If we take two points $w_1$ and $w_2$ contained in different line segments between points in $G_{j-1}(u_0)$ and $G_j(u_0)$, then by Theorem~\ref{t:markov}, $Y(w_1)$ and $Y(w_2)$ are (conditionally) independent, which is in accordance to our construction. So it suffices to consider the (conditional) distribution of $(Y(w_1),Y(w_2))$ when $w_1,w_2\in(u,v]$, $u\in G_{j-1}(u_0)$, $v\in G_j(u_0)$, and $u\sim v$. This (conditional) distribution depends only on $Y(u)$, cf.\ Theorem~\ref{t:markov}, and it is straightforwardly seen to be  
a bivariate normal distribution with mean and covariance matrix given by \eqref{e:M2} and \eqref{e:M22}, respectively.	
%	simulated $Y(u)$ for all $u\in \cup_{i=0}^{j-1}G_i(u_0)$, we then need to verify that $Y(w)$ for $w$ in the half open line segment between $u$ and $v$ with $v\in G_j(u_0)$ has the correct distribution. If we take two points $w_1$ and $w_2$ contained in different line segments between points in $G_{j-1}(u_0)$ and $G_j(u_0)$, then by Theorem~\ref{t:markov}, $Y(w_1)$ and $Y(w_2)$ are conditionally independent given all $Y(u)$ for $u\in\cup_{i=0}^{j-1}G_i(u_0)$, so it suffices to consider the conditional distribution of one of these, say $Y(w_1)$ located on the segment from $u_1\in G_{j-1}(u_0)$ to $v_1\in G_{j}(u_0)$. By Theorem~\ref{t:markov}, $Y(w_1)$ is conditionally independent of $Y(u)$ for all $u\in\cup_{i=0}^{j-1}G_i(u_0)\backslash\{u_1\}$ given $Y(u_1)$. Thus we consider the conditional distribution of $Y(w_1)$ given $Y(u_1)$, which is a normal distribution with mean and covariance functions given by \ref{e:M2}. 

	%Therefore the output of the algorithm has the correct distribution. 
\end{proof}


In practice, when using Algorithm~\ref{a:2} for simulation, a discretization on each line segment is needed. For some integer $n_j>0$ (which may depend on $l_j$) and $i\in\{0,1,\ldots,{n_j}\}$, define $s_i=i l_j/{n_j}$. Then, for
each $u\in L_j$, if $u$ is the midpoint between $s_{i-1}$ and $s_i$, we approximate $Y(u)$ by the average $(Y(u_j(s_{i-1}))+Y(u_j(s_i)))/2$, and otherwise if $u$ is closest to $u_j(s_i)$, we approximate $Y(u)$ by 
$Y(u_j(s))$. Further,
to generate the ${n_j}-1$ normal variables $Y(u_j(s_i))$, $i=1,\ldots,{n_j}-1$, we start by generating the variable $Y(u_j(s_1))$ in accordance to \eqref{e:M2} and \eqref{e:M22} with $w=w_1=w_2=s_1$. Then we can add $u_j(s_1)$ to the vertex set, whereby we split $l_j$ into the two line segments given by this new vertex and the endpoints of $l_j$. Hence, if $n_j>1$, we can repeat the procedure when generating $Y(u_j(s_2))$, and so on until all the $n_j-1$ normal variables have been generated.  

%If $L$ is not a tree, we use a triadic scheme as follows. For each $u\in L_j$, we approximate $Y(u)$ by   
%$Y(u_j(s))$ if $u$ is closest to $u_j(s_i)$ with now $s_i=i l_j/3^{n_j}$ and $i\in\{0,1,\ldots,3^{n_j}\}$; or, if $u$ is the midpoint between $s_{i-1}$ and $s_i$, we approximate $Y(u)$ by the average $(Y(u_j(s_{i-1}))+Y(u_j(s_i)))/2$. To generate the $3^{n_j}-1$ normal variables $Y(u_j(s_i))$, $i=1,\ldots,3^{n_j}-1$, we start by generating the two variables $Y(u_j(s_i))$, $i=3^{n_j-1},2\times 3^{n_j-1}$, in accordance to \eqref{e:M-a} and \eqref{e:C-a}. Then we can add $u_j(s_i)$ for $i=3^{n_j-1},2\times 3^{n_j-1}$ to the vertex set, whereby we split $l_j$ into the three line segments given by these two new vertices and the endpoints of $l_j$. Hence, if $n_j>1$ we can repeat the procedure of generating two normal variables on each of these three line segments, and so on until all the $3^{n_j}-1$ normal variables have been generated.   

%It is important to notice that 
Theorems~\ref{t:markov} and \ref{t:simalgotree} do not hold if $L$ is not a tree. Indeed, a GP on the circle $\mathbb S^1$ (which is equivalent to a loop of length $2\pi$) with exponential covariance function, considering four arbitrary distinct points on $\mathbb S^1$, it can be shown that the GP is not  
Markov by calculating the precision matrix of the GP at these four points. On the other hand, letting $c(u,v) = a \cosh(b(d_\mathcal{G}(u,v) -\pi))$ for $a,b>0$,
\cite{Pitt71} verified that the GP on $\mathbb S^1$ with covariance function $c$ is Markov, but considering two arbitrary points on a tree together with a point on the path connecting them, it can be shown that the GP on the tree with covariance function $c$ is not Markov.
Consequently, we cannot have a covariance function only depending on the geodesic distance which makes an arbitrary linear network Markov. Moreover, 
 in general, if we want to simulate a GP with an exponential covariance function on a linear network which is not a tree, we cannot rely on Markov properties. Instead we have to use the straightforward but slower Algorithm~\ref{a:1}.
 
For other covariance functions than the exponential, %Theorem~\ref{t:simalgotree} can be used in connection to 
we may use the following theorem, which follows immediately from \eqref{e:bernstein} and the central limit theorem.

\begin{thm}\label{t:3}
Suppose $d$ is a metric on $L$ so that $(u,v)\mapsto\exp(-sd(u,v))$ for $(u,v)\in L^2$ is a well-defined correlation function for all $s>0$, and let $Y$ be a zero mean GP on $L$ with
	covariance function 
	\begin{equation}\label{e:c-ber}
	c(u,v)=\sigma^2
	\int\exp(-sd(u,v))\,\mathrm dF(s)
	\end{equation}
	where $\sigma>0$ and $F$ is a CDF with $F(s)=0$ for $s<0$. For an integer $n>0$ and $i=1,\ldots,n$, suppose we generate first $S_i$ from $F$ and second $Y_i$ as a zero mean GP on $L$ with covariance function $\sigma^2\exp(-S_i d(u,v))$ so that $(S_1,Y_1),\ldots,(S_n,Y_n)$ are independent. Let $\bar Y_n=\sum_{i=1}^n Y_i/n$. Then $\sqrt n \bar Y_n$ is a zero mean stochastic process on $L$ with covariance function $c$. As $n\rightarrow\infty$, $\sqrt n \bar Y_n$ approximates $Y$ in the sense that any finite dimensional distribution of $\sqrt n \bar Y_n$ converges in distribution towards the corresponding finite dimensional distribution of $Y$.  
\end{thm}   

Theorem~\ref{t:3} allows simulation of any GP with a covariance function of the form \eqref{e:c-ber},
if a simulation algorithm for $F$ is available. For the case of a tree, Theorem~\ref{t:3} combined with Algorithm~\ref{a:2} gives the following simulation algorithm.

\begin{alg}\label{a:3}
Let the situation be as in Theorem~\ref{t:3} and suppose that $L$ is a tree and $d=d_{\mathcal R}=d_{\mathcal G}$. 
Choose an integer $n>0$ and independently for $i=1,...,n$ make the following steps. 
\begin{description}
		\item[(I)] Generate $S_i$ from $F$.
		\item[(II)] Using Algorithm~\ref{a:2} generate $Y_i$ as a zero mean GP on $L$ with covariance function $\sigma^2\exp(-S_i d(u,v))$.
\end{description}
Output 	$\sqrt n \bar Y_n$ as an approximate simulation of 	a zero mean GP with covariance function given by \eqref{e:c-ber}.
\end{alg}

Conditioned on $S_1,...,S_n$, the output in Algorithm~\ref{a:3} is a zero mean GP with covariance function
\begin{equation}\label{e:approxc}
c(u,v|S_1,...,S_n)=\frac{\sigma^2}{n}\sum_{i=1}^n\exp(-S_i d(u,v)).
\end{equation}	
The running times of Algorithms~\ref{a:1}--\ref{a:3} for simulating a GP on a tree are compared in the following example,  and the  choice of $n$ in Algorithm~\ref{a:3} needed to obtain that \eqref{e:approxc} is a good approximation of the covariance function \eqref{e:c-ber} is considered in the example thereafter.

\begin{ex} 
Let $L$ be the dendrite tree in Figure~\ref{fig:data}.
Table~\ref{tab:algotimes} shows the times used in our implementation for obtaining a simulation of a GP with either an exponential covariance function or a covariance function with inverse gamma Bernstein CDF, cf.\ \eqref{e:f_2}. 
Specifically, $\tau=\phi=\sigma=1$ but the running times do not depend on this choice of the parameters. However,
the running times depend heavily on the number of points $\#D$ in the grid $D\subset L$, so we consider two different situations, one with $\#D=387$ (this choice is used later in Sections~\ref{s:den-an} and \ref{s:simstudy}) and another with $\#D=1863$. Since Algorithm~\ref{a:3} depends on the choice of $n$, various values of $n$ are also shown in Table~\ref{tab:algotimes}. The first two rows in the table show that Algorithm~\ref{a:2} is faster than Algorithm~\ref{a:1} for simulating a GP with exponential covariance function for both choices of grids, but the difference is far more clear for the fine grid. For the inverse gamma Bernstein CDF in the case $\#D=387$, the running time of Algorithm~\ref{a:1} is roughly the same as for Algorithm~\ref{a:3} with $n=20$, while in the case $\#D=1863$ Algorithm~\ref{a:3} is much faster than Algorithm~\ref{a:1} even with $n=200$. This illustrates that Algorithm~\ref{a:2} is faster than Algorithm~\ref{a:1}, while it depends on the choice of $\#D$ and $n$ whether Algorithm~\ref{a:1} or \ref{a:3} is fastest.

\begin{table}
	\centering
	\begin{tabular}{l|c||c|c}
		Covariance function & Algorithm & $\#D = 387$ & $\#D = 1863$ \\
		\hline\hline
 		Exponential & Algorithm~\ref{a:1} & 0.129 s & 13.7 s \\
 		Exponential & Algorithm~\ref{a:2} & 0.0118 s &   0.0166 s \\
 		Inverse gamma Bernstein CDF & Algorithm~\ref{a:1} & 0.157 s & 14.5 s\\
 		Inverse gamma Bernstein CDF & Algorithm~\ref{a:3} (n=20) & 0.162 s & 0.255 s\\
 		Inverse gamma Bernstein CDF & Algorithm~\ref{a:3} (n=50) & 0.400 s & 0.638 s\\
 		Inverse gamma Bernstein CDF & Algorithm~\ref{a:3} (n=200) &  1.61 s & 2.48 s\\
	\end{tabular}
	\caption{Running times of Algorithms~\ref{a:1}--\ref{a:3} used for simulating a GP on two different grids on a tree.} 
	\label{tab:algotimes}
\end{table}
\end{ex}

\begin{ex} The choice of $n$ in Algorithm~\ref{a:3} leading to a good approximation depends on the choice of covariance function, so in 
continuation of the previous example, we consider two inverse gamma Bernstein CDFs with parameters 
$(\tau,\phi)=(1,1)$ and $(\tau,\phi)=(5,5)$. For these choices of parameters and for $n=20,50,200$, 
Figure~\ref{fig:algo3} shows the correlation functions along with 100 approximations given by \eqref{e:approxc}. For $n=20$ the approximations show a lot of variability around the true correlation functions, while for $n=200$ the approximations are much closer to the true correlation functions. Also it is evident in the figure that a higher $n$ is required to obtain a good approximation for $(\tau,\phi)=(1,1)$ than for $(\tau,\phi)=(5,5)$. This is expected since the variance of $\Gamma^{-1}(\tau,\phi)$ is infinite if and only if $0<\tau\le2$.


\begin{figure}
	\begin{center}
		\includegraphics[width=6cm]{"Figs/sim-corfunc-nsim20-par11.pdf"}
		\includegraphics[width=6cm]{"Figs/sim-corfunc-nsim20-par55.pdf"}\\
		\includegraphics[width=6cm]{"Figs/sim-corfunc-nsim50-par11.pdf"}
		\includegraphics[width=6cm]{"Figs/sim-corfunc-nsim50-par55.pdf"}\\
		\includegraphics[width=6cm]{"Figs/sim-corfunc-nsim200-par11.pdf"}
		\includegraphics[width=6cm]{"Figs/sim-corfunc-nsim200-par55.pdf"}
	\end{center}
	\caption{Correlations functions with inverse gamma Bernstein CDFs (black curves) with $(\tau,\phi)=(1,1)$ (left column) and $(\tau,\phi)=(5,5)$ (right column). The grey curves in each plot show 100 approximated correlation functions given by \eqref{e:approxc} with $n=20$ (upper row), $n=50$ (middle row), and $n=200$ (lower row).}
	\label{fig:algo3}
\end{figure}

\end{ex}

If $L$ is not a tree, Algorithms~\ref{a:2} and \ref{a:3} cannot be used, so we use the straightforward Algorithm~\ref{a:1} instead, provided of course that $Y$ is specified by a valid covariance function, meaning that Algorithm~\ref{a:1} may work for $d=d_{\mathcal{R}}$ but not for $d=d_{\mathcal{G}}$ as illustrated in the following example.

\begin{ex}  Figure~\ref{fig:chicago-GP} shows examples of simulations of zero mean GPs defined on the Chicago street network in Figure~\ref{fig:data} and with various isotropic covariance functions $c(u,v)=\sigma^2r_0(d(u,v))$ where $\sigma=1$ and in order that $c$ is valid we take $d=d_{{\mathcal R}}$, cf.\ Theorem~\ref{t:1}.  In the first two plots (the top row), $r_0(t)=\exp(-st)$ is 
an isotropic exponential correlation function with $s=0.1$ or $s=0.01$, and the next two plots (the middle row) relate to Theorem~\ref{t:3} with the Bernstein CDF given by a $\Gamma(1,100)$-distribution or a $\Gamma^{-1}(2,0.01)$-distribution, cf.\ Example~\ref{ex:0}. The densities for these Bernstein CDFs are shown in the left panel in the bottom row, and the last plot shows the   
corresponding correlation functions for $t\le200$ feet. 
%Since the resistance metric is used, the covariance functions for the GPs in plots 1--4 are well-defined, cf.\ Theorem~\ref{t:1}. 
The top row shows the scaling effect of the parameter $s$ for the exponential correlation function. Plots 2--4 are comparable, since $s$ has mean 0.01 in all three cases. %the mean values of $s$ all agree. 
Since $\max d_{\mathcal R}\approx 675$ feet, the last plot
indicates a rather strong correlation in the GPs in plots 2--4 and shows that the correlation is smallest when $r$ is fixed and rather similar when using the $\Gamma$ or $\Gamma^{-1}$ distribution although these distributions are rather distinct (e.g.\ the variance is finite for $\Gamma(1,100)$ but infinite for $\Gamma^{-1}(2,0.01)$). Accordingly, in plot 2 we see less smoothness than in plots 3 and 4 which show a similar degree of smoothness. 

\begin{figure}
	\begin{center}
	\includegraphics[width=6.5cm]{"Figs/chicago-GPexp0.1.pdf"}
	\includegraphics[width=6.5cm]{"Figs/chicago-GPexp0.01.pdf"}\\
%	\includegraphics[width=6cm]{"Figs/chicago-GPexp0.001.pdf"}
	\includegraphics[width=6.5cm]{"Figs/chicago-GPgamma.pdf"}
	\includegraphics[width=6.5cm]{"Figs/chicago-GPinvgamma.pdf"}\\
	\includegraphics[width=6cm]{"Figs/chicago-densities.pdf"}
	\includegraphics[width=6cm]{"Figs/chicago-covfuncs.pdf"}
	\end{center}
	\caption{Simulation of zero mean GPs on the Chicago street network with an isotropic exponential covariance function $c(u,v)=r_0(d_{\mathcal{R}}(u,v))$. Top row: When $r_0(t)=\exp(-st)$ with parameter $s=0.1$ (left) or $s=0.01$ (right). Middle row: When $r_0$ has Bernstein CDF given by a gamma distribution (left) or inverse gamma distribution (right), where in both cases the mean of $s$ is 0.01. Bottom row: For plots 2--4, the corresponding densities for $F$ and correlation functions, where the curves in solid, dashed, and dotted correspond to plots 2--4, respectively.}
	\label{fig:chicago-GP}
\end{figure}

\end{ex}

\section{Point processes and some of their characteristics}\label{s:3}

Throughout this section we 
consider a general metric space $(S,d)$ equipped with its Borel $\sigma$-algebra $\mathcal B$ and
a reference measure $\nu$. Sometimes this is specified to be $S\in\{\mathbb R^k,\mathbb S^k,L\}$ with $d$ and $\nu$ as given in Section~\ref{s:new-setting}.

%The setting for point processes used in this paper is given in Section~\ref{s:3.1}. Since Cox and the other point process models treated in this paper possess useful moment properties, Section~\ref{s:3.2} provides the general definition of first and higher order intensity functions. In particular, we focus on the pair correlation function and
%%, depending on the choice of space, 
%the related $K$-function defined in Section~\ref{s:3.2}. These and other functional characteristics for point processes become useful for statistical inference as discussed in Sections~\ref{s:est-check} and \ref{s:appl}.   

%\section{Setting}\label{s:2.1}

%The state space for the points of the point processes considered in this paper is a metric space $(S,d)$ equipped with a reference measure $\nu$. I concentrate on the following specific cases (a)--(c).
%\begin{description}
%\item[(a)] $S=\mathbb R^k$ is the $k$-dimensional Euclidean space, $d(u,v)=\|u-v\|$ is the usual Euclidean distance, and $\mathrm d\nu(u)=\mathrm du$ is Lebesgue measure.
%\item[(b)] $S=\mathbb S^k=\{u\in \mathbb R^{k+1}\,|\, \|u\|=1\}$ is the $k$-dimensional unit sphere, $d(u,v)=d_\mathcal{G}(u,v)=\arccos(u\cdot v)$ is the geodesic (orthodromic, great-circle, or shortest path) metric, and $\nu=\nu_k$ is the surface ($k$-dimensional Hausdorff) measure, see e.g.\ 
%\citet[][Chapter~1]{MR3060033} or Appendix~\ref{a:A1}. Here, $u\cdot v$ is the usual inner product of $u$ and $v$.
%\item[(c)] $S=L=\cup_{i=1}^m L_i$ is a linear network, that is, a connected set given by a finite union of closed line segments $L_i$ of length $l_i\in (0,\infty)$, $i=1,\ldots,m$, and which can only overlap at the endpoints of the line segments. Furthermore, $d=\delta$ is some `natural' metric $\delta$ and $\mathrm d\nu(u)=\mathrm d_L(u)$ is arc length measure.  
%\end{description} 
%In connection to (c) some remarks are in order. For disconnected networks, definitions and results may be applied
%separately to each connected component of the network. For each line segment $L_i$, there are two possible arc length parametrisations. I assume one is chosen and given by $u_i(t)=(1-t/l_i)a_i+(t/l_i)b_i$, $t\in[0,l_i]$, where $a_i$ and $b_i$ are the endpoints of $L_i$. The definitions and results in this paper will not depend on this choice. It is used when calculating arc length measure restricted to $L_i$: For Borel sets $A\subseteq L_i$, $\mathrm d_L(A)=\int_0^{l_i}1(u_i(t)\in A)\,\mathrm dt$, where $1(\cdot)$ denotes the indicator function. I have mainly two cases of natural metrics in mind when $S=L$: $\delta$ is given by the geodesic metric $d_\mathcal{G}$ or the resistance metric $d_\mathcal{R}$ on $L$ as developed in \citet{AnderesEtAl}. Appendix~\ref{a:A2} provides the definition and a specific expression useful for the calculation of $d_\mathcal{R}(u,v)$ for all $u,v\in L$, cf.\ \eqref{e:d-R}. As discussed in Sections~\ref{s:2.2}--\ref{s:covfct},
%the resistance metric has several advantages over the geodesic metric in the isotropic case of covariance functions, i.e., when considering covariance functions which only depend on the metric. Moreover, $d_\mathcal{G}\le d_\mathcal{R}$, with equality if and only if $L$ is a tree. Indeed there are other interesting metric when $S=L$, including the least-cost metric \citep{RakshitEtAl}, %Rakshit, Nair, Baddeley (2017), Spat. Stat.
%but to the best of my knowledge parametric models for isotropic covariance functions have so far only been developed when $\delta=d_\mathcal{G}$, $\delta=d_\mathcal{R}$, or $\delta$ is the usual Euclidean distance (which is usually not a natural metric on a linear network). Finally, (c) may be extended to the more abstract case of a graph with Euclidean edges, i.e., graphs with edges viewed as line or curve segments allowing us to consider points on the graph which are either vertices or points on an edge, cf.\ \citet{AnderesEtAl}.

\subsection{Point process setting}\label{s:3.1}

%Consider again a general metric space $(S,d)$, where $S=L$ is of our primary interest, cf.\ Section~\ref{s:new-setting}.
 We restrict attention to point processes whose realisations can be viewed as locally finite subsets $x$ of $S$: For  
 $x,A\subseteq S$, define $x_A= x\cap A$ and let $n(x_A)$ denote the cardinality of $x_A$. We let the state space of 
 a point process be $$\mathcal N= \{x\subset S\,|\,n(x_A)<\infty \mbox{ whenever $A\subseteq S$ is bounded}\}$$ 
 and equip  $\mathcal N$ with the smallest $\sigma$-algebra $\mathcal F$ such that the mapping $x\mapsto n(x_A)$ is measurable whenever $A\subseteq S$ is a bounded Borel set. So
 by a point process is meant a random variable $X$ with values in $\mathcal N$ \citep[in the terminology of point process theory, $X$ is a simple locally finite point process, see e.g.][]{DaleyVere-Jones}. Then, for any bounded Borel set $A\subseteq S$, the count $N(A)=n(X_A)$ is a random variable.
 
We need the following notions. % where $\nu$ is a reference measure on $S$, cf.\ Section~\ref{s:new-setting}.  

We say that $X$ is a Poisson process with intensity function $\rho:S\mapsto[0,\infty)$ if for any bounded Borel set $A\subseteq L$, $N(A)$ is Poisson distributed with finite mean $\int_A\rho(u)\,\mathrm d\nu(u)$, and conditioned on $N(A)$, the points in $X_A$  are independent and each point has a density proportional to $\rho$ with respect to $\nu$ restricted to $A$. 

For every $u\in L$, we let $X_u$  be 
the point process which follows the reduced Palm distribution of $X$ at $u$, that is,
 \[\mathbb E\sum_{u\in X}h(X\setminus\{u\},u)=\int\rho(u)\mathbb E h(X_u,u)\,\mathrm d\nu(u)\]
 for any non-negative measurable function $h$ defined on $\mathcal N\times S$ (equipped with the product $\sigma$-algebra of $\mathcal F$ and $\mathcal B$). Intuitively, $X_u$ follows the distribution of $X\setminus\{u\}$ conditioned on that $u\in X$ \citep[see e.g.][Appendix C]{MW04}. If $\rho(u)=0$, 
 %it makes not much sense to consider the reduced Palm distribution, since 
 $X_u$ may follow an arbitrary distribution. If   
  $X$ is a Poisson process with intensity function $\rho$, then $X$ and $X_u$ are identically distributed whenever 
 $\rho(u)>0$.
 
 Let $X_1$ denote the Poisson process with intensity 1. Suppose $S\in\{\mathbb S^k,L\}$ and the distribution of $X$ is absolutely continuous with respect the distribution of $X_1$ (this makes sense since $S$ is bounded and hence $X$ is finite).  Let $f$ be the density of (the distribution of) $X$ with respect to (the distribution of) $X_1$.
 If $\rho(u)>0$, then $X_u$ has a density $f_u$ with respect to  $X_1$ such that
\begin{equation}\label{e:Palmdensity}
 f_u(\{u_1,\ldots,u_n\})=f(\{u,u_1,\ldots,u_n\})/\rho(u)\quad\mbox{for $n=1,2,\ldots$ and pairwise distinct }u_1,\ldots,u_n\in S\setminus\{u\}.
 \end{equation} 

\subsection{Moment and invariance properties of point processes}\label{s:3.2}

%Let $S$ and the reference measure $\nu$ be as in one of the cases (a)--(c) in Section~\ref{s:new-setting}. 
The point process $X$ has $n$-th order intensity function $\rho(u_1,\ldots,u_n)$ (with respect to the $n$-fold product measure of $\nu$) if this is a non-negative Borel function so that
\begin{equation}\label{e:1}
\mathbb E\left[N(A_1)\cdots N(A_n)\right]=\int_{A_1}\cdots\int_{A_n}\rho(u_1,\ldots,u_n)\,\mathrm d\nu(u_1)\cdots\,\mathrm d\nu(u_n)<\infty
\end{equation}
for every pairwise disjoint, bounded, Borel sets $A_1,\ldots,A_n\subset S$ ($\rho(u_1,\ldots,u_n)$ is also called the $n$-th order product density for the $n$-th order reduced moments measure). Thus, 
 $\rho(u_1,\ldots,u_n)$ is a locally integrable function, which is almost everywhere unique on $S^n$ (with respect to the $n$-fold product measure of $\nu$). 
In the following, for simplicity nullsets are ignored, so the non-uniqueness of $\rho(u_1,\ldots,u_n)$ is ignored. Moreover, when we write $\rho(u_1,\ldots,u_n)$ it is implicitly assumed that the $n$-th order intensity function exists. 
%\eqref{e:1} is satisfied  

In particular, $\rho(u)$ is the usual intensity function. The point process is said to be (first-order) homogeneous if $\rho(u)=\rho_0$ is constant. For instance, this is the case if $S=\mathbb R^k$ and $X$ is stationary, i.e., its distribution is invariant under translations in $\mathbb R^k$; or if $S=\mathbb S^k$ and $X$ is isotropic, i.e., its distribution is invariant under rotations on $\mathbb S^k$. A similar example is not easy to specify if $S=L$, since there is no natural (transitive) group action on a linear network. 

Instead of the second order intensity function, one usually considers the pair correlation function (pcf) given by
\[g(u,v)=\frac{\rho(u,v)}{\rho(u)\rho(v)},\]
 setting $\frac{a}{0}=0$ for $a\ge0$, and one usually assumes that
 \begin{equation}\label{e:iso-asump}
 g(u,v)=g_0(d(u,v))
 \end{equation}
 is isotropic. For a Poisson process, 
 \begin{equation}\label{e:Poisson-rho-n}
 \rho(u_1,\ldots,u_n)=\rho(u_1)\cdots\rho(u_n),
 \end{equation}
  so $g=1$. One often interprets $g_0(t)>1$ as repulsion/inhibition and $g_0(t)<1$ as attraction/clustering for point pairs at distance $t$ apart but one should be careful with this interpretation as $t$ grows. %if the distance $u$ and $v$ are `distant apart'. 

The specific models in this paper are typically attractive ($g>1$) and satisfies the following stronger property: For $n=2,3,\ldots$, $i=1,\ldots,n-1$, and any pairwise distinct $u_1,\ldots,u_n\in S$,
\begin{equation}\label{e:pos-cor}
\rho(u_1,\ldots,u_n)\ge \rho(u_1,\ldots,u_i)\rho(u_{i+1},\ldots,u_n),
\end{equation}
or equivalently, for any pairwise disjoint, bounded, Borel sets $A_1,A_2,\ldots\subset S$, 
\[\mathbb E\left[N(A_1)\cdots N(A_n)\right]\ge \mathbb E\left[N(A_1)\cdots N(A_i)\right]\mathbb E\left[N(A_{i+1})\cdots N(A_n)\right].\]
So \eqref{e:pos-cor} means that the counts $N(A_1),N(A_2),\ldots$ are positively correlated at all orders, and for brief we shall say that $X$ is positively correlated at all orders.  

%FLG.\ TEKST OG DE NAESTE SECTIONS FLYTTES ELLER SKRIVES OM, SAA DER ER MERE FOKUS PAA LINEARE NETVAERK... 
Non-parametric estimation of $\rho$ and $g$ are discussed in e.g.\ \citet{spatstat} and the references therein; see also \citet{ShawEtAl} when $S=\mathbb R^k$; \citet{LawrenceEtAl} when $S=\mathbb S^k$; and \citet{RakshitEtAl} and \citet{RakshitEtAl19} when $S=L$. 
For non-parametric estimation of the pcf, usually
kernel methods are used under the assumption \eqref{e:iso-asump} of isotropy (but we do not need to assume that $X$ is (first-order) homogeneous).  
Since kernel methods are sensible to the choice of bandwidth, 
%as discussed in more detail in Sections~\ref{s:3.2.1} and \ref{s:3.2.3}, 
popular alternatives  are given by 
%straightforward non-parametric 
estimators %(which avoid using kernel methods) 
of the $K$-function % accumulated versions of $g$ as defined 
defined in 
%\eqref{e:K1}--\eqref{e:K3}
 \eqref{e:K1new}--\eqref{e:K2new} below. On the other hand, for the specific parametric models in this paper, we have simple expressions for $g_0$ but not for $K$,
 and
 since the $K$-function is an accumulated version of $g_0$, it may be harder to interpret.
 % (we return to this at the end of Section~\ref{s:3.2.3}).
 More details on these points are given in Sections~\ref{s:3.2.1} and \ref{s:3.2.3} below.
 
%For $u\in S$, I consider below mean values of $X\setminus\{u\}$ conditioned on $u\in X$; formally, this will be expectations with respect to the reduced Palm distribution of $X$ at the point $u$, which is well-defined for $\nu$-almost all $u\in S$, see e.g.\ \citet{}, \citet{},...  
%Moreover, $1(\cdot)$ denotes the indicator function. 

\subsubsection{$K$-functions in the classical cases of Euclidean spaces and spheres}\label{s:3.2.1}

Let either $(S,d)=(\mathbb R^k,\|\cdot\|)$ and $\nu$ be Lebesgue measure or $(S,d)=(\mathbb S^k,d_{\mathcal G})$ and $\nu$ be surface measure, cf.\ cases (a)--(b) in Section~\ref{s:new-setting}. Suppose again that the pcf %$g(u,v)=g_0(d(u,v))$ 
is isotropic. Following \citet{BMW}, \citet{LawrenceEtAl}, and \citet{MollerRubak}, we say that $X$ is second-order intensity reweighted stationary (SOIRS) if $S=\mathbb R^k$ and second order intensity reweighted
isotropic (or pseudo/correlation isotropic) if $S=\mathbb S^k$. 
This property is satisfied if $S=\mathbb R^k$ and $X$ is stationary and isotropic or if $S=\mathbb S^k$ and $X$ is isotropic. 

Following again \citet{BMW}, \citet{LawrenceEtAl}, and \citet{MollerRubak}, we define
the (inhomogeneous) $K$-function for $t\ge0$ by %is for an arbitrary $u\in S$ given by
\begin{equation}\label{e:K1new}
K(t)=\int_{d(u,v)\le t} g_0(d(u,v))\,\mathrm d\nu(v)
\end{equation}
which does not depend on the choice of $u\in S$. 
So
\[K(t)={K_{\mathbb R^k}(t)=\sigma_{k-1} \int_0^t r^{k-1} g_0(r)\,\mathrm dr} \quad \text{if } S=\mathbb R^k \text{ and } t\ge0\]
where $\sigma_{k-1}=2\pi^{k/2}/\Gamma(k/2)$ is the surface area of the $(k-1)$-dimensional unit sphere, and
\[K(t)={K_{\mathbb S^k}(t)=\sigma_{k-1}\int_0^t g_0(\vartheta)\sin^{k-1}\vartheta\,\mathrm d\vartheta} \quad \text{if } S=\mathbb S^k \text{ and } 0\le t\le\pi.\]
 Thus, for $k=1$, $K_{\mathbb R}(t)=K_{\mathbb S}(t)$ agree for every $t\in[0,\pi]$.

%\subsubsection{The case $S=\mathbb R^k$}\label{s:3.2.1}
%
%Suppose $S=\mathbb R^k$ and $g(u,v)=g_1(u-v)$ depends only on the lag $u-v$. Then $X$ is said to be second-order intensity reweighted stationary \citep[SOIRS, cf.][]{BMW}. For example, $X$ is SOIRS if it is stationary. For any bounded Borel set $A\subset\mathbb R^k$ of positive volume $|A|$ and $t\ge0$,
%the $K$-function is defined by
%%\begin{align}
%%K_{\mathbb R^k}(t)
%%&= %\frac{1}{|W|}\int_W\mathbb E\left[\sum_{v\in X\setminus\{u\}}1(\|u-v\|\le t)/\rho(v)\,\bigg |\,u\in X\right]\,\mathrm du\label{e:K1a}\\
%%\frac{1}{|W|}\mathbb E\sum_{u\in X_W}\sum_{v\in X\setminus\{u\}}\frac{1(\|u-v\|\le t)}{\rho(u)\rho(v)}\label{e:K1a}\\
%%&=\int_{\|u\|\le t}g_0(u)\,\mathrm du
%%\label{e:K1}
%%\end{align}
%\begin{equation}K_{\mathbb R^k}(t)
%=\frac{1}{|A|}\mathbb E\sum_{u\in X_A}\sum_{v\in X\setminus\{u\}}\frac{1(\|u-v\|\le t)}{\rho(u)\rho(v)}
%=\int_{\|u\|\le t}g_1(u)\,\mathrm du.
%\label{e:K1}
%\end{equation}
%%where $1(\cdot)$ denotes the indicator function.
%More precisely, this is the extension of Ripley's $K$-function \citep{ripley76,ripley77} for the stationarity case to the inhomogeneous $K$-function introduced in \citet{BMW} -- `inhomogeneous $K$-function' since $X$ does not need to be first-order homogeneous in \eqref{e:K1}. 
%If $g(u,v)=g_0(\|u-v\|)$ is isotropic, then $X$ is SOIRS and
% \begin{equation}K_{\mathbb R^k}(t)=\sigma_{k-1} \int_0^t r^{k-1} g_0(r)\,\mathrm dr\end{equation}
% where $\sigma_{k-1}=2\pi^{k/2}/\Gamma(k/2)$ is the surface area of the $(k-1)$-dimensional unit sphere (cf.\ Appendix~\ref{a:A1}). 
%
%%Intuitively, if $u$ is a uniformly selected point in $W$ and one considers the intensity-reweighted distribution of $X\setminus\{u\}$ conditioned on $u\in X$, $K_{\mathbb R^k}(t)$ is the expected number of further points within distance $t$. If $X$ is stationary with intensity $\rho_0$, $X$ is SOIRS and for an arbitrary chosen $u\in\mathbb R^k$,
%%\[\rho_0K_{\mathbb R^k}(t)=\mathbb E\left[\sum_{v\in X\setminus\{u\}}1(\|u-v\|\le t)/\rho(v)\,\bigg |\,u\in X\right]\] 
%%can be interpreted as the mean number of further points in $X$ within Euclidean distance $t$ from a typical point in $X$, cf.\ \citet{}. 
% 
%\subsubsection{The case $S=\mathbb S^k$}\label{s:3.2.2} 
% 
%Suppose $S=\mathbb S^k$ and $g(u,v)=g_0(d_\mathcal{G}(u,v))$ is isotropic. Then $X$ is said to be second order intensity reweighted
%isotropic (or pseudo/correlation isotropic) and for any Borel set $A\subseteq\mathbb S^k$ of positive surface area and $0\le t\le\pi$,  
%the 
%$K$-function is in accordance with \eqref{e:K1} defined by
%%\begin{align}
%%K_{\mathbb S^k}(t)
%%&= %\frac{1}{\nu_k(W)}\int_W\mathbb E\left[\sum_{v\in X\setminus\{u\}}1(d_{\mathcal G}(u,v)\|\le t)/\rho(v)\,\bigg |\,u\in X\right]\,\mathrm d\nu_k(u) \label{e:K2a}
%%\frac{1}{\nu_k(W)}\mathbb E\sum_{u\in X_W}\sum_{v\in X\setminus\{u\}}\frac{1(d_{\mathcal G}(u,v)\|\le t)}{\rho(u)\rho(v)}
%% \label{e:K2a}\\
%%%&=\int_{d_\mathcal{G}(u,v)\le t} g_0(d_\mathcal{G}(u,v))\,\mathrm d\nu_k(v)\\
%%&=\sigma_{k-1}\int_0^t g_0(\vartheta)\sin^{k-1}\vartheta\mathrm d\vartheta. \label{e:K2}
%%\end{align}
%\begin{equation}
%K_{\mathbb S^k}(t)
%=\frac{1}{\nu_k(A)}\mathbb E\sum_{u\in X_A}\sum_{v\in X\setminus\{u\}}\frac{1(d_{\mathcal G}(u,v)\le t)}{\rho(u)\rho(v)}
%=\sigma_{k-1}\int_0^t g_0(\vartheta)\sin^{k-1}\vartheta\,\mathrm d\vartheta, \label{e:K2}
%\end{equation}
%cf.\ \citet{LawrenceEtAl} and \citet{MollerRubak}. Note that \eqref{e:K1} and \eqref{e:K2} agree if $k=1$.
%%where %$u\in\mathbb S^k$ is arbitrary chosen and 
%%the last equality follows from \eqref{e:a1} in Appendix~\ref{a:A1}. % [$K(t)$ does not depend on this choice). 
%%Thus $K_{\mathbb S^k}$ has a similar interpretation as $K_{\mathbb R^k}$. 
%%In particular,  
%Isotropy of $X$ implies that the pair correlation function is isotropic and hence that $X$ is correlation isotropic.
%%for an arbitrary chosen point $u\in\mathbb S^k$,
%%\[\rho_0K_{\mathbb S^k}(t)=\mathbb E\left[\sum_{v\in X\setminus\{u\}}1(d_{\mathcal G}(u,v)\|\le t)/\rho(v)\,\bigg |\,u\in X\right]\] 
%%is the mean number of further points in $X$ within geodesic distance $t$ from a typical point in $X$, cf.\ \citet{}. REF UD HER??

\subsubsection{$K$-functions in the case of linear networks}\label{s:3.2.3}
%\subsubsection{The case $S=L$}\label{s:3.2.3}

Suppose $S=L$ and $g(u,v)=g_0(\delta(u,v))$ where $\delta$ is a metric on $L$ (since we consider derivatives below, it is convenient to switch from the previous notation $d$ to $\delta$). Then $X$ is said to be $\delta$-correlated, cf.\ \citet{RakshitEtAl}; if  $\delta=d_{\mathcal G}$, $X$ is also said to be second-order reweighted pseudostationary \citep{AngEtAl}.
%\citep[][said $X$ to be second-order reweighted pseudostationary when $\delta=d_{\mathcal G}$]{AngEtAL}. 
%Typically for a Cox process model (though there are exceptions), $g_0> 1$ is a decreasing function, and so $g_0(d_{\mathcal G}(u,v))\le g_0(d_{\mathcal R}(u,v))$. This means that for two such Cox models with the same intensity function and the same $g_0$-function but using different metrics $\delta=d_{\mathcal G}$ or $\delta=d_{\mathcal R}$, the pair correlation functions indicate more clustering for the model using the resistance metric. On the other hand, typically for a DPP model, $g_0< 1$ is an increasing function, so $g_0(d_{\mathcal G}(u,v))\ge g_0(d_{\mathcal R}(u,v))$ and comparing two DPP models the same intensity function and the same $g_0$-function but using different metrics $\delta=d_{\mathcal G}$ or $\delta=d_{\mathcal R}$, the pair correlation functions indicate more repulsion for the model using the resistance metric. Furthermore, when such Cox and DPP models are specified by an isotropic covariance function, larger classes of parametric models for isotropic covariance function are available with respect to the resistance metric than the geodesic metric, cf.\ Section~\ref{s:2.3}.  
%Therefore, using the resistance metric rather than the geodesic metric in such situations may allow for more flexibility when modelling clustering or repulsiveness in a point pattern dataset. MOVE THIS DISCUSSION TO A FINAL SECTION??? 
Following \citet{RakshitEtAl} and defining $R=\inf_{u\in L}\sup_{v\in L}d(u,v)$, 
the $K$-function is given by
\begin{equation}\label{e:K2new}
K(t)=K_L(t)=\int_0^t g_0(r)\,\mathrm dr,\quad 0\le t\le R.
\end{equation}
%Note that $K_L=K_{\mathbb R}/2$ 
%if $L$ is a straight line segment which is broken into $m$ %$m\ge2$ 
%line segments.
Note that $K$ depends only on $g_0$, but $g$ depends on both $g_0$ and $\delta$.
We have $K_L=K_{\mathbb R}/2$ 
if $L$ is a straight line segment which is broken into $m$ line segments. If $X$ is a Poisson process, then $K(t)=t$.

Non-parametric estimation of $K$ was carefully studied in \citet{RakshitEtAl} under the following technical assumption for the metric.  
 % when assuming $S=L$ and $X$ is $\delta$-correlated. 
Suppose $\delta$ is regular, meaning that for every $u\in L$, 
$\delta(u,v)$ is a continuous function of $v\in L$ and 
there is a finite set $N\subset L$ such that for $i=1,\ldots,m$ and all $v\in L_i\setminus N$, the Jacobian 
\begin{equation}\label{e:Jacobian}
J_\delta(u,v)=|(\mathrm d/\mathrm dt) \delta(u,v)|
\end{equation}
exists and is non-zero where $t=\|v-a_i\|$. 
% % except a $\nu_L$-nullset, 
%%the following conditions (A)--(B) are satisfied. 
%\begin{itemize}
%\item $\delta(u,v)$ is a continuous function of $v\in L$, 
%\item there is a finite set $N\subset L$ such that for $i=1,\ldots,m$ and all $v\in L_i\setminus N$, the Jacobian $J_\delta(u,v)=|(\mathrm d/\mathrm dt) \delta(u,u_i(t))|$ %|(\partial/\partial v)\delta(u,v)|$ 
%exists and is non-zero where $v=u_i(t)$. 
%%$(\partial/\partial v)\delta(u,v)=(\mathrm d/\mathrm dt) \delta(u,u_i(t))$ (this partial derivative is unique up to a sign, so $J_\delta(u,v)$ does not depend on the chosen arc length parametrization of the line segment $L_i$, cf.\ Section~\ref{s:2.1}).
%\end{itemize}
%%(This definition of regularity is slightly different from the one given in \citet{RakshitEtAl} where all $u\in L$ and all $v\in L\setminus N$, with $N$ a finite set,
%% but in fact only (A)--(B) will be needed when establishing the last equality in \eqref{e:K3} below.)
 Both $d_\mathcal{G}$ and $d_\mathcal{R}$ are regular, where $J_{d_\mathcal{G}}=1$ and a useful expression for the calculation of $J_{d_\mathcal{R}}$ is given in the following corollary which follows immediately from Proposition~\ref{p:calc}.
 
 
% with $N=\cup_{i=1}^m\{a_i,b_i\}$ given by the set of line segment endpoints. 
% %Since $d_{\mathcal R}\le d_{\mathcal G}$, it follows that $J_{d_{\mathcal R}}\le J_{d_\mathcal{G}}=1$, where
%We have $J_{d_\mathcal{G}}=1$ and a useful expression for the calculation of $J_{d_\mathcal{R}}$ is given in the following proposition. %Proposition~\ref{p:weight} below.
%% Appendix~\ref{a:A2}, cf.\ \eqref{e:der-d-R}. 

\begin{cor}\label{c:weight}
For all $u\in L_j$ and $v\in L_i$ with $u\not=v$, %so that $i=j$ and $u\not=v$ if $L_i\cap L_j\not=\emptyset$, 
using a notation as in Proposition~\ref{p:calc}, we have
 \begin{align}
\frac{\mathrm d}{\mathrm d t}d_{\mathcal{R}}(u,v) =
%=&\,-2\frac{\|v-b_i\|}{l_i^2}\Sigma(a_i,a_i)+2\frac{\|v-a_i\|}{l_i^2}\Sigma(b_i,b_i)+2\frac{\|v-b_i\|-\|v-a_i\|}{l_i^2}\Sigma(a_i,b_i)\nonumber\\
%&+2\frac{\|u-b_j\|}{l_jl_i}\Sigma(a_j,a_i)-2\frac{\|u-b_j\|}{l_jl_i}\Sigma(a_j,b_i)
%+2\frac{\|u-a_j\|}{l_jl_i}\Sigma(b_j,a_i)-2\frac{\|u-a_j\|}{l_jl_i}\Sigma(b_j,b_i)\nonumber\\
%&+\frac{\|v-b_i\|-\|v-a_i\|}{l_i}+2\frac{\|u-a_i\|}{l_i}1(u\in L_i,\,\|u-a_i\|\|v-b_i\|<\|u-b_i\|\|v-a_i\|)\nonumber\\
%&-2\frac{\|u-b_i\|}{l_i}1(u\in L_i,\,\|u-a_i\|\|v-b_i\|>\|u-b_i\|\|v-a_i\|). 
\begin{cases}
2A_i(t-s)+1 & \text{if }i=j,\ t> s,\\
2A_i(t-s)-1 & \text{if }i=j,\ t< s,\\
2A_it+B_{ij}(s) & \text{if }i\not=j.
\end{cases}
\label{e:der-d-R} %-\frac{1(i\not=j)}{l_i}\|v-a_i\|.
\end{align}
%Moreover, $d_{\mathcal{R}}(u,u_i(t))$ is a concave function of $t$.
\end{cor}
%\begin{proof} %Note that if $u\in L_i$ then $t\not=\|u-a_i\|$ because $u\not= v$.
%Since $u_i(t)=(1-t/l_i)a_i+(t/l_i)b_i$, $\|v-a_i\|=t$, and $\|v-b_i\|=l_i-t$, we get
%\[ %\frac{\partial}{\partial v}\|v-b_i\|^2=
%\frac{\mathrm d}{\mathrm dt}\|u_i(t)-b_i\|^2=2(t-l_i)=-2\|v-b_i\|,\quad \frac{\mathrm d}{\mathrm dt}\|u_i(t)-b_i\|=-1.\]
%%and
%%\[ %\frac{\partial}{\partial v}\|v-b_i\|=
%%\frac{\mathrm d}{\mathrm dt}\|u_i(t)-b_i\|=-1. \]
%From this and \eqref{e:d-R} we easily obtain \eqref{e:der-d-R} since %a straightforward calculation gives 
%\begin{align}
%%\frac{\partial}{\partial v}d_{\mathcal{R}}(u,v)=&\,\frac{\mathrm d}{\mathrm d t}d_{\mathcal{R}}(u,u_i(t))\nonumber\\
%\frac{\mathrm d}{\mathrm d t}d_{\mathcal{R}}(u,u_i(t))=&\,\frac{\mathrm d}{\mathrm d t}\Bigg\{
%\frac{(l_i-t)^2}{l_i^2}\Sigma(a_i,a_i)+\frac{t^2}{l_i^2}\Sigma(b_i,b_i)+\frac{2t(l_i-t)}{l_i^2}\Sigma(a_i,b_i)\nonumber\\
%&-\frac{2\|u-b_j\|(l_i-t)}{l_jl_i}\Sigma(a_j,a_i)-\frac{2\|u-b_j\|t}{l_jl_i}\Sigma(a_j,b_i)\nonumber\\
%&-\frac{2\|u-a_j\|(l_i-t)}{l_jl_i}\Sigma(b_j,a_i)-\frac{2\|u-a_j\|t}{l_jl_i}\Sigma(b_j,b_i)\nonumber\\
%&+\frac{t(l_i-t)}{l_i}-2\frac{1(u\in L_i)}{l_i}\min\{\|u-a_i\|(l_i-t),\|u-b_i\|t\}
%%\frac{1(i=j)}{l_i}\left[t(l_i-t)-2\min\{\|u-a_i\|(l_i-t),\|u-b_i\|t\}\right]+\frac{1(i\not=j)}{l_i}t(l_i-t)
%\Bigg\}.\nonumber
%\end{align}
%%whereby  easily follows.
%%Furthermore,
%%\[\left(\frac{\mathrm d}{\mathrm d t}\right)^2 d_{\mathcal{R}}(u,u_i(t))=\frac{2}{l_i^2}\left[\Sigma(a_i,a_i)+\Sigma(b_i,b_i)-2\Sigma(a_i,b_i)\right]-\frac{2}{l_i}=\frac{2}{l_i}\left[\frac{d_V(a_i,b_i)}{l_i}-1\right]\le0,\]
%%so $d_{\mathcal{R}}(u,u_i(t))$ is a concave function of $t$. Thereby we obtain \eqref{e:min}. 
%\end{proof} 


%Now, %define $R=\inf_{u\in L}\sup_{v\in L}d(u,v)$, and 
%The interpretation of $K_L$ is a bit harder than that of $K_{\mathbb R^k}$ and $K_{\mathbb S^k}$, but 
%Rakshit, Nair, Baddeley (2017), Spat. Stat. in order to obtain \eqref{e:K3}.

Some final remarks are in order. 

For $u\in L$ and $0\le t\le R$, define
\[w_\delta(u,t)=1\bigg/\sum_{v\in L:\, \delta(u,v)=t}1/J_\delta(u,v).\]
%\citet{RakshitEtAl} remark that $J_{d_{\mathcal G}}\le1$ and hence that $w_{\mathcal R}(u,t)\le w_{\mathcal G}(u,t)=1$. 
This is a weight which accounts for the geometry of the linear network when shifting from arc length measure on $L$ to Lebesgue measure on the positive half-line \citep[][Propositions 1 and 2]{RakshitEtAl}. 
For $\delta=d_{\mathcal{G}}$, we have 
$w_{\mathcal G}(u,t)=1/\#\{v\in L\,|\,\delta(u,v)=t\}$, since $J_{d_\mathcal{G}}=1$, and  
%$w_{\mathcal R}(u,t)\le w_{\mathcal G}(u,t)=1/\#\{v\in L\,|\,\delta(u,v)=t\}$, since $d_{\mathcal R}\le d_{\mathcal G}$ and $J_{d_\mathcal{G}}=1$. 
for $\delta=d_{\mathcal{R}}$, once the matrix $\Sigma$ from Section~\ref{s:res-met} has been calculated, $w_{d_{\mathcal{R}}}$ is quickly calculated from \eqref{e:der-d-R}.

It follows from \eqref{e:1}, \eqref{e:K2new}, and \citet[][Equation (8)]{RakshitEtAl}  that 
for any Borel set $A\subseteq L$ of positive arc length measure, %$u\in L$ 
%and $0\le t\le R$, %the $K$-function is defined by
\begin{equation}
K(t)=
\frac{1}{\nu_L(A)}\mathbb E\sum_{u\in X_A}\sum_{v\in X\setminus\{u\}}\frac{1(\delta(u,v)\le t)w_\delta(u,\delta(u,v))}{\rho(u)\rho(v)}.
%=\int_0^t g_0(r)\,\mathrm dr . 
\label{e:K3}
\end{equation}
Non-parametric estimators of $K$ are based on omitting the expectation symbol in \eqref{e:K3}, possibly after elaborating on the right hand side in \eqref{e:K3} in order to realize how correction factors can be included in order to adjust for edge effects, cf.\ \citet{RakshitEtAl}.
 
%We can rewrite \eqref{e:K3} 
In terms of Palm probabilities, for $d_L$-almost all $u\in L$ with $\rho(u)>0$,
\[K(t)=\frac{1}{|L|}\mathbb E \sum_{v\in X_u}\frac{1(\delta(u,v)\le t)w_\delta(u,\delta(u,v))}{\rho(v)}.\]
In general
%\citep[see again][]{RakshitEtAl}, 
the weight makes it hard to interpret this expression of $K$. %, except that $K$ is an accumulated counterpart of $g_0$.  
 
%For the calculation of the weight, % $w_\delta$ in \eqref{e:K3}, which also appears in non-parametric estimators of $K_L$, 
%the calculation of the Jacobian 
%%$(\mathrm d/\mathrm dt) \delta(u,u_i(t))$
%%the Jacobian $J_\delta(u,v)=|(\mathrm d/\mathrm dt) \delta(u,u_i(t))|$ %|(\partial/\partial v)\delta(u,v)|$ 
%is needed. %$(\partial/\partial v)\delta(u,v)$ is needed 
%%for all $u\in L$ and all $v=u_i(t)\in L_i\setminus N$ with $i=1,\ldots,m$, where $N\subset L$ is finite.
%As shown in Proposition~\ref{p:weight}, if $\delta=d_{\mathcal{R}}$, we can take $N=V$ and
%once the matrix $\Sigma$ from Section~\ref{s:res-met} has been calculated, the %resistance metric and the 
%Jacobian and hence the weight %$J_\delta(u,v)=|(\mathrm d/\mathrm dt) \delta(u,u_i(t))|$ are 
%are quickly calculated from %\eqref{e:d-R} and 
%\eqref{e:der-d-R}. % below.

\subsection{Estimation and model checking}\label{s:est-check} %ther functional summaries}

%TBD....
%
%%\subsubsection{Remarks about the pair correlation and $K$-functions and other functional summaries used for estimation and model checking}\label{s:remarks} %$K_{\mathbb R^k}$, $K_{\mathbb S^k}$, and $K_{L}$}
%
%SHORTEN AND MODIFY THIS OLD TEXT WHEN Cronie et al.\ have replied to my emails....
%
%The interpretation and the derivation of non-parametric edge-corrected estimators of $K_{\mathbb R^k}$, $K_{\mathbb S^k}$, and $K_{L}$ 
%are based on elaborating on the right hand expressions of the first equalities in \eqref{e:K1}--\eqref{e:K3}.
%The references in Sections~\ref{s:3.2.1}--\ref{s:3.2.3} provide this interpretation in terms of conditional (Palm) distributions, which is easiest to understand if $X$ is homogeneous and $S=\mathbb R^d$ or $S=\mathbb S^d$: Then, for $K=K_{\mathbb R^k}$ or $K=K_{\mathbb S^k}$,  
% $\rho_0K(t)$ is the mean number of further points within distance $t$ from a typical point in $X$. The interpretation of $K_L$ is harder because of the weight in \eqref{e:K3}.
%For the exact expressions of the edge-corrected non-parametric estimators, see again the  
% references in Sections~\ref{s:3.2.1}--\ref{s:3.2.3}; see also \citet{spatstat} and \citet{ShawEtAl}. 
 
For parametric families of Cox point process models the most common estimation methods are based on the intensity, pair correlation, or $K$-functions using either minimum contrast estimation, composite likelihood, or Palm likelihoods, see \citet[][]{MW07,MW17} and the references therein. For the analyses in Section~\ref{s:appl}, we use minimum contrast estimation for fitting the Cox process models in Section~\ref{s:models} to various datasets in Section~\ref{s:appl} using the pcf, i.e., if $g_0$ depends on a parameter, we estimate this parameter by minimizing the integral
\begin{equation}\label{e:contrast}
D(g_0,\hat g_0) = \int_{a_1}^{a_2} \left|g_0(t)^q-\hat{g}_0(t)^q\right|^p \mathrm{d}t
\end{equation}
where $\hat g_0$ is a non-parametric estimate of $g_0$ \citep[see e.g.\ (31) in][]{RakshitEtAl} and $0\leq a_1 < a_2$ and $p,q>0$ are user-specified values (see Section~\ref{s:appl}). The models in Section~\ref{s:models} all have nice expressions of the pcf but not of the $K$-function. Moreover, the models used for the data analyses in Section~\ref{s:appl} include a parameter for the intensity function, which $g_0$ does not depend on, and this parameter is estimated by a simple moment method; in the simplest case homogeneity is assumed and the intensity is estimated by $\hat{\rho}={n(x)}/{|L|}$. 

For model checking other functional summary statistics are needed when $\rho$, $g$, or $K$ and their corresponding non-parametric estimators have been used for estimation. 
For the case $S=L$, \citet{CronieEtAl} suggested analogies to the so-called $F,G,J$-functions %and their main result
%$F$ (the empty space), $G$ (the nearest-neighbour distribution), and $J$-functions
% 
\citep[introduced in][when $S=\mathbb R^k$]{Inhom-J-fct} which account for the geometry of the linear network. \citet[][]{CronieEtAl} 
showed that their $F,G,J$-functions make good sense under a certain condition called intensity reweighted moment pseudostationarity
(IRMPS): $X$ is IRMPS if $\inf\rho>0$ and $\delta$ is a regular metric on $L$ such that for $n=2,3,\ldots$, any pairwise distinct $u_1,\ldots,u_n\in L$, and any $u\in L$, $g(u_1,\ldots,u_n)=\rho^{(n)}(u_1,\ldots,u_n)/[\rho(u_1)\cdots\rho(u_n)]$ is of the form
\begin{equation}\label{e:IRMPS}
g(u_1,\ldots,u_n)=g_0(\delta(u,u_1),\ldots,\delta(u,u_n))
\end{equation}
for some function $g_0$. This condition is satisfied if  $X$ is either  a Poisson process or  %$L$ is a line segment and 
 a log Gaussian Cox process (LGCP) with a stationary pair correlation function (which is usually not a natural assumption, cf.\ Section~\ref{s:LGCP}). Apart from these examples 
\citet[][]{CronieEtAl} did not %(correctly) 
 verify any other cases of models where IRMPS is satisfied \citep[in Section~\ref{s:LGCP} we correct some mistakes in][]{CronieEtAl}. At least
%, and in our opinion IRMPS  is a very restrictive assumption. % (see again Section~\ref{s:LGCP}). 
\citet[][]{CronieEtAl} demonstrated % for examples of planar linear networks 
the practical usefulness of their empirical estimator of the $J$-function for both a Poisson process, a simple
sequential inhibition (SSI) point process, and a LGCP. We show in Section~\ref{s:LGCP} that IRMPS is in general not satisfied for the LGCP. For the SSI point process in \citet[][]{CronieEtAl} it is hard to evaluate $g(u_1,\ldots,u_n)$ and hence to check the assumption of IRMPS. 

Alternatively, \citet{HeidiMe} introduced three purely empirical summary functions obtained by modifying the empirical $F,G,J$-functions for (inhomogeneous)
point patterns on a Euclidean space to linear
networks. Briefly, the modification consists of replacing the Euclidean space with
the linear network, introducing the shortest path distance instead of the Euclidean
distance, and adapting the notion of an eroded set to linear networks. \citet{HeidiMe} demonstrated the usefulness of their empirical $F,G,J$-functions for model checking although underlying theoretical functions are missing. We also use these functions for the data analyses in Section~\ref{s:appl}. Specifically, we will consider a concatenation of the three functions and validate a fitted model using a 95\% global envelope (i.e., a 95\%  confidence region for the concatenated function) and a $p$-value obtained by the 
global envelope test (based on the extreme rank length) as described in \citet{Mylly1}, \citet{Mylly2}, and \citet{Mylly3}. 
For the dendrite spine locations datasets analysed in \citet{HeidiMe} it is concluded that the results based on such a global envelope test are consistent with the results obtained if instead the summary functions from {\citet[][]{CronieEtAl}} are used.  





%   
%
% such as the $F$, $G$, and $J$-functions whose definitions depend on the space $S$ and the assumed invariance properties: Let $u\in S$ be an arbitrary fixed position in $S$. If $S=\mathbb R^k$ and $X$ is stationary or if $S=\mathbb S^k$ and $X$ is isotropic, the empty space function $F(t)$ is the probability that $X$ has no points within distance $t$ from $u$, the nearest-neighbour function $G(t)$ is the probability that $X$ has no further points within distance $t$ when conditioning on $u\in X$, and $J(t)=(1-G(t))/(1-F(t))$ provided $F(t)<1$; see e.g.\ \citet{BaddeleyEtAl}, \citet{LawrenceEtAl}, \citet{MollerRubak} and the references therein. 
%Inhomogeneous versions of these functions are studied in \citet{Inhom-J-fct} when $S=\mathbb R^k$, in \citet{LawrenceEtAl} when $S=\mathbb S^k$, and in \citet{CronieEtAl} when $S=L$.

%To validate a fitted model when considering one or more empirical summary functions, the most popular method uses simulations to obtain so-called global envelopes (which are confidence regions) and $p$-values for global envelope tests based on the extreme rank length as described in \citet{Mylly1}, \citet{Mylly2}, and \citet{Mylly3}. In the data analysis in section~\ref{s:appl} we use the global envelope test based on a concatenation of the empirical $F,G,J$-functions introduced in \citet{HeidiMe}.


%In general, for the calculation of $K_{\mathbb R^k}$, $K_{\mathbb S^k}$, and $K_{L}$ given by the integrals in \eqref{e:K1}--\eqref{e:K3}, numerical methods are needed, though for some special models closed form expressions are known, including the case of a Poisson process where  $g=1$. Numerical methods are usually also needed for the $F$, $G$, and $J$-functions when $S=\mathbb R^k$ and $X$ is stationary or if $S=\mathbb S^k$ and $X$ is isotropic, where again the Poisson process is an exception. This is illustrated by the expressions of $F$ and $G$ in \eqref{e:F} and \eqref{e:G} below. %Section~\ref{s:LGCP}.

%\section{Covariance functions}\label{s:covfct} % Isotropic covariance functions in Euclidean space and beyond}\label{s:covfct}
%
%Let $c$ be a complex covariance function defined on the product space $S\times S$, which means the following. 
%For $x,y\in\mathbb R$, let $\overline{x+iy}=x-iy$ denote complex conjugation where $i^2=-1$. Then $c$ is Hermitian, that is, $c(u,v)=\overline{c(v,u)}$ for all $u,v\in S$, and $c$ is positive semi-definite, that is, $\sum_{j,\ell=1}^n z_j\overline{z_\ell}c(u_j,u_\ell)\ge0$ for all $z_1,\ldots,z_n\in\mathbb C$, all pairwise distinct $u_1,\ldots,u_n\in S$, and $n=1,2,\ldots$. 
%
%Usually in this paper $c$ is real. Then the conditions above reduce to that $c$ is symmetric and $\sum_{j,\ell=1}^n x_j{x_\ell}c(u_j,u_\ell)\ge0$ for all $x_1,\ldots,x_n\in\mathbb R$, $u_1,\ldots,u_n\in S$, and $n=1,2,\ldots$. The case where $c$ is complex becomes relevant when studying DPPs. 
%
%Often $c$ is assumed to be isotropic, that is, $c(u,v)=c_0(d(u,v))$ depends only on the metric distance $d(u,v)$ for $u,v\in S$, where $c_0$ is also called a covariance function. Then $c$ is real, 
%$\sigma^2=c_0(0)$ is the variance, and for $\sigma>0$, $r_0(r)=c_0(r)/\sigma^2$ defines an isotropic correlation function. 
%
%Suppose $S=\mathbb R^k$. Then $c$ is stationary if $c(u,v)$ depends only on the lag $u-v$. This is the case if $c$ is isotropic or just geometric anisotropic, meaning that $c(u,v)=c_0(\|D\mathcal O(u-v)\|)$ where $D$ is a $k\times k$ diagonal matrix with positive diagonal elements and $\mathcal O$ is a $k\times k$ rotation matrix. 
%%Isotropy is a natural condition, since any stationary and anisotropic covariance function can be obtained from some isotropic covariance function using some rotation followed by some rescaling, see e.g.\ \citet{Goovaerts}. 
%Examples of popular parametric models for isotropic correlation functions are given in ~\ref{tab:covariances-sphere} below. 
%See \citet{Schoenberg38} and \citet{Gneiting01} for detailed studies of these and other continuous isotropic covariance functions. 
%
%\begin{table}
%    \centering
%    \begin{tabular}{l||l|l}
%        Model & Correlation function $r_0(r)$ & Range of shape and smoothness parameters\\Table
%        \hline
%        \hline
%        Powered exponential & $\exp\left(-r^\alpha/\phi\right)$ & $\alpha\in(0,2]$ if $S=\mathbb R^k$; $\alpha\in(0,1]$ if $S=\mathbb S^k$ or $S=L$\\
%        Mat{\'e}rn & $\frac{2}{\Gamma(\nu)}\left( \frac{r}{2\phi} \right)^{\nu} K_{\nu}\left(\frac{r}{\phi}\right)$ &  $\nu>0$  if $S=\mathbb R^k$; $0< \nu \leq \frac{1}{2}$ if $S=\mathbb S^k$ or $S=L$\\
%        Generalized Cauchy & $(1 + (\frac{r}{\phi})^{\alpha})^{-\tau/\alpha}$ & $\tau > 0$; $\alpha \in (0,2]$ if $S=\mathbb R^k$; $\alpha \in (0,1]$ if $S=\mathbb S^k$ or $S=L$\\
%        Dagum & $1 - ((\frac{r}{\phi})^{\tau}/(1+(\frac{r}{\phi})^{\tau}))^{\frac{\alpha}{\tau}}$ & $\tau \in (0,2]$ if $S=\mathbb R^k$; $\tau \in (0,1]$ if $S=\mathbb S^k$ or $S=L$; $\alpha \in (0,\tau)$
%    \end{tabular}
%    \caption{Four parametric models for an isotropic correlation function $r_0(r)$. Here, $\Gamma$ is the gamma function, $K_{\nu}$ is the modified Bessel function of the second kind, $\phi$ is a scale parameter, $\tau$ is a 
%shape parameter, and $\alpha$ and $\nu$ are smoothness parameters. The correlation functions are well-defined for $k=1,2,\ldots$ and all $\phi>0$ but the range of shape and smoothness parameters depends on the model and the space $S$. For $S=L$, distance is measured by the resistance metric.
%    % and $t_{+} := \max\{t,0\}$ for $t\in\mathbb R$. For the powered exponential, Mat{\'e}rn,  generalized Cauchy, Dagum, multiquadric, and %sine power models, $k\in \{1,2,\ldots\}$, whilst for the spherical, Askey, $C^2$-Wendland, and $C^4$-Wendland models, $k\in\{1,2,3\}$.
%    } 
%    \label{tab:covariances-sphere}
%\end{table}
%
%Further detailed studies of continuous isotropic covariance functions are given in \citet{Schoenberg} and \citet{gneiting2013strictly} when $S=\mathbb S^k$. Examples are seen in Table~\ref{tab:covariances-sphere}, where compared to the case $S=\mathbb R^k$, the ranges of shape and smoothness parameters are more restrictive.
%%sphere: Mat{\'e}rn model (Gneiting)...
%%multiquadric covariance function 
%%which is isotropic and given by
%%\begin{equation}\label{e:multiquadric}
%%c_{(\sigma,\delta,\tau)}(r) = \sigma^2\left(\frac{(1 − \delta)^2}{1 + \delta^2 − 2\delta\cos r}\right)^\tau,\quad r \in[0,\pi],
%%\end{equation}
%%where $(\sigma,\delta,\tau)\in(0,\infty)\times(0,1)\times(0,\infty)$. 
%
%%\begin{table}
%%    \centering
%%    \begin{tabular}{l||l|l}
%%        Model & Correlation function $c_0(r)$ & Parameter range \\
%%        \hline
%%        \hline
%%        Powered exponential & $\exp\left(-r^\alpha/\phi\right)$ & $\alpha\in(0,1]$, $\phi>0$\\
%%        Mat{\'e}rn & $\frac{2}{\Gamma(\nu)}\left( \frac{r}{2\phi} \right)^{\nu} K_{\nu}\left(\frac{r}{\phi}\right)$ & $0< \nu \leq \frac{1}{2}$, $\phi>0$\\
%%        Generalized Cauchy & $(1 + (\frac{r}{\phi})^{\alpha})^{-\tau/\alpha}$ & $\phi > 0$, $\alpha \in (0,1]$, $\tau > 0$\\
%%        Dagum & $1 - ((\frac{r}{\phi})^{\tau}/(1+(\frac{r}{\phi})^{\tau}))^{\frac{\alpha}{\tau}}$ & $\phi > 0$, $\tau \in (0,1]$ \\
%%        multiquadric & $\left(\frac{(1-\delta)^2}{1+\delta^2-2\delta\cos r}\right)^\tau$ & $\delta\in(0,1)$, $\tau>0$\\
%%        Sine power & $1- \sin(r/2)^{\alpha}$ & $\alpha \in (0,2)$ \\
%%        Spherical & $(1 + \frac{1}{2}\frac{r}{\phi})(1-\frac{r}{\phi})_{+}^{2}$ & $\phi > 0$  \\
%%        Askey & $(1-\frac{r}{\phi})_{+}^{\tau}$ & $\phi > 0$, $\tau \geq 2$ \\
%%        $C^{2}$-Wendland & $(1 + \tau\frac{r}{\phi})(1-\frac{r}{\phi})_{+}^{\tau}$ & $\phi \in (0,\pi]$, $\tau \geq 4$ \\
%%        $C^{4}$-Wendland & $(1 + \tau\frac{r}{\phi} + \frac{\tau^{2}-1}{3}\frac{r^{2}}{\phi^{2}})(1-\frac{r}{\phi})_{+}^{\tau}$ & $\phi \in (0,\pi]$, $\tau \geq 6$ \\
%%    \end{tabular}
%%    \caption{Parametric models for an isotropic correlation function $r_0(r)$ on $\mathbb{S}^{k}$, where $0\le r\le\pi$, $\Gamma$ is the gamma function, $K_{\nu}$ is the modified Bessel function of the second kind, and $t_{+} := \max\{t,0\}$ for $t\in\mathbb R$. For the powered exponential, Mat{\'e}rn,  generalized Cauchy, Dagum, multiquadric, and sine power models, $k\in \{1,2,\ldots\}$, whilst for the spherical, Askey, $C^2$-Wendland, and $C^4$-Wendland models, $k\in\{1,2,3\}$.} 
%%    \label{tab:covariances-sphere}
%%\end{table}
%
%For $S=L$, many of the commonly used isotropic covariance functions, including those in Table~\ref{tab:covariances-sphere}, are valid with respect to the resistance metric for any graph with Euclidean edges, whilst they are only valid with respect to the geodesic metric in more special cases, cf.\ \citet{AnderesEtAl}. For the examples in Table~\ref{tab:covariances-sphere}, the range of the parameters are the same  when $S=\mathbb S^k$ or $S=L$.

%ALLE DE FLG. SECTIONS SKAL KORTES NED SAA FOKUS MERE ER PAA LINEARE NETVAERK...


\section{Cox processes driven by transformed Gaussian processes}\label{s:cox}

%Section~\ref{s:cox-general} treats Cox processes in a general setting, since moment and many other properties of Cox processes driven by transformed Gaussian processes do not depend on which space we consider. The dependence of the space occurs when we want to specify well-defined parametric models of covariance functions, simulate the underlying GPs, and use statistical inference procedures, which for the case $S=L$ are the topics of Sections~\ref{s:GP-cov-fct} and \ref{s:appl}. Section~\ref{s:models} considers specific models for Cox processes driven by transformed Gaussian processes.

%LGCP, ICP, and PCPP models are well-studied when $S=\mathbb R^k$ or $S=\mathbb S^k$, and most of their properties immediately extend to the case $S=L$ as discussed in the following. The differences are basically how to specify well-defined parametric models of covariance functions, to simulate the underlying GPs, and to use statistical inference procedures, which for the case $S=L$ are the topics of Sections~\ref{s:GP-cov-fct} and \ref{s:appl}.

\subsection{General setting}\label{s:cox-general}

%Consider again the general setting of Section~\ref{s:3.1} as used in Section~\ref{s:3.2} until Section~\ref{s:3.2.1}. 

Let $(S,d)$ be a metric space with reference measure $\nu$ % as considered in Section~\ref{s:3.1} 
and 
let $\Lambda=\{\Lambda(u)\,|\,u\in S\}$ be a non-negative stochastic process so that with probability 1, for any bounded Borel set $A\subseteq S$,  the measure $\xi(A)=\int_A\Lambda(u)\,\mathrm d\nu(u)$ is finite. Throughout this section we assume $X$ is a Cox process driven by $\Lambda$, that is, $X$ is a point process on $S$ which conditioned on $\Lambda$ is almost surely a Poisson process with intensity function $\Lambda$ (with respect to $\nu$). 
%In other words, with probability 1, for any bounded Borel set $A\subseteq S$, $N(A)|\Lambda$ is Poisson distributed with mean $\mu(A)$, and for any integer $n>0$, conditioned on $\Lambda$ and $N(A)=n$, the $n$ points in $X_A$ are independent and each point has a density proportional to $\Lambda$ with respect to $\nu$ restricted to $A$. 
Usually in applications $\Lambda$ is unobserved, and so if only one point pattern dataset is observed the Cox process $X$ is indistinguishable from the inhomogeneous Poisson process $X|\Lambda$  \citep[for a discussion of which of the two models is most appropriate, see][Section 5.1]{MW04}. 

Henceforth, assume $\mathbb E\Lambda(u)$ is a locally integrable function with respect to $\nu$. Then the Cox process $X$ is well-defined, since almost surely $\xi(A)<\infty$ for all bounded Borel sets $A\subseteq S$. % and has intensity function $\rho(u)=\mathbb E\Lambda(u)$.
By conditioning on $\Lambda$ it follows from \eqref{e:1} and \eqref{e:Poisson-rho-n} that
\begin{equation}\label{e:Cox-rho-n}
\rho(u_1,\ldots,u_n)=\mathbb E\left[\Lambda(u_1)\cdots\Lambda(u_n)\right].
\end{equation}

Let $W\subseteq S$ be a bounded Borel set with $\nu(W)>0$; we think of $W$ as an observation window. Then  $X_W$ is a Cox process driven by $\Lambda$ restricted to $W$, and $X_W$ has a density given by
\begin{equation}\label{e:densityCox}
f(\{u_1,\ldots,u_n\})=\mathbb E\left[\exp\left(\int_W (1- \Lambda(u))\,\mathrm d\nu(u)\right)\prod_{i=1}^n\Lambda(u_i)\right]\quad\mbox{for pairwise distinct }u_1,\ldots,u_n\in W
\end{equation}
with respect to $X_1\cap W$, where $X_1$ is the unit rate Poisson process, cf.\ Section~\ref{s:3.1}.
In particular, if $S\in\{\mathbb S^k,L\}$, $W=S$, and $u\in S$ with $\rho(u)>0$, then $X_u$ has density
\begin{equation}\label{e:PalmdensityCox} 
f_u(\{u_1,\ldots,u_n\})=\mathbb E\left[\exp\left(\int_W (1- \Lambda(u))\,\mathrm d\nu(u)\right)\frac{\Lambda(u)}{\rho(u)}\prod_{i=1}^n\Lambda(u_i)\right]
\end{equation}
with respect to $X_1$, cf.\ \eqref{e:Palmdensity}.
In general the densities in \eqref{e:densityCox} and \eqref{e:PalmdensityCox} are intractable because the expected values are difficult to evaluate. Instead we  
exploit \eqref{e:Cox-rho-n} for calculating the intensity and pair correlation functions and for making inference as demonstrated in the following sections. 
% see e.g.\ \citet{}...  

As pointed out in \citet{MW07} it is useful to write $\Lambda$ as
\begin{equation}\label{e:Lambda}
\Lambda(u)=\rho(u)\Lambda_0(u)
\end{equation}
where $\Lambda_0=\{\Lambda_0(u)\,|\,u\in S\}$ is a non-negative `residual' stochastic process with $\mathbb E\Lambda_0(u)=1$ whenever $\rho(u)>0$.
Then
\[\rho(u_1,\ldots,u_n)=\rho(u_1)\cdots\rho(u_n)\mathbb E\left[\Lambda_0(u_1)\cdots\Lambda_0(u_n)\right]\]
and %in particular
 \begin{equation}\label{e:g-Lambda0}
 g(u,v)=\mathbb E\left[\Lambda_0(u)\Lambda_0(v)\right].
 \end{equation}
%(provided the pair correlation function exists). 
Typically, it is only $\rho(u)$ which is allowed to depend on covariate information, whilst $\Lambda_0$ is considered to account for unobserved covariates or other effects which has not been successfully fitted by a Poisson process with intensity function $\rho$, see e.g.\ \citet{MW07} and \citet{Diggle}. 

Usually, assumptions on $\Lambda_0$ are imposed so that $g(u,v)=g_0(d(u,v))$ is isotropic; this is the case if $S=\mathbb R^k$ and
$\Lambda_0$ is stationary and isotropic, or if $S=\mathbb S^k$ and $\Lambda_0$ is isotropic. 
Furthermore, for specific models of $\Lambda_0$ it usually happens that 
\begin{equation}\label{e:all-orders}
\mathbb E\left[\Lambda_0(u_1)\cdots\Lambda_0(u_n)\right]\ge\mathbb E\left[\Lambda_0(u_1)\cdots\Lambda_0(u_i)\right]\mathbb E\left[\Lambda_0(u_{i+1})\cdots\Lambda_0(u_n)\right]
\end{equation}
for $n=2,3,\ldots$, $i=2,\ldots,n$, and all pairwise distinct $u_1,\ldots,u_n\in S$, meaning that
$X$ is positively correlated at all orders, 
 cf.\ \eqref{e:pos-cor}. 


%Apart from the model classes studied in Section~\ref{s:cox-cov} below, shot-noise Cox processes play a key role for statistical applications. A shot-noise Cox process (SNCP) with a homogeneous residual random field is given by
%\begin{equation}\label{e:SNCP}
%\Lambda_0(u)=\frac{1}{\rho_Y}\sum_{y\in Y}
%\end{equation}
%where 
%\citet studied SNCPs in detail when $S=\mathbb R^k$, and most of the results ... See also...  Clusters....

\subsection{Models}\label{s:models}

Consider a GP $Y=\{Y(u)\,|\,u\in S\}$ with mean function $\mu$ and covariance function $c$, and let $Y_1,\ldots,Y_h$ be independent copies of $Y$.
In the remainder of this paper we study the following models.
\begin{itemize}
\item $X$ is a log Gaussian Cox process (LGCP) if 
\begin{equation}\label{e:defLGCP}
\Lambda_0(u)=\exp(Y(u))
\end{equation}
and $\mu(u)=-c(u,u)/2$ for all $u\in S$. The latter condition is required 
 since we want $\mathbb E\Lambda_0(u)=1$.  
\item $X$ is an interrupted Cox process (ICP) if 
 $\mu=0$ and 
\begin{equation}\label{e:defInter}
\Lambda_0(u)=\Pi(u)(1+2c(u,u))^{h/2}
\end{equation}
for all $u\in S$ where we define $\Pi(u)=\exp(-\sum_{i=1}^h Y_i(u)^2)$ (this definition differs slightly from the one used in \cite{FredMe}, which includes a factor $1/2$ inside the exponential function). Since $\mathbb E\Pi(u)=(1+2c(u,u))^{-h/2}$, we have $\mathbb E\Lambda_0(u)=1$. 
\item  $X$ is a permanental Cox point process (PCPP) if $\mu=0$, $c(u,u)=1$, and % (so $c$ is a correlation function). If 
\begin{equation}\label{e:defPCPP}
\Lambda_0(u)=\frac{1}{h}\sum_{i=1}^h Y_i(u)^2
\end{equation}
for all $u\in S$. Since the sum in \eqref{e:defPCPP} is $\chi^2(h)$-distributed, $\mathbb E\Lambda_0(u)=1$.
\end{itemize}
In all cases, the distribution of $X$ is completely specified by $(\rho,c)$ and in the case of an ICP or PCPP the value of $h$. 
Note that the intensity function $\rho$ can be any non-negative locally integrable function with respect to $\nu$.


% (as used in Section~\ref{s:3.2} until Section~\ref{s:3.2.1}). 

%In the following, we assume \eqref{e:Lambda} and adapt the most common families of Cox processes when $S=\mathbb R^k$ or $S=\mathbb S^k$ to the case $S=L$, where the distribution of $\Lambda_0$ is specified by a covariance function $c$. Then the Cox process $X$ is specified by $(\rho,c)$ where the intensity function $\rho$ can be any non-negative locally integrable function with respect to $\nu$.

%\section{Cox processes specified by a covariance function}\label{s:cox-cov}

LGCP, ICP, and PCPP models are well-studied when $S\in\{\mathbb R^k,\mathbb S^k\}$, and most of their properties immediately extend to the case $S=L$ as discussed in the following. 
%The differences are basically how to specify well-defined parametric models of covariance functions, to simulate the underlying GPs, and to use statistical inference procedures, which for the case $S=L$ are the topics of Sections~\ref{s:GP-cov-fct} and \ref{s:appl}. 

\subsubsection{Log Gaussian Cox processes}\label{s:LGCP}

Let $X$ be a LGCP, cf.\ \eqref{e:defLGCP}. \citet{LGCP} studied LGCPs in the case $S=\mathbb R^k$, and \citet{CM} in the case $S=\mathbb S^k$. As in \citet{LGCP} and \citet{CoeurjollyEtAl}, we obtain for  the general setting of the space $S\in\{\mathbb R^k,\mathbb S^k,L\}$ the following results. 

 For any integer $n\ge2$ and any pairwise distinct $u_1,\ldots,u_n\in S$,
\begin{equation}\label{e:n-LGCP}
\rho(u_1,\ldots,u_n)=\rho(u_1)\cdots\rho(u_n)\exp\bigg(\sum_{1\le i<j\le n}c(u_i,u_j)\bigg).
\end{equation}
%which follows from the well-known expression of the moment generating function for a normal distribution.
In particular, the LGCP is determined by $\rho$ and 
\begin{equation}\label{e:pairLGCP}
g=\exp(c),
\end{equation} 
i.e., by its first and second order moment properties. Note that $g$ is isotropic if and only if $c$ is isotropic, and %if $c(u,v)=c_0(d(u,v))$ is isotropic, 
then $\rho(u_1,\ldots,u_n)$ depends only on the inter-point distances $d(u_i,u_j)$, $1\le i<j\le n$.  
Moreover, in most specific models, including those in Table~\ref{tab:covariances-sphere}, $c\ge0$ or equivalently $X$ is positively correlated at all orders, cf.\ \eqref{e:pos-cor} and \eqref{e:n-LGCP}. 
%$c>0$ and hence $g>1$ or more generally
%\[\rho(u_1,\ldots,u_n)>\rho(u_1,\ldots,u_i)\rho(u_{i+1},\ldots,u_n),\quad i=1,\ldots,n-1,\]
%which reflects clustering/aggregation in $X$. 

%Along similar lines as in \citet[][Theorem~1]{CoeurjollyEtAl}  % Coeurjolly et al. (2017, Theorem 1), 
%it follows that 
For any $u\in S$ with $\rho(u)>0$, the reduced Palm distribution of $X_u$ is a LGCP with intensity function $\rho(v|u)=\rho(v)\exp(c(u,v))$ but the pcf is still $g(v,w|u)=g(v,w)=\exp(c(v,w))$ for $v,w\in S$. This follows from \eqref{e:densityCox} and \eqref{e:PalmdensityCox} in the cases $S=W=\mathbb S^k$ and $S=W=L$; if $S=\mathbb R^k$, see \citet{CoeurjollyEtAl}. Consequently, if $c$ is isotropic, the $K$-functions of $X$ and $X_u$ agree. We return to this result in Section~\ref{s:7.4}. 

%Therefore,
%the $K$-functions based on $X$ and $X_u$ agree, so
%if $X$ is observed within a bounded window $W\subseteq S$, empirical estimators $\hat K$ and $\hat K_u$ for $u\in X_W$ are expected to be close. This may be exploited when fitting a LGCP model, but to the best of our knowledge this possibility has yet not been investigated...... {\color{red} JAKOB, dette er noget vi boer undersoege i afsnittet om anvendelser med maaske ogsaa her ifm et simulations studie...!!}
%
%
%TBD: The 
%special structure in \eqref{e:n-LGCP} can be exploited to construct a third-order moment characteristic  
% which
%is useful for model checking, see \citet[][Theorem 2]{LGCP}. This can be extended to the more general setting considered in the present paper...  {\color{red} JAKOB, maaske dropper vi det og bare naevner det til sidst i diskussions afsnittet... eller hvad synes du? }

Assuming $S=L$ is a linear network, let us return to the concept of IRMPS as defined by \eqref{e:IRMPS}. \citet{CronieEtAl} noticed that IRMPS is satisfied  for the LGCP if $\inf\rho>0$ and for all $u_1,u_2,u\in L$, %it is satisfied that 
\begin{equation}\label{e:strong-cond}
c(u_1,u_2)=c_1(\delta(u,u_1),\delta(u,u_2))
\end{equation} 
for some function $c_1$ \citep[in our notation; see][Equation (29)]{CronieEtAl}. This statement is true due to \eqref{e:n-LGCP}, however, in our opinion \eqref{e:strong-cond} is a very strong condition, since we are not aware of any good examples where it is satisfied unless $L$ is isometric to a closed interval and $\delta$ is usual (Euclidean/geodesic/resistance) distance. 

Incidentally, in \citet[][Lemma 2]{CronieEtAl} the metric $\delta$ is assumed to be origin independent; they did not define the meaning of `origin independent' but we have been informed (by personal communication) that they had in mind that \eqref{e:strong-cond} should be  satisfied and %unfortunately 
when they let $\delta=d_{\mathcal R}$ be the resistance metric \citep[in the text after Lemma 2 in][]{CronieEtAl},
they admit that they misunderstood the meaning of origin independent as used in \citet[][Proposition 2]{AnderesEtAl}. Moreover, 
%even when using their meaning of origin independent, 
the proof of \citet[][Lemma 2]{CronieEtAl} is incorrect (they claim that $\delta(u',u_1)=\delta(u'',u_1)$ for any $u',u'',u_1\in L$, which is obviously not correct if $\delta=d_{\mathcal R}$, and which seems wrong in general for another choice of metric). 

%Moreover, along similar lines as in \citet[][Theorem~1]{CoeurjollyEtAl}  % Coeurjolly et al. (2017, Theorem 1), 
%it can be shown that for any $u\in S$, $X\setminus\{u\}|u\in X$ (formally the reduced Palm distribution of $X$ at $u$) is a LGCP with intensity function $\rho(v|u)=\rho(v)\exp(c(u,v))$ but the pair correlation function is still $g(v,w|u)=g(v,w)=\exp(c(v,w))$ for $v,w\in S$.
%Because of these results, statistical inference procedures for estimation in parametric models are mostly based on the moment properties of LGCPs although there are important exceptions of likelihood based inference as discussed below. 


%\subsection{The case $S=\mathbb R^k$}

%TBD: SKAL FLG.\ MED? Suppose $S=\mathbb R^k$. Then 
% $X$ is SOIRS if and only if $c(u,v)=C(u-v)$ is stationary, and $X$ is stationary if and only if $\rho=\rho_0$ is constant and $c$ is stationary. If $X$ is stationary, a
%  straightforward calculation gives 
%\begin{equation}\label{e:F}
%F(t)=\sum_{n=1}^\infty \frac{(-1)^{n-1}\rho_0^n}{n!}\int_{\|u_1\|\le t}\cdots\int_{\|u_n\|\le t}\exp\bigg(\sum_{1\le i<j\le n}C(u_i-u_j)\bigg)\,\mathrm du_1\,\cdots\,\mathrm du_n 
%\end{equation}
%and
%\begin{equation}\label{e:G}
%G(t)=\sum_{n=1}^\infty \frac{(-1)^{n-1}\rho_0^n}{n!}\int_{\|u_1\|\le t}\cdots\int_{\|u_n\|\le t}\exp\bigg(\sum_{i=1}^nC(u_i)+\sum_{1\le i<j\le n}C(u_i-u_j)\bigg)\,\mathrm du_1\,\cdots\,\mathrm du_n 
%\end{equation}
%provided conditions apply so that certain expectations, sums, and integrals can be interchanged by Fubini's theorem.
%Thus, in general $F$, $G$, and $J$ are not expressible on closed form. Palm-LGCP paper: $J\le1$. !!! SAA DET SKAL NOK MED... `POINTEN' ER AT VI MAA TY TIL Van Lieshout's inhomogene $J$-fkt, men det har Cronie et al.\ ikke haft held med ifm en LGCP naar $S=L$.... 

%LGCPs have found numerous statistical applications and form perhaps the most popular subclass of Cox processes when $S=\mathbb R^k$: 
% In most statistical applications $c$ or equivalently $g$ is assumed to be isotropic, and as the likelihood function is intractable, 
%minimum contrast or composite likelihood methods based on $\rho$ and $g$ are usually used for estimation, see \citet[][]{MW07,MW17} and the references therein. See also \citet{MT} for the case where $c$ is only assumed to be geometric anisotropic, cf.\ Section~\ref{s:2.1}. Likelihood-based MCMC methods include 
%missing data MCMC methods for maximum likelihood \citep{Geyer,MW04} and Bayesian approaches \citep{RueEtAl,GT}. 

%\subsection{The case $S=\mathbb S^k$}
%
%Suppose $S=\mathbb S^k$. Then $X$ is isotropic if and only if $\rho=\rho_0$ is constant and $c(u,v)=c_0(d_{\mathcal G}(u,v))$ is isotropic, in which case \eqref{e:F} and \eqref{e:G} hold after making obvious modifications. \citet[][Propositions 1 and 4]{CM} study the smoothness properties of $\Lambda_0$ under conditions satisfied by a number of isotropic covariance functions including those in Table~\ref{tab:covariances-sphere} and the multiquadric model 
%given by
%\begin{equation}\label{e:multiquadric}
%c_{(\sigma,\delta,\tau)}(t) = \sigma^2\left(\frac{(1-\delta)^2}{1 + \delta^2-2\delta\cos t}\right)^\tau,\quad t \in[0,\pi],
%\end{equation}
%where $(\sigma,\delta,\tau)\in(0,\infty)\times(0,1)\times(0,\infty)$.
%
%I am only aware about a few papers with statistical applications of LGCPs when $S=\mathbb S^k$: \citet{SimpsonEtAl} % Simpson et al. (2016)
%consider a 
%Bayesian analysis of inhomogeneous LGCPs where 
%$c$ is an extension of the Mat{\'e}rn model in Tabel~\ref{tab:covariances-sphere} given by an infinite series based on spherical
%harmonics \citep[Simpson][Section 6.4]{}, and %(Simpson, 2009, Section 6.4). 
%using the R-INLA approach \citep{VGR} % REF MED?
%as summarized in 
%\citet[][Section~4.2]{CM}. Minimum contrast estimation based on $\rho$ and $g$ is applied in \citet{CM} and they demonstrate how a thinning procedure of an inhomogeneous LGCP gives a homogeneous LGCP so that empirical estimates of $F$, $G$, and $J$ apply for model checking. 
%Specifically, for a large dataset of sky positions of galaxies \citep[which has previously been modelled by a another kind of Cox process in][]{LawrenceEtAl}, \citet{CM} fit an inhomogeneous LGCP with a multiquadric covariance function.
% 
%%establish a simple condition ensuring almost surely that $Z$ is H{\"o}lder continuous (of some order) and hence that $\Lambda_0$ is continuous, where the condition is satisfied for a number of isotropic covariance functions including those in Table~\ref{tab:covariances-sphere}, cf.\ \citet[][Propositions 1 and 4]{CM}.
%%
%%where the focus is on isotropic covariance functions % CM = Francisco and my paper...
%%which are expressible on closed form (like in Table~\ref{tab:covariances-sphere}) so that it is easy to work with 
%%$g = \exp(c)$. \citet{CM} establish a simple condition ensuring almost surely that $Z$ is H{\"o}lder continuous (of some order) and hence that $\Lambda_0$ is continuous, where the condition is satisfied for a number of isotropic covariance functions including those in Table~\ref{tab:covariances-sphere}, cf.\ \citet[][Propositions 1 and 4]{CM}. Moreover, they fit an inhomogeneous LGCP with a so-called multiquadric covariance function for a dataset of sky positions of galaxies \citep[which has previously been modelled by a another kind of Cox process in][]{LawrenceEtAl} and show how a thinning procedure gives a homogeneous LGCP so that functional summary statistics easily applies for model checking.  

%\subsection{The case $S=L$}
%
%For $S=L$, the potential of using LGCPs with isotropic covariance functions \citep[as developed in][]{AnderesEtAl} has seemingly yet not been much explored in the literature, at least not in connection to real data analysis.  ADD SOMETHING ABOUT \citep{CronieEtAl}?? 
%
%%: Their covariance model for the underlying GRF
%%is an extension of the Mat ́ern covariance function defined as the covariance function at any fixed
%%time to the solution to a stochastic partial differential equation which is stationary in time and
%%isotropic on the sphere. Using the R-INLA approach, a finite approximation of their LGCP based
%%on a triangulation of the sphere is used, which for computational efficiency requires α := ν + d/2
%%to be an integer, where ν > 0 is the shape parameter (the expression α = ν − d/2 in their paper
%%is a typo). Simpson et al. (2016) considered an inhomogeneous LGCP defined only on the worlds
%%oceans, with d = 2 and α = 2, so ν = 1. In the original Mat ́ern model (Gneiting, 2013, Table 1),
%%0 < ν ≤ 0.5 and the covariance function has a nice expression in term of a Bessel function, but
%%otherwise the covariance function can only be expressed by an infinite series in terms of spherical
%%harmonics (Simpson, 2009, Section 6.4). In contrast, we have focused on covariance functions c
%%which are expressible on closed form so that we can easily work with the pair correlation function
%%g = exp(c). Indeed it could be interesting to apply the R-INLA approach as well as other likelihood
%%based methods of inference (Møller and Waagepetersen, 2017), though they are certainly much
%%more complicated than the simple inference procedure considered in the present paper. 

%Finally, when $S=L$...

\subsubsection{Interrupted Cox processes}

Let $X$ be an ICP, cf.\ \eqref{e:defInter}.
%Suppose $X$ is a Cox process with residual random field $\Lambda_0(u)=\Pi(u)/\mathbb E\Pi(u)$ where 
%$\Pi=\{\Pi(u)\,|\,u\in S\}\subset[0,1]$ is a random field so that $\mathbb E\Pi(u)>0$ for all $u\in S$. 
Then $X$ conditioned on $\Pi$ is obtained by an independent thinning of a Poisson process $Z$ on $S$ with intensity function $\rho_Z(u)=\rho(u)\left(1+c(u,u)\right)^{h/2}$, where the selection probabilities are given by $\Pi$. In the terminology of \citet{Stoyan}, $X$ is an interrupted point process. 

%Following an example in \citet{}, a
\citet{FredMe} studied the ICP (as well as many other examples of interrupted point processes) when $S=\mathbb R^k$. For our general state space setting, assuming for ease of presentation that $c(u,v)=\sigma^2r_0(d(u,v))$ is isotropic with $r_0(0)=1$, we obtain in a similar way as in \citet{FredMe} that the mean selection probability is constant and given by
\begin{equation}\label{e:pms}
p_{\mathrm{ms}}=(1+2\sigma^2)^{-h/2}
\end{equation}
and the pcf is given by
\begin{equation}\label{e:g0ICP}
g_0(t)=\left(\frac{(1+\sigma^2)^2}{(1+\sigma^2)^2-\sigma^4 r_0(t)^2}\right)^{h/2}.
\end{equation}
Third and higher-order moment results are less simple to express, but it can be proven that $X$ is positively correlated at all orders.
As $\sigma$ increases from 0 to infinity, then $p_{\mathrm{ms}}$ decreases from 1 to 0, whilst $g_0(t)$ increases from 1 to $(1-r_0(t)^2)-{h/2}$
if $r_0(t)\not=0$, which shows a trade-off between the degree of thinning and the degree of clustering. 
To understand how $g_0(t)$ depends on $h$ it
 is natural to fix the value of $p_{\mathrm{ms}}\in(0,1)$. 
 % $$\mathbb E\Pi(u)=\rho(u)/\rho_Y(u)=(1+\sigma^2)^{-h/2},$$
   Then
  \[g_0(t)=\left(1+\left(1-p_{\mathrm{ms}}^{h/2}\right)^2r_0(t)^2\right)^{-h/2}\]
  is a strictly decreasing function of $h$ whenever $r_0(t)\not=0$.
%   , as $h$ increases, $\sigma^2=q^{-h/2}-1$ decreases toward 0 and 
%%\[[g(u,v)=\left(\frac{\left(1+1/c(u,u)\right)\left(1+1/c(v,v)\right)}{\left(1+1/c(u,u)\right)\left(1+1/c(v,v)\right)-\sigma^4 r(u,v)^2}\right)^{h/2}\] 
%$g_0(t)$ decreases toward 1. 
Consequently, taking $h=1$ is natural if we wish to model as much clustering as possible when keeping the degree of thinning fixed. 

%In addition, assume that 
%\begin{eqnarray}\label{e:Pi}
%\Pi(u)=\exp\bigg(-\sum_{i=1}^h Z_i(u)^2\bigg)
%\end{eqnarray}
%where $Z_i=\{Z_i(u)\,|\,u\in S\}$, $i=1,\ldots,h$ are independent identically distributed zero mean GRFs with covariance function $c$. 
%It follows immediately from \eqref{e:g-Lambda0} and \eqref{e:Pi} that 
%\begin{equation}\label{e:selprob}
%\rho_Y(u)=\rho(u)\left(1+c(u,u)\right)^{h/2}
%\end{equation}
% and
%\[g(u,v)=\left(\frac{\left(1+c(u,u)\right)\left(1+c(v,v)\right)}{\left(1+c(u,u)\right)\left(1+c(v,v)\right)-c(u,v)^2}\right)^{h/2}.\]
%Thus $g(u,v)\ge1$, with equality if and only if $c(u,v)=0$. This reflects that the thinning causes clustering/aggregation in $X$. To better understand this behaviour, suppose for simplicity that $c(u,v)=\sigma^2r_0(d(u,v))$
%is isotropic and $\sigma^2>0$. Then the ratio of intensities of $X$ and $Y$ is given by $q=(1+\sigma^2)^{-h/2}$, cf.\ \eqref{e:selprob}, and $g(u,v)=g_0(d(u,v))$ is isotropic with
%\[g_0(t)=\left(\frac{(1+\sigma^2)^2}{(1+\sigma^2)^2-\sigma^4 r_0(t)^2}\right)^{h/2}.\]
%As $\sigma$ increases from 0 to infinity, then $q$ decreases from 1 to 0, whilst $g_0(t)$ increases from 1 to $(1-r_0(t)^2)-{h/2}$
%if $r_0(t)\not=0$, which shows a trade-off between the degree of thinning and the degree of clustering. 
%To understand how $g_0(t)$ depends on $h$ it
% is natural to fix the value of $q\in(0,1)$. 
% % $$\mathbb E\Pi(u)=\rho(u)/\rho_Y(u)=(1+\sigma^2)^{-h/2},$$
%   Then
%  \[g_0(t)=\left(1+\left(1-q^{h/2}\right)^2r_0(t)^2\right)^{-h/2}\]
%  is a strictly increasing function of $h$ whenever $r_0(t)\not=0$.
%%   , as $h$ increases, $\sigma^2=q^{-h/2}-1$ decreases toward 0 and 
%%%\[[g(u,v)=\left(\frac{\left(1+1/c(u,u)\right)\left(1+1/c(v,v)\right)}{\left(1+1/c(u,u)\right)\left(1+1/c(v,v)\right)-\sigma^4 r(u,v)^2}\right)^{h/2}\] 
%%$g_0(t)$ decreases toward 1. 
%Consequently, taking $h=1$ is natural if one wishes to model as much clustering as possible. 
%
%Third order moments...???
%
%%If $c(u,v)=\sigma^2r_0(d(u,v))$
%%is isotropic and $\sigma^2>0$, then $g(u,v)=g_0(d(u,v))$ is isotropic with
%%\[g_0(t)=\left(\frac{(1+1/\sigma^2)^2}{(1+1/\sigma^2)^2-r_0(t)^2}\right)^{h/2}.\]
%%For fixed $t$ with $r_0(t)\not=0$, $g_0(t)$ is a strictly increasing function of both $\sigma^2$, $|r_0(t)|$, and $h$.  


For the case $S=L$ and $r_0(t)=\exp(-t/\phi)$ specifying an isotropic exponential covariance function, expressions of $K_L$ when $h=1,2,\ldots,5$ were given in \citep[][Appendix A]{HeidiMe}. Although these expressions were given for $d=d_{\mathcal G}$, they remain true for a general metric $d$ because the pcf depends only on $h$, $\sigma^2$, and $r_0(t)=\exp(-t/\phi)$, cf.\ \eqref{e:g0ICP} and our comment after \eqref{e:K2new}. 

 

%The definition of an interrupted point process in Stoyan (1979...) extends immediately as follows. 
%
%Let $\Pi=\{\Pi(u)\,|\,u\in S\}\subset[0,1]$ be a random field and $Y$ a homogeneous Poisson point process on $S$ with intensity $\rho_Y>0$, where $\Pi$ and $Y$ are independent. Assume $X$ is an independent thinning of $Y$ with retention probabilities, i.e., $X=\{y\in Y\,|\,R(y)\le\Pi(y)\}$ where $R=\{R(u)\,|\,u\in S\}$ is independent of $(Y,\Pi)$ and consists of independent random variables $R(u)$ which are uniformly distributed on $[0,1]$. In the terminology of \citet{}, $X$ is an interrupted point process.  
%
%Fred and me... extends previous work by Dietrich Stoyan on interrupted point processes. Dependent thinning of a DPP...(Later)... or here a Poisson process...!
%
%Dendritic spines, which are small protrusions on the dendrites of a neuron,
%are of interest in neuroscience as they are related to cognitive processes such
%as learning and memory. We analyse the distribution of spine locations on
%six different dendrite trees from mouse neurons using point process theory for
%linear networks. Besides some possible small-scale repulsion, we find that two
%of the spine point pattern datasets may be described by inhomogeneous Poisson 
%process models, while the other point pattern datasets exhibit clustering
%between spines at a larger scale. To model this we propose an inhomogeneous
%Cox process model constructed by thinning a Poisson process on a linear network with retention probabilities determined by a spatially correlated random
%field. For model checking we consider network analogues of the empirical F-,
%G-, and J-functions originally introduced for inhomogeneous point processes
%on a Euclidean space. The fitted Cox process models seem to catch the clus-
%tering of spine locations between spines, but also posses a large variance in the
%number of points for some of the datasets causing large confidence regions for
%the empirical F- and G-functions.
%
%Heidi and me: K-fct (Appendix)...

\subsubsection{Permanental Cox point processes}

Let $X$ be a PCPP, % with parameters $\alpha=h/2$ and $c$, 
cf.\ \eqref{e:defPCPP}. The case $S=\mathbb R^k$ was studied in \citet{Macchi:75} and \citet{MacCM}; using the parametrization in the latter paper, $X$ is a PCPP with parameters $\alpha=h/2$ and $C(u,v)=\sqrt{\rho(u)\rho(v)}c(u,v)/\alpha$. 
The process can also be defined in the general state space case, see \citet{ShiraiTakahashi}.

To stress that $c$ is a correlation function, we shall write $c=r$. For $n=1,2,\ldots$ and $u_1,\ldots,u_n\in S$, define the $\alpha$-weighted
permanent by
\[{\mathrm{per}}_\alpha[r](u_1,\ldots,u_n)=\sum_{\pi}\alpha^{\#\pi}r(u_1,u_{\pi_1})\cdots r(u_n,u_{\pi_n})
\]
where the sum is over all permutations $\pi=(\pi_1,\ldots,\pi_n)$ of $(1,\ldots,n)$ and $\#\pi$ is the number of cycles. 
%This is a polynomial of degree $n$ in $α$ in which the coefficient of degree $r$ is the sum of products over permutations having exactly $r$ cycles. 
The usual permanent corresponds to $\alpha = 1$ \citep{Minc}, in which case $X$ is also called a Boson process \citep{Macchi:75}. We have
\begin{equation*}\label{e:uha}
\rho(u_1,\ldots,u_n)=\rho(u_1)\cdots\rho(u_n)\mathrm{per}_\alpha[r](u_1,\ldots,u_n)/\alpha^n
%=\mathrm{per}_\alpha[c](u_1,\ldots,u_n),
\end{equation*}
from which it can be verified that 
 $X$ is positively correlated at all orders. It also follows that the degree of clustering is a decreasing function of $\alpha$. In particular,
\begin{equation}\label{e:pairPCPP}
g(u,v)=1+r(u,v)^2/\alpha.
\end{equation}
Thus $g\le 1+1/\alpha\le 3$, which reflects the limitation of modelling clustering by a PCPP.  

\citet{Valiant} showed that exact computation of permanents of general
matrices is a $\#$P (sharp P) complete problem, so no deterministic polynomial time algorithm is available. 
%However, polynomial-time algorithms exist for certain special cases, such as general fixed-rank matrices (Barvinok 1996), and for approximate Monte-Carlo computation of general non-negative matrices (Jerrum et al. 2004). 
For most statistical purposes, approximate computation of
permanent ratios is sufficient, and analytic approximations are available for large $\alpha$. However, as $\alpha\rightarrow\infty$, $X$ tends to a Poisson process and the process becomes less and less interesting for the purpose of modelling clustering. 

The PCPP can be extended to the case where $\alpha\ge0$ and $c$ is not necessarily a covariance function (in which case we loose the connection to Gaussian and Cox processes), see \citet{MacCM} and \citet{ShiraiTakahashi}. Indeed the process also extends to the case where $\alpha$ is a negative integer whereby a (weighted) determinantal point process (DPP) is obtained. We return to DPPs in Section~\ref{s:7.1}. 

\subsection{An index of cluster strength}
Consider again a LGCP, ICP, or PCPP $X$ defined on the metric space $S$ and with an isotropic covariance function $c(u,v)=\sigma^2r_0(d(u,v))$ where $r_0(0)=1$, $h=1$ if $X$ is an ICP, and $\sigma=1$ if $X$ is a PCPP. 
To quantify how far $X$ is 
from a Poisson process we follow 
\citet{BadEtAl2022} in defining 
\begin{equation}\label{e:varphi}
\varphi=g_0(0)-1.
\end{equation}
If $X$ has constant intensity $\rho$,
the total variation distance between the distribution of $X$ and that of a Poisson process with intensity $\rho$ is at most $\rho\nu(S)\sqrt\varphi$. This follows since 
\citet[][Lemma~9]{BadEtAl2022} immediately extends to our setting.

\citet{BadEtAl2022} used $\varphi$ to describe the degree of clustering for another class of Cox processes, namely Neyman-Scott point processes (with the cluster size following a Poisson distribution): the degree of clustering increases as $\varphi$ increases. 
Although our Cox processes are not cluster point processes, $\varphi$ has a similar interpretation since realizations of the processes may look more or less clustered. More precisely, 
it follows from \eqref{e:pairLGCP}, \eqref{e:g0ICP}, and \eqref{e:pairPCPP} that $\varphi$ is equal to
\begin{equation}\label{e:varphiICP}
\exp(\sigma^2)-1\mbox{ if $X$ is a LGCP},\quad 
\frac{1+\sigma^2}{\sqrt{1+2\sigma^2}}-1\mbox{ if $X$ is an ICP},\quad 
2/h \mbox{ if $X$ is a PCPP,}
\end{equation}
and $\varphi+1$ is the maximal value of the pcf.
%\begin{align}
%\varphi&=\exp(\sigma^2)-1&\mbox{if $X$ is a LGCP},\nonumber\\
%\varphi&=\left(\frac{1+\sigma^2}{\sqrt{1+2\sigma^2}}\right)-1 &\mbox{if $X$ is a ICP},     \label{e:varphiICP}\\
%\varphi&=2/h&\mbox{if $X$ is a PCPP}.\nonumber
%\end{align}
Thus, if $X$ is a LGCP or an ICP, $\varphi$ is a strictly increasing function of $\sigma>0$, where $X$ approaches a Poisson process as $\sigma\rightarrow0$, while realizations of $X$ become more and more clustered as $\sigma$ grows. If instead $X$ is a PCPP, the degree of clustering decreases as $h\in\{1,2,...\}$ increases, where $X$ approaches a Poisson process as $h\rightarrow\infty$. 

The index $\varphi$ for `cluster strength' is used in Sections~\ref{s:simstudy} and \ref{s:7.3}. 


\section{Application of statistical inference procedures}\label{s:appl}

For the analyses of the real and simulated datasets in this section, we have used the functions \texttt{mincontrast} and \texttt{linearpcf} from \texttt{spatstat} for calculating the minimum contrast estimates and the non-parametric estimate of the pcf (the latter is only available in the case of the geodesic metric), respectively, as well as a range of other minor functions from \texttt{spatstat} for handling point processes on linear networks. Furthermore, we have used the function \texttt{global\_envelope\_test} from the \texttt{GET} package for global envelope tests. The rest of the code (such as the estimation of the pcf in the case of the resistance metric, or the estimation of the $F$, $G$, and $J$-functions) we have implemented ourselves, and it is available in our package \texttt{coxln}. 

\subsection{Analysis of Chicago crime dataset}

\cite{AngEtAl} used the empirical $K$-function together with simulations to show that the Chicago crime dataset (shown in the left panel of Figure~\ref{fig:data} and available in \texttt{spatstat}) was more clustered than a homogeneous Poisson process while an inhomogeneous Poisson process with log-quadratic intensity fitted by maximum likelihood provided a better fit. They remarked that the latter model was only provisional and  shown for demonstration. 

For the analysis of the street crimes in this paper we fitted instead all three model classes (LGCP, ICP, and PCPP) given in Section~\ref{s:cox}. %, i.e., LGCP, ICP and PCPP, for modeling the data. 
For the Gaussian process underlying the Cox process models, we used for simplicity a constant mean and an isotropic exponential covariance function with metric $d=d_{\mathcal R}$ 
(since the network is not a 1-sum of trees and cycles, the geodesic metric will not give well-defined models). Thus we have (up to) four unknown parameters in each model: $\rho$, the intensity of the point process; $\sigma^2$, the variance of the Gaussian process (this is not  a parameter in the PCPP model); $s$, the scaling of the exponential covariance function; and $h$, the number of Gaussian processes used in the model (this is not  a parameter in the LGCP model). 
%Note that $h$ is not included in the LGCP model, and $\sigma^2$ is not included in the PCPP model.

To estimate the parameters of all three models, we used the unbiased estimate $\hat{\rho} ={n(x)}/{|L|}= 0.00372$ for the intensity and estimated the remaining parameters by the 
minimum contrast method based on the pair correlation function, cf.\ \eqref{e:contrast}. We estimated $g_0$ non-parametrically on the interval $r\in[0,100]$, but since the shape of the non-parametric estimate mostly resembled the shapes of the theoretical pair correlation functions for the three models on the interval $r\in[20,100]$, in \eqref{e:contrast} we %base minimum contrast estimation on this interval. That is, in \eqref{e:contrast}, we 
let $a_1=20$ and $a_2=100$ while $p$ and $q$ were given by the default values in \texttt{spatstat}. 

The estimated parameter values are shown in Table~\ref{tab:chicago-estimates}. Figure~\ref{fig:chicago-pcf} shows the non-parametric estimate together with the estimated pair correlation functions for each of the three models, where % The figure indicates that 
the estimated pair correlation functions are very similar for the LGCP and the ICP, while the pair correlation function for the PCPP deviates from the other two at low distances. %Moreover, for the intensity of all three models, we use the unbiased estimate $\hat{\rho} ={n(x)}/{|L|}= 0.00372$. 

%Since the intensity parameter is not included in the pair correlation functions, this parameter is estimated by $\hat\rho = {n(x)}/{|L|}$ resulting in the same estimate for all three models given by $\hat{\rho} = 0.00372$. 

\begin{table}
	\centering
	\begin{tabular}{l||c|c|c|c}
		Model & $\hat\rho$ & $\hat{\sigma}^2$ & $\hat s$ & $\hat h$ \\
		\hline\hline
		LGCP  & 0.00372 & 1.70 & 0.0213  & -- \\
		ICP   & 0.00372 & 22.8 & 0.00747 & 2 \\
		PCPP   & 0.00372 &   --  & 0.00988 & 1
	\end{tabular}
	\caption{Parameter estimates for the LGCP, ICP, and PCPP models for the Chicago crime dataset.  A dash  indicates that the parameter is not present in the model.} 
	\label{tab:chicago-estimates}
\end{table}

\begin{figure}
	\begin{center}\includegraphics[width=7cm]{"Figs/chicago-pcf.pdf"}\end{center}
	\caption{Pair correlation functions for the Chicago crime dataset: Non-parametric estimate (solid curve) and curves estimated by minimum contrast for a LGCP (dashed curve), an ICP (dotted curve), and a PCPP (dashed-dotted curve) model.}
	\label{fig:chicago-pcf}
\end{figure}

The left column of Figure~\ref{fig:chicago-sim} shows a simulation from each of the three fitted models. None of the simulations show strong deviation from the data, although the ICP does show a tendency to have small densely packed clusters of points that are not present in the data. Moreover, the right column of Figure~\ref{fig:chicago-sim} shows the 95\% global envelopes based on a concatenation of the empirical $F,G,J$-functions (based on 999 simulations) discussed at the end of Section~\ref{s:est-check}. The LGCP provides the best fit with a $p$-value of 0.155 for the global envelope test. The ICP provides a rather bad fit to the data with a $p$-value of just 0.004, and the envelope shows a clear discrepancy between data and the model at distances below 50, where the discrepancy indicates more clustering in the fitted ICP model than in the data.
%reveals that in particular the $G$ and $J$-functions show a clear discrepancy between data and the model at distances below 50. 
The PCPP model shows a decent fit with a $p$-value just below 0.05 due to the $F$-function going outside the envelope at distance 100. 
%In conclusion, the LGCP model provides the best fit. % when using the exponential covariance function.

\begin{figure}
	\begin{center}
		\includegraphics[width=6.5cm]{"Figs/chicago-sim-lgcp.pdf"}
		\includegraphics[width=6cm]{"Figs/chicago-FGJglobal-lgcp.pdf"}
		\includegraphics[width=6.5cm]{"Figs/chicago-sim-icp.pdf"}
		\includegraphics[width=6cm]{"Figs/chicago-FGJglobal-icp.pdf"}
		\includegraphics[width=6.5cm]{"Figs/chicago-sim-PPP.pdf"}
		\includegraphics[width=6cm]{"Figs/chicago-FGJglobal-PPP.pdf"}
	\end{center}
	\caption{Left column: Simulations of point patterns under the estimated LGCP (upper row), ICP (middle row) and PCPP (lower row) models for the Chicago crime dataset. The width of the lines shows the simulation of the underlying transformed Gaussian processes (the random intensity functions), while the points show the simulated point pattern. Right column: 95\% global envelopes for the fitted LGCP (upper row), ICP (middle row), and PCPP (lower row) models for the Chicago crime dataset. The $p$-values of the global envelope test are shown above each plot.}
	\label{fig:chicago-sim}
\end{figure}
 
%{\color{red} Hvorfor er der ingen tekst vedr simulationer i venstre kolone i Figure~\ref{fig:chicago-sim}?} 

%\begin{figure}
%	%\begin{center}
%		\includegraphics[width=6cm]{"Figs/chicago-FGJglobal-lgcp.pdf"}
%		\includegraphics[width=6cm]{"Figs/chicago-FGJglobal-icp.pdf"}
%		\includegraphics[width=6cm]{"Figs/chicago-FGJglobal-PCPP.pdf"}
%	%\end{center}
%	\caption{Global envelope tests for the fit of the LGCP, ICP, and PCPP models for the Chicago crime dataset. The $p$-values of the test are shown above each plot.}
%	\label{fig:chicago-FGJ}
%\end{figure}

\subsection{Dendrite data}\label{s:den-an}

\citet{HeidiMe} analysed several point pattern datasets given by spine locations on different dendrite trees which were identified by linear networks. For each dataset they fitted an inhomogeneous ICP with $h=1$, $c$ an isotropic exponential covariance function with metric $d=d_{\mathcal G}=d_{\mathcal R}$, and an inhomogeneous intensity function given by a constant intensity $\rho_{\mathrm{mb}}>0$ on the main branch and a different constant intensity $\rho_{\mathrm{sb}}>0$ on the side branches. 

In this section we restrict attention to the dataset for dendrite number five in \citet{HeidiMe} (the right panel of Figure~\ref{fig:data}) modelled by 
an ICP with $h=1$ and considering different covariance function models. Hence we
estimated the intensity parameters by $\hat\rho_{\mathrm{mb}}= {n(x_{\mathrm{mb}})}/{|L_{\mathrm{mb}}|}=0.119$ and $\hat\rho_{\mathrm{sb}} = {n(x_{\mathrm{sb}})}/{|L_{\mathrm{sb}}|}=0.184$ where $L_{\mathrm{mb}}$ denotes the main branch, $L_{\mathrm{sb}}$ the union of side branches, and $x_{\mathrm{mb}}$ and $x_{\mathrm{sb}}$ the point patterns restricted to these two subsets of $L$. 

%Specifically, they used a two step estimation procedure: First, given a parametric model for $\rho$ where the parameter ranges independently of $(\sigma,\beta)$, the parameter is estimated by maximizing a first-order composite likelihood function (the Poisson likelihood), second a minimum contrast procedure based on $g_0$ was used for estimating $\sigma>0$ and $\phi>0$, where they plugged-in the estimated intensity function. Simulation studies in \citet{HeidiMe} indicated that 1) it can be difficult to estimate $\sigma^2$ and $h$ simultaneously, as an increase in $h$ can be balanced out by an increase in $\sigma^2$, 2) with the default bandwidth and kernel from the spatstat-package for estimating $g_0$, minimum contrast estimation based on $g_0$ generally performs better than if it is based on $K$ \citep[which is consistent with results in][]{FredMe}, and 3) minimum contrast estimation is far more reliable than (second-order) composite likelihood estimation. 

%We use the dataset for dendrite number five, which is shown in the right panel of Figure~\ref{fig:data}, and start out by making a model similar to the one used in \citet{HeidiMe}; that is, we use an ICP model with an exponential covariance function and an inhomogeneous intensity function given by a constant intensity $\rho_b>0$ on the main branch and a different constant intensity $\rho_s>0$ on the side branches.

First, we let $c(u,v)=\sigma^2\exp(-sd(u,v))$ be an isotropic exponential covariance function with $d=d_{\mathcal G}=d_{\mathcal R}$. Then we estimated $(\sigma^2,s)$ following the recommendations in \citet{HeidiMe}, i.e., we %fix $h=1$ and 
used the minimum contrast method with  $a_1=0$, $a_2=50$, $p=2$, and $q=1$ in \eqref{e:contrast}.
%\citep[in \eqref{e:contrast} we use $a_1=0$, $a_2=50$, $p=2$, and $q=1$ which match the choice in ][]{HeidiMe}. 
The obtained estimate $(\hat\sigma^2,\hat s)$ is given in Table~\ref{tab:neu-estimates} and the corresponding estimated pcf is shown in Figure~\ref{fig:neu-pcf} (the dotted curve). The estimate 
deviates a bit from that obtained in \citet{HeidiMe}, which is due to a different implementation of the non-parametric estimation of the pair correlation function.
Figure~\ref{fig:neu-FGJ} (left panel) shows the 95\% global envelope based on a concatenation of the empirical $F,G,J$-functions (based on 999 simulations). This shows a satisfactory fit with a $p$-value of $0.26$ for the global envelope test. The $p$-value deviates substantially from the value obtained in \citet{HeidiMe}, which may be due to the different estimates of $(\sigma^2,s)$ and a different number of simulations used in the global envelope test.

\begin{table}
	\centering
	\begin{tabular}{l||c|c}
		Covariance function &  $\hat{\sigma}^2$ & Other estimates\\
		\hline\hline
		Exponential   & $3.90$ & $\hat{s}=0.0356$ \\
		Gamma Bernstein CDF  & $3.91$ &  $(\hat\tau,\hat\phi)=(22.6, 626)$\\ 
		Inverse gamma Bernstein CDF & $3.90$ &  $(\hat\tau,\hat\phi)=(163, 5.79)$ \\
		Generalised inverse Gaussian Bernstein CDF & $3.90$ & $(\hat\psi,\hat\chi,\hat\lambda) = (213, 35.8, -499)$ \\
		Inverse gamma Bernstein CDF with fixed $\tau$  & $4.63$ & $(\hat\tau,\hat\phi)=(2, 0.0188)$ \\
	\end{tabular}
	\caption{Parameter estimates for the ICP models with various covariance functions for the dendrite dataset.} 
	\label{tab:neu-estimates}
\end{table}
%\begin{table}
%	\centering
%	\begin{tabular}{l||c|c|c|c|c}
%		Covariance function & $\hat\rho_{\mathrm{mb}}$ & $\hat\rho_b$ & $\hat h$ & $\hat{\sigma}^2$ & Other estimates\\
%		\hline\hline
%		Exponential  & $0.119$ & $0.184$ & $1$ & $3.90$ & $\hat{s}=0.0356$ \\
%		Gamma & $0.119$ & $0.184$ & $1$ & $3.91$ &  $(\hat\tau,\hat\phi)=(22.6, 626)$\\ 
%		Inverse gamma & $0.119$ & $0.184$ & $1$ & $3.90$ &  $(\hat\tau,\hat\phi)=(163, 5.79)$ \\
%		Generalised inverse Gaussian & $0.119$ & $0.184$ & $1$ & $3.90$ & $(\hat\psi,\hat\chi,\hat\lambda) = (213, 35.8, -499)$ \\
%		Inverse gamma with fixed $\tau$ & $0.119$ & $0.184$ & $1$ & $4.63$ & $(\hat\tau,\hat\phi)=(2, 0.0188)$ \\
%	\end{tabular}
%	\caption{Parameter estimates for the ICP models with various covariance functions for the dendrite dataset.} 
%	\label{tab:neu-estimates}
%\end{table}

\begin{figure}
	\begin{center}
		\includegraphics[width=6cm]{"Figs/neu-pcf.pdf"}\end{center}
	\caption{Pair correlation functions for the dendrite dataset: Non-parametric estimate (solid curve), and curves estimated by minimum contrast for ICP with exponential covariance  (dotted curve) and covariance function with inverse gamma Bernstein density with fixed $\tau=2$  (dotted-dashed curve).}
	\label{fig:neu-pcf}
\end{figure}

\begin{figure}
	\begin{center}
	\includegraphics[width=6cm]{"Figs/neu-FGJglobal-exp.pdf"}
	\includegraphics[width=6cm]{"Figs/neu-FGJglobal-invgammatau2.pdf"}
	\end{center}
	\caption{95\% global envelopes for the fitted ICP models with an exponential covariance function (left) and a covariance function with an inverse gamma Bernstein CDF with $\tau=2$ (right) for the dendrite dataset. The $p$-values of the global envelope test are shown above each plot.}
	\label{fig:neu-FGJ}
\end{figure}

Second, we let $c$ be one of the three covariance functions given in Example~\ref{ex:0}, i.e., those with a gamma, an inverse gamma, and a generalised inverse Gaussian Bernstein CDF. We used the same minimum contrast procedure as above for estimating the parameters, thereby obtaining the estimates in Table~\ref{tab:neu-estimates}. 
%The parameters $\rho_{\mathrm{mb}}$, $\rho_b$, and $h$ are equal to the case using the exponential covariance function as $\rho_{\mathrm{mb}}$ and $\rho_b$ are obtained from the same calculations and $h$ is fixed at the same value. 
For all three cases of the estimated Bernstein CDFs, the variance is close to zero meaning that the distributions are almost degenerate and the corresponding covariance functions are very close to exponential covariance functions. Since the mean values of these distributions are very close to $\hat s$ from the estimated exponential covariance function, all three estimated covariance functions are very close to the estimated exponential covariance function, thus leading to the same estimated model. 
This suggests possible problems with the estimation procedure, which we will explore in a simulation study in Section~\ref{s:simstudy}, or with unidentifiability of the parameters of the Bernstein CDF.
 
To try out a model which do not become almost identical to the model using the exponential covariance function, we considered a covariance function with an inverse gamma Bernstein CDF with $\tau=2$ fixed (corresponding to an inverse gamma distribution with an infinite variance). 
Table~\ref{tab:neu-estimates} shows the minimum contrast estimate and Figure~\ref{fig:neu-pcf} shows 
 the corresponding estimated covariance function. This covariance function has a larger variance parameter $\sigma^2$ and a heavier tail than the estimated exponential covariance function. The 95\% global envelope in Figure~\ref{fig:neu-FGJ} (right panel) and the $p$-value of $0.243$ for the global envelope test reveal that the covariance function with inverse gamma density and fixed parameter $\tau=2$ provides a similar good fit as the exponential covariance function. This observation is further discussed in Section~\ref{s:7.3}.

\subsection{Simulation study}\label{s:simstudy}

The analysis of the dendrite dataset showed the problem that we in practice get an exponential covariance function when we fit the covariance functions given in Example~\ref{ex:0}. It should be noted that all these covariance functions have the exponential covariance function as a limiting case, which makes it possible to get arbitrarily close to an exponential covariance function in the estimation procedure, while this was not the case when the $\tau$ parameter was fixed in the analysis of the dendrite data. 

%To explore whether this is a general problem or simply applies to the dendrite data, we will make a simulation using a covariance function which is clearly different from the exponential case to see if we still get an exponential covariance function from the estimation procedure. Specifically we take the network used in the dendrite data and simulate an ICP with homogeneous intensity and covariance function with inverse gamma Bernstein CDF, where the parameters used are shown in Table~\ref{tab:sim-estimates}. The left panel of Figure~\ref{fig:sim-pcf} shows the simulation along with the underlying GP. Note that the chosen parameters result in an infinite variance for the inverse gamma distribution, while the exponential covariance function corresponds to a variance of zero. In the estimation we fix $h=1$ at its true value, and focus on the estimation of the parameters $\sigma^2$, $\tau$, and $\phi$ using the minimum contrast method based on the pair correlation function. Table~\ref{tab:sim-estimates} shows the estimated parameters and the right panel of Figure~\ref{fig:sim-pcf} shows the corresponding  estimated pcf together with the
%true pcf and a non-parametric estimate of the pcf. We see that the estimated pcf fits well with the non-parametric estimate, but both of these deviate from the true pcf. Again the estimated parameters correspond to an inverse gamma distribution with a variance close to 0, implying an estimated covariance function which is close to an exponential covariance function.

%\begin{table}
%	\centering
%	\begin{tabular}{l||c|c|c|c|c}
%		Parameters & $\rho$ & $h$ & $\sigma^2$ & $\tau$ & $\phi$ \\
%		\hline\hline
%		True parameters & $1$ & $1$ & $5$ & $2$ & $0.02$ \\
%		Estimates & $1.470$ & $1$ & $3.59$ & $141.4$ &  $6.56$\\
%		Average estimates & $1.017$ & $1$ & $16.32$ & $117.4$ & $2.67$
%	\end{tabular}
%	\caption{True parameter values used in the simulations, estimates obtained by minimum contrast estimation from a single simulation, and average estimates obtained from 1000 simulations.} 
%	\label{tab:sim-estimates}
%\end{table}
%
%\begin{figure}
%	\begin{center}
%		\includegraphics[width=7cm]{"Figs/simstudy-1sim-sim.pdf"}
%		\includegraphics[width=7cm]{"Figs/simstudy-1sim-pcf.pdf"}
%	\end{center}
%	\caption{Left: A simulation of the ICP model with inverse gamma Bernstein CDF. Right: The non-parametric estimate of the pcf for the simulation (solid curve), the pcf for the true ICP model (dotted-dashed curve), and the pcf for the estimated ICP model (dotted curve).}
%	\label{fig:sim-pcf}
%\end{figure}

%Repeating the above with new simulations using the same parameters, we see that it is quite common that the non-parametric estimate of the pcf deviates from the true pcf, suggesting that the non-parametric estimate of the pcf does not capture the various parameters in the used non-exponential covariance functions well, and so the minimum contrast estimation procedure also fails at estimating these. {\color{red} Hvad menes mon med `quite common'? Og er det virkelig estimationsmetoden, der er problemt ? Det er jo ikke ligefrem konklusion i den nylige artikel af Adrian \& Co, hvor det skyldes problemer med identificerbarhed af parametrene. Derfor kunne det vaere interessant at se, hvad der sker hvis pcf er meget laengere vaek fra 1 end det ses i Fig. 7.} We also see that it is quite common {\color{red} Hvad menes mon med `quite common'?}that estimated parameters correspond to an exponential covariance function. We confirm this observation by making 1000 simulations, estimating the parameters for each, and calculating the averages of the estimated parameters, which are shown in Table~\ref{tab:sim-estimates}. The average estimates also correspond to an inverse gamma distribution with a variance close to zero, suggesting that this is indeed the typical behavior.
%
%{\color{green} (Quite common vil jeg her vurdere til omkring 50\% af gangene, men det er meget subjektivt hvornår det ikke-parametrise estimat og den estimerede models pcf er tæt på hinanden, men langt fra den rigtige pcf anvendt til simulation.  Hvis jeg kigger på en pcf der ligger meget længere væk fra 1, fx ved at gange $\sigma^2$ med 100 så ændrer det ikke rigtigt på noget. Jeg har følgende forslag til omskrivning af ovenstående paragraf.)}


%{\color{red}
%Repeating the above with new simulations using the same parameters, we see that it is quite common (roughly $50\%$) that the non-parametric estimate of the pcf deviates from the true pcf, suggesting that the non-parametric estimate of the pcf does not capture the various parameters in the used non-exponential covariance functions well. The reason for this may be that the parameters of the used covariance function are not identifiable from the data. Making 1000 simulations, we estimate the parameters of each simulation, and calculate the averages of the estimated parameters, which are shown in Table~\ref{tab:sim-estimates}. The average estimates also correspond to an inverse gamma distribution with a variance close to zero, suggesting that this is indeed the typical behavior.}


%To explore whether this was a general problem or simply applied to the dendrite data, we made a number of simulations using a covariance function which was clearly different from the exponential case to see if we still got exponential covariance functions from the estimation procedure. Specifically we took the network used in the dendrite data and simulated 1000 simulations of an ICP with homogeneous intensity and covariance function with inverse gamma Bernstein CDF with the parameters $(\rho, h, \sigma^2, \tau, \phi) = (1, 1, 5, 2, 0.02)$. The mean number of points  is $\mathbb E N(L) = \rho|L| = 736$. For each simulation we fitted two models using minimum constrast estimation: an ICP with an inverse gamma Bernstein CDF with $h=1$ and $\tau=2$ fixed, and an ICP with an exponential covariance function with $h=1$ fixed. For each simulation, we then calculated the distance from the non-parametric estimate to each of the two pcfs for the estimated models, say $D_\text{IG}$ and $D_\text{Exp}$, where we used $D$ in \eqref{e:contrast} for calculating distances. Figure \ref{fig:sim-ratiohist} shows a histogram of the ratio $D_\text{IG}/D_\text{Exp}$. In $71.7\%$ of the simulations, the non-parametric estimate is closer to the model using the exponential covariance function than the model using the inverse gamma Bernstein CDF with $\tau=2$, showing that there indeed is a tendency for getting an estimated model corresponding to an exponential covariance function. 


%\begin{figure}
%	\begin{center}
%		\includegraphics[width=7cm]{"Figs/simstudy-1000sim-ratiohist.pdf"}
%	\end{center}
%	\caption{A histogram of the ratio of $D_\text{IG}$ and $D_\text{Exp}$ calculated from 1000 simulations.}
%	\label{fig:sim-ratiohist}
%\end{figure}



To explore whether this was a general problem or simply applied to the dendrite dataset, we made a number of simulations using covariance functions which were different than the exponential case to see if we still obtained exponential covariance functions from the estimation procedure. Specifically we took the network used in the dendrite dataset and simulated 1000 simulations of an ICP with a homogeneous intensity function and a covariance function with inverse gamma Bernstein CDF. We did this for various combinations of parameters, where $\sigma^2\in\{10^{-1},10^0,10^1,10^2,10^3,10^4\}$ and $\tau\in\{1.1,1.5,2,5\}$ while the rest of the parameters were given by $(\rho, h, \phi) = (1, 1, 0.02)$. The mean number of points  is $\mathbb E N(L) = \rho|L| = 736$ for all the chosen parameter settings. For each simulation we fitted two models using minimum constrast estimation (using the same values of $a_1,a_2,q,p$ as in Section~\ref{s:den-an}): an ICP with an inverse gamma Bernstein CDF with $h$ and $\tau$ fixed at their true values, and an ICP with an exponential covariance function with $h=1$ fixed. 
For each simulation, we calculated the non-parametric estimate of the pcf and used $D$ in \eqref{e:contrast} for calculating its distance to each of the two pcfs for the estimated models; using an obvious notation, these distances are denoted $D_\text{Exp}$ and $D_\text{IG}$. For each combination of parameters $\sigma^2$ and $\tau$, Table~\ref{tab:comparison} shows the percentage of simulations where $D_\text{IG}>D_\text{Exp}$, i.e., the cases where the non-parametric estimate of the pcf resembles an exponential covariance function more than a covariance function with inverse gamma Bernstein CDF with $\tau$ equal to the value used in the simulations. To quantify how far the ICP is from a Poisson process, the values of $\varphi$ given by \eqref{e:varphiICP} and $p_\textrm{ms}$ given by \eqref{e:pms} are also shown in Table~\ref{tab:comparison}. The row in the table corresponding to $\sigma^2=10^{-1}$ contains values close to $50\%$, which makes sense since $g_0(0)-1=0.00416$ or $p_\mathrm{ms}=0.912$ reveal that the simulated processes are almost Poisson.
% and the two pair correlation functions are almost identical. 
In the rest of the table the non-parametric estimate are typically closest to an exponentiel covariance function, showing that there indeed is a tendency for getting an estimated model corresponding to an exponential covariance function (although the values seem to decrease in the last row). 

\begin{table}
	\centering
	\begin{tabular}{c|c|c|c|c|c|c}
		$\sigma^2$ & $\varphi$ & $p_\textrm{ms}$ & $\tau=1.1$ & $\tau=1.5$ & $\tau=2$ & $\tau=5$  \\
		\hline\hline
		$10^{-1}$  & 0.00416 & 0.912 & 56.6\% & 54.3\% & 53.3\% & 48.3\% \\
		$10^0$     & 0.155   & 0.577 & 81.1\% & 84.6\% & 86.8\% & 81.6\% \\
		$10^1$     & 1.40    & 0.218 & 74.3\% & 80.1\% & 85.8\% & 89.2\% \\
		$10^2$     & 6.12    & 0.0705& 76.1\% & 82.2\% & 82.6\% & 87.1\% \\
		$10^3$     & 21.4    & 0.0223& 82.5\% & 79.8\% & 79.0\% & 81.0\% \\
		$10^4$     & 69.7    &0.00707& 55.6\% & 60.1\% & 66.1\% & 77.3\% \\
	\end{tabular}
	\caption{The percentage of simulations of an ICP with $D_\text{IG}>D_\text{Exp}$ for various combinations of $\sigma^2$ and $\tau$ values, where the corresponding values of $\varphi$ and $p_\mathrm{ms}$ are shown.} 
	\label{tab:comparison}
\end{table}

We also made similar simulation studies for LGCPs and PCPPs. These showed similar values to those in Table~\ref{tab:comparison} for the LGCP, while the values where significantly smaller for the PCPP (typically in the range $30\%$-$40\%$ when $\varphi=2$ is as large as possible in this model).

The results of this section are further discussed in Section~\ref{s:7.3}.


%Since it seemed to be a problem to estimate the parameters in the inverse gamma Bernstein CDF, we also tried to fix the parameter $\sigma^2=5$ at its true value, and focused on estimating the two parameters $\phi$ and $\tau$ in this setup. We then repeated the above experiment to see whether the non-parametric estimate of the pcf was closest to the estimated exponential covariance function or the covariance function with inverse gamma Bernstein CDF with $\tau=2$, where both of the estimated pcfs had $\sigma^2=5$ fixed at its true value. The result showed that $79.8\%$ of the 1000 simulations was closest to the exponential case, thus showing no improvement. 



%MODIFICER OG FLYT TIL NAESTE AFSNIT: The findings of this simulation study suggest that the minimum contrast procedure based on the pair correlation function does not work well for estimating the parameters in the covariance functions given in Examples~\ref{ex:0}, while it works fine for the exponential covariance function. Minimum contrast estimation based on other summary statistics or different estimation procedures may work better, but we will leave this topic for future research.


\section{Concluding remarks}\label{s:conclusion}

%Our contributions can be summarized as follows. We have given an overview on and a comparison of the geodesic metric and the resistance metric on a linear network. %in Section~\ref{s:2} 
%This was needed for constructing isotropic covariance functions and GPs %in Section~\ref{s:GP-cov-fct} 
%which %in Section~\ref{s:cox}
% was transformed to random intensity functions for Poisson processes, resulting in new Cox point process models on linear networks with isotropic pair correlation functions. %In particular, we contributed with various new results:
% %In Section~\ref{s:res-met}, 
% Further, the resistance metric and the associated Jacobian in \eqref{e:jacobian} %$J_\delta(u,v)=|(\mathrm d/\mathrm dt) \delta(u,u_i(t))|$ 
% are quickly calculated from \eqref{e:d-R1}, \eqref{e:d-R2}, and \eqref{e:der-d-R} once the matrix $\Sigma$ has been calculated \citep[][do not specify any such results]{RakshitEtAl}. Thereby we can easily work with non-parametric estimates of the $K$-function associated to the resistance metric, cf.\ \eqref{e:K3}.
%  %In Section~\ref{s:sim-GP}, Theorem~\ref{t:simalgotree} gives a fast 
%Furthermore, Algorithm~\ref{a:2} gives a fast way for simulating a Gaussian process with an exponential covariance function when $L$ is a tree, and Theorem~\ref{t:3} extends this to an approximate simulation algorithm for arbitrary covariance functions with known Bernstein CDF; in relation to this, Example~\ref{ex:0} introduces three classes of covariance functions specified by the Bernstein CDF. 
% In particular, we have studied the model classes of log Gaussian Cox, interrupted Cox, and permanental Cox point processes in detail.
% Finally, we have discussed how minimum contrast estimation and global envelope tests for model checking work when these models are applied for real data analysis and simulation studies on linear networks. 
% Incidentally, when discussing model checking we corrected some considerations in \citet{CronieEtAl} in relation to the concept of IRMPS, cf.\ Sections~\ref{s:est-check} and \ref{s:LGCP}.
%%These new contributions give rise to several comments as follows.

Our results and considerations on point processes on linear networks in the previous sections give rise to various conclusions and open research problems which we discuss in Sections~\ref{s:7.1}--\ref{s:7.4} below.

\subsection{New point process models}\label{s:7.1}

\citet{Stationary} mentioned the lack of repulsive models on linear networks. An interesting case is a determinantal point process (DPP) which to the best of our knowledge has yet not been investigated in connection to linear networks. 
 For any given covariance function $c$ on $L$, a well-defined DPP will be specified by the density 
\begin{equation}\label{e:dppdens}
f(\{u_1,...,u_n\})\propto \mathrm{det}\left(\{c(u_i,u_j)\}_{i,j=1,...,n}\right)
\end{equation}
with respect to the unit rate Poisson process on $L$. A DPP is a model for repulsiveness (e.g.\ $g\le1$) and it {possesses} a number of appealing properties.
See
\citet{FredEtAl} (and the references therein) who considered DPPs with state space $S=\mathbb R^k$ but many things easily are modified to the case $S=L$. 

Several facts are interesting: If $c$ is isotropic then the pcf of the DPP is isotropic, and e.g.\ 
$c$ could be specified by one of 
the isotropic covariance functions given in Table~\ref{tab:covariances-sphere} and Example~\ref{ex:0}. 
The normalizing constant which is omitted in \eqref{e:dppdens} can be expressed in terms of the eigenvalues of a spectral decomposition of $c$. This spectral decomposition is also needed when specifying the $n$-th order intensity functions and an efficient simulation algorithm of the DPP.  However, it is an open problem how to determine the eigenvalues and eigenfunctions of the spectral representation, in particular if $c$ is required to be isotropic. Until this problem has been solved, we need to work with the unnormalized density (in the right hand side of \eqref{e:dppdens}) and to use MCMC methods for approximating the normalizing constant \citep{Geyer,MW04}.

The recent paper \citet{BSW} on a new class of Whittle–Mat{\'e}rn GPs on compact metric graphs is interesting for several reasons including the following. 
The model class is flexible and contains differentiable GPs; the precision matrix at the vertices can be quickly calculated; and
Markov properties of the GP can be used for computationally efficient inference. Thus the Whittle–Mat{\'e}rn GPs may 
 serve as an interesting alternative to the GPs considered in the present paper. However, we believe the Whittle–Mat{\'e}rn GPs  are less useful 
 when constructing Cox processes driven by transformed GPs because the
Whittle–Mat{\'e}rn covariance function is not isotropic (which is a drawback for several reasons, cf.\ items (I)--(V) in Section~\ref{s:intro}).   Specifically, the variance function of a Whittle–Mat{\'e}rn GP is not constant and \citet{BSW}
 used 
  the geodesic (rather than the resistance) metric. 

\subsection{The choice of metric for isotropic covariance and pair correlation functions}\label{s:7.2}

Isotropic covariance functions are available for larger classes of linear networks with respect to the resistance metric than the geodesic metric, cf.\ Theorem~\ref{t:1}. Therefore, we are able to obtain LGCPs, ICPs, and PCPPs for larger classes of linear networks with respect to the resistance metric than the geodesic metric.

Suppose $L$ is not a tree and that both $c_0(d_{\mathcal G}(u,v))$ and $c_0(d_{\mathcal R}(u,v))$ are well-defined isotropic covariance functions.
(Recall that $d_{\mathcal G}\ge d_{\mathcal R}$ with equality if and only if $L$ is a tree, cf.\ (C) in Theorem~\ref{t:res-m}.) 
Suppose also that $c_0> 1$ and $c_0$ is a decreasing function, so $c_0(d_{\mathcal G}(u,v))\le c_0(d_{\mathcal R}(u,v))$. This is 
typical for the covariance function of the GP underlying a LGCP, ICP, or PCPP
(an exception is a LGCP if $c_0$ can be negative but usually {$c_0\ge0$}). Consider the two LGCPs, ICPs, or PCPPs obtained by using the same type of transformed isotropic GP (or GPs) with the same mean and $c_0$-function but using the different metrics $d=d_{\mathcal G}$ and $d=d_{\mathcal R}$. Using an obvious notation, let us denote 
these two Cox processes by $X_{\mathcal G}$ and $X_{\mathcal R}$ and their corresponding pair correlation functions by 
$g_{\mathcal G}(u,v)=g_0(d_{\mathcal G}(u,v))$ and $g_{\mathcal R}(u,v)=g_0(d_{\mathcal R}(u,v))$. It follows from \eqref{e:pairLGCP}, \eqref{e:g0ICP}, and \eqref{e:pairPCPP} that   
 {$g_{\mathcal G}(u,v)\le g_{\mathcal R}(u,v)$}, indicating that $X_{\mathcal R}$ is able to model a higher degree of clustering than  
 $X_{\mathcal G}$.
 
% Thus, for two such Cox processes with the same decreasing $g_0$-function but given by the different metrics $d_{\mathcal G}$ and $d_{\mathcal R}$, respectively, we have that $g_0(d_{\mathcal R}(u,v))\ge g_0(d_{\mathcal G}(u,v))$. Consequently, the Cox process using the resistance metric is able to model a higher degree of clustering. Therefore, considering LGCPs, ICPs, and PCPPs with covariance functions for the underlying GPs given by the parametric classes of covariance functions in Table~\ref{tab:covariances-sphere} and Example~\ref{ex:0}, we see that using the resistance metric gives a larger degree of flexibility for modelling clustering
%
%ranging from the Poisson case ($g=1$) to 
%
%the pcf indicates more clustering for the model using the resistance metric. 

On the other hand, for a DPP, we have always that $g_0\le 1$, and typically $g_0< 1$ and $g_0$ is an increasing function. So for two such DPPs with the same $g_0$-function but given by the different metrics $d_{\mathcal G}$ and $d_{\mathcal R}$, respectively, the DPP using $d_{\mathcal R}$ is able to a higher degree of repulsiveness than the DPP using $d_{\mathcal G}$. 

In conclusion, using the resistance metric rather than the geodesic metric allows more well-defined LGCP, ICP, and PCPP models as well as more
flexibility for modelling clustering and repulsiveness. 

\subsection{Estimation}\label{s:7.3}

%{\color{red} Check om dette afsnit skal opdateres, naar alt er paa plads mht sim studiet i Sec 6.3.}

For parameter estimation in this paper we used minimum contrast estimation based on the pcf which worked well when we used an exponential covariance function, but this approach was not able to identify the parameters in the covariance functions given by the Bernstein CDFs in Example~\ref{ex:0} unless we fixed a subparameter, cf.\ Sections~\ref{s:den-an} and \ref{s:simstudy}. We are not sure if this is a problem of unidentifiability or caused by the choice of estimation procedure for the following reasons. 

On the one hand,
\cite{HeidiMe} concluded that for the dendrite data minimum contrast estimation based on the pcf performs better than if it is based on the $K$-function, and the alternative method based on  composite likelihood estimation is less reliable than minimum contrast estimation for the exponential covariance function, thus suggesting that this approach will not improve estimation for more complex covariance functions either.
This is also in line with the comments in \citet{BadEtAl2022} on estimation procedures. 
In particular, \citet[][Section 12]{BadEtAl2022} concluded that Neyman-Scott point process models are poorly identified under weak clustering, irrespective of the model parametrization, and this is not due to faults in currently available fitting
methods (they considered the Euclidean state space case but we do not expect the choice of space to be of importance for this conclusion).
%Moreover, for Neyman–Scott cluster process models, \citet{BadEtAl2022} showed that cluster scale is unidentifiable when the model is close to a Poisson process (they considered the Euclidean state space case but we do not expect the choice of space to be of importance). 

On the other hand, if the problem is unidentifiability, one idea for estimating the parameters of the Bernstein CDFs could be to  
to use shrinkage estimators involving a penalty on cluster scale as was done in \citet{BadEtAl2022} using the parameter $\varphi$ in \eqref{e:varphi}.
%we leave it as an open problem to investigate how to estimate the parameters of the Bernstein CDFs, where the idea in \citet{BadEtAl2022} of introducing an index of cluster strength and using shrinkage estimators involving a penalty on cluster scale could be a source of inspiration. 
%In our setting with an isotropic covariance function for the residual GP, the cluster strength parameter in \citet{BadEtAl2022} is given by $\varphi=g_0(0)-1$, which for LGCPs, ICPs, and PCPPs is equivalent to the variance $\sigma^2$, cf.\ \eqref{e:pairLGCP}, \eqref{e:g0ICP}, and \eqref{e:pairPCPP}: as $\sigma^2$ grows, realizations of the Cox process look more and more clustered.
However, we also remade the simulation studies behind Table~\ref{tab:comparison} in the case where we fixed the parameter $\sigma^2$ at its true value, which did not improve the estimation of the parameters in the Bernstein CDF, thus suggesting this approach might not work.

A more careful study of whether the problem is the choice of estimation procedure or unidentifiability, extending 
the methods of \citet{BadEtAl2022} to our setup, is left for future research. 
More ambitiously than minimum contrast, composite likelihood, and other moment based estimation procedures \citep{MW17}, future work may also investigate likelihood-based inference for parametric point process models on linear networks. This could include the adaptation of missing data MCMC methods for maximum likelihood \citep{Geyer,MW04} and Bayesian approaches \citep{RueEtAl,GT} to linear networks. 

\subsection{Miscellaneous}\label{s:7.4}
  
The special results in Section~\ref{s:LGCP} for LGCPs and any choice of the state space $S\in\{\mathbb R^k,\mathbb S^k,L\}$ could possibly be exploited for developing techniques of model fitting or checking, where we have the following results in mind and 
assume the covariance function of the underlying GP is isotropic. The $K$-functions based on $X$ and $X_u$ agree, so if $X$ is observed within a bounded window $W\subseteq S$, empirical estimators $\hat K$ and $\hat K_u$ for $u\in X_W$ are expected to be close. 
Moreover, the special structure in 
\eqref{e:n-LGCP} implies that 
\[\frac{\rho(u_1,u_2,u_3)}{\rho(u_1)\rho(u_2)\rho(u_3)}=g_0(d(u_1,u_2))g_0(d(u_1,u_3))g_0(d(u_2,u_3)).\]
This structure was exploited in \citet[][]{LGCP} to construct a third-order moment characteristic which is useful for model checking when $S=\mathbb R^k$, but it remains to consider the case $S\in\{\mathbb S^k,L\}$. %This can be extended to the more general setting considered in the present paper...

We have given various references for non-parametric estimation of the intensity, pair correlation, and $K$-functions on linear networks, cf.\ Sections~\ref{s:3.2} and \ref{s:3.2.3}. In the inhomogeneous case these estimators depend `locally' on the intensity function at the individual
observed points. Recently, when $S=\mathbb R^k$, \citet{ShawEtAl} demonstrated the advantages of introducing new global estimators over the
existing local estimators. Again it remains to consider the case $S\in\{\mathbb S^k,L\}$.



%Finally, let us mention the possibility of developing other kind of Cox processes than those considered in the present paper. The determinantal shot noise Cox processes introduced in  .... for the case $S=\mathbb R^k$ immediately extend to more general state spaces, including when $S\in\{\mathbb S^k,L\}$. These processes involve a  

%When $S=\mathbb R^k$, LGCPs have found numerous statistical applications and form perhaps the most popular subclass of Cox processes: In most statistical applications $c$ or equivalently $g$ is assumed to be isotropic, and as the likelihood function is intractable, minimum contrast or composite likelihood methods based on $\rho$ and $g$ are usually used for estimation, see \citet[][]{MW07,MW17} and the references therein. See also \citet{MT} for the case where $c$ is only assumed to be geometric anisotropic, cf.\ {\color{red} Section~\ref{s:2.1}}. Likelihood-based MCMC methods include missing data MCMC methods for maximum likelihood \citep{Geyer,MW04} and Bayesian approaches \citep{RueEtAl,GT}. These approaches can be extended to the case of linear networks as alternatives to minimum contrast estimation.










%HUSK: Determinantal shot noise Cox processes....

%HUSK:Cox-Strauss...

%HUSK: Mail fra Jesper: Har nu set på den artikel, som Ottmar sendte et link til: https://arxiv.org/pdf/2205.06163.pdf : Det er ret så langhåret. De får defineret en ny slags Gaussiske processer på en metrisk graf med en slags Matern (Whittle-Matern) kovariansfunktion: Det er alt sammen meget teknisk men ender op med at være hurtigt til beregninger i Markov tilfældet alpha=1,2,…, hvor de i praksis ser på alpha = 1 eller 2. Hurtigt fordi de så kender præcisionsmatricen for knuderne, hvor denne er sparse (mange nuller), og når de betinger med alle knuderne/endepunkterne (mere præcist betinger med de tilsvarende stokastiske variable) , så er der betinget uafhængighed mellem kanterne og den betingede fordeling på en kant afhænger kun af kantens endepunkter; og betinges med  kantens endepunkter, så stemmer kovariansfunktionen stemmer overens med en sædvanlig 1D Matern, når der tilsvarende betinges i denne. Altså en betinget uafhængighedsstruktur som sjovt nok er lig med den jeg havde i tankerne, da jeg forkert fik formuleret en sætning om simulation under en eksponentiel kovariansfunktion.
%Heldigvis betyder det ikke det store for vores arbejde: Deres ” Whittle-Matern”-kovariansfunktion er nemlig ikke isotrop (de betragter den geodætiske metrik og understreger selv at fx variansen ikke er konstant, se fx to af deres figurer). For alpha=1 (svarende til nu=1/2) er deres ” Whittle-Matern”-kovariansfunktion heller ikke eksponentiel – se deres simple eksempler, når grafen er et linjestykke eller en cirkel. Det skyldes, så vidt jeg forstår, at en randbetingelse for knuderne (fx a la’ Kirchhoff’s lov, se side 4; eller hvad der senere bruges af eksempler) er nødvendig, når de definer Whittle-Matern processen via (2) eller (7). Til gengæld har de så en vedefineret gaussisk proces for ethvert nu>0. 



%\section*{Address} Jesper M\o ller, Department of Mathematical Sciences, Aalborg University, Skjernvej 4A, DK-9220 Aalborg \O, Denmark. E-mail: jm@math.aau.dk

\appendix

%\section{Appendix}\label{a:A} %Details for Section~\ref{s:setting}
%
%\subsection{Surface measure on the unit sphere}\label{a:A1}
%
%For $k=1,2,\ldots$, the surface measure $\nu_k$ on %the the $k$-dimensional unit sphere
%$\mathbb{S}^k$ can be defined recursively as follows \citep[see e.g.][Chapter~1]{MR3060033}.
%For $k=1$ and
%$u=(\cos\theta,\sin\theta)$ with $0\le\theta<2\pi$, 
%$\mathrm d\nu_1(u)=\mathrm
%d\theta$ is the usual Lebesgue measure on $[0,2\pi)$. For $k\ge2$ and
% $u=(v\sin\vartheta,\cos\vartheta)$ with $v\in\mathbb S^{k-1}$ and
%  $\vartheta\in[0,\pi]$, 
%\begin{equation}\label{e:a1}
%\mathrm d\nu_k(u)=\sin^{k-1}\vartheta\,\mathrm d\nu_{k-1}(v)\,\mathrm d\vartheta.
%\end{equation}
%In particular, 
%% instance, if $d=1$ and
%% $\bx=(x_1,x_2)=(\cos\theta,\sin\theta)$ with $0\le\theta<2\pi$, then
%% $\mathrm d\nu_1(\bx)=\mathrm
%% d\theta$, the usual Lebesgue measure on $[0,2\pi)$; and
%if $k=2$ and 
%$u=(\sin\vartheta\cos\varphi,\sin\vartheta\sin\varphi,\cos\vartheta)$
%where $\vartheta\in[0,\pi]$ is the polar latitude 
%and $\varphi\in[0,2\pi)$ is the polar longitude, 
%$\mathrm d\nu_2(u)=\sin\vartheta\,\mathrm d\varphi\,\mathrm d\vartheta$. % For a general description of $\nu_d$, see e.g.\ 
%% \cite[Chapter~1]{MR3060033}.
%Note that $\mathbb{S}^k$ has 
%surface measure
%$\sigma_k=\nu_k(\mathbb{S}^k)=\frac{2\pi^{(k+1)/2}}{\Gamma((k+1)/2)}$ 
%% $\sigma_d=\nu(\mathbb{S}^d)=2\pi^{(d+1)/2}/\Gamma((d+1)/2)$ 
% ($\sigma_1=2\pi$, $\sigma_2=4\pi$).
% 
%\subsection{Resistance metric on a linear network}\label{a:A2} 

%%This appendix describes the resistance metric for a graph with Euclidean edges \citep{AnderesEtAl} in the special case of
%% a linear network $L=\cup_{i=1}^m L_i$ as given in Section~\ref{s:2.3}. Thereby a graph $G$ with nodes $V$ given by the endpoints and edges given by $L_1,\ldots,L_m$ is obtained. 
%% 
%% Since the resistance metric on $L$ is an extension of the classic (effective) resistance metric $d_V$ defined on $V$, I start by recalling what this is:
%% As in Section~\ref{s:2.3}, let 
%%$a_i$ and $b_i$ denote the endpoints of $L_i$, $i=1,\ldots,m$, so $V=\cup_{i=1}^m\{a_i,b_i\}$.
%%Let $u_0\in V$ be an arbitrarily chosen vertex called the origin.
%%Consider any $u,v\in V$ and write $u\sim v$ if $u,v\in L_i$ for some $i$. Define the so-called conductance function by
%%\[{\mathrm{con}}(u,v)=\begin{cases*}
%%                    1/\|u-v\| & if  $u\sim v$  \\
%%                     0 & otherwise
%%                 \end{cases*}
%%                 \] 
%%and define a matrix $\Delta$ with rows and columns indexed by $V$ so that its entry $(u,v)$ is given by
%%\[\Delta(u,v)=\begin{cases*}
%%                    1+c(u) & if  $u=v=u_0$  \\
%%                     c(u) & if $u=v\not=u_0$\\
%%                     -{\mathrm{con}}(u,v) & otherwise
%%                 \end{cases*}
%%                 \] 
%%where $c(u)=\sum_{w\in V:\,w\sim u}{\mathrm{con}}(u,w)$ is the sum of the conductances associated to the edges incident to vertex $u$.
%%In fact $\Delta$ is symmetric and strictly positive definite, and it is similar to %what would be called 
%%the `Laplacian matrix'
%%from electrical network theory \citep[when viewing $G$ as an electrical network over the nodes with resistors given by the length of each line segment, see e.g.][]{Kigami,JorgensenPearse} 
%%except that $D$ has the additional $1$ added at entry $(u_0,u_0)$ which makes $\Delta$ invertible. Let $B_0$ be a Gaussian vector indexed by $V$ and having zero mean and covariance matrix $\Sigma=\Delta^{-1}$. Then
%%the resistance metric on $V$ is the variogram
%%\[d_V(u,v)=\mathbb{V}{\mathrm{ar}}(B_0(u)-B_0(v))=\Sigma(u,u)+\Sigma(v,v)-2\Sigma(u,v)\quad\mbox{for } u,v\in V.\]
%%It can be shown that %the definition of 
%%$d_V$ does not depend on the choice of origin $u_0$.
%%
%%Now, extend $B_0$ by linear interpolation to a zero mean Gaussian random field (GRF) $Z_0$ on $L$ so that %$Z_0(u)=B_0(u)$ if $u\in V$ and
%%\[Z_0(u)=\frac{\|u-b_i\|}{l_i}B_0(a_i)+\frac{\|u-a_i\|}{l_i}B_0(b_i)\quad\mbox{for }u\in L_i.\]
%%For $i=1,\ldots,m$, define a zero mean Brownian bridge $B_i$ on $L_i$ so that  %$B_i(a_i)=B_i(b_i)=0$ and
%%\[\mathbb{C}{\mathrm{ov}}(B_i(u),B_i(v))=\min\{\|u-a_i\|\|v-b_i\|,\|v-a_i\|\|u-b_i\|\}/l_i\quad\mbox{for $u,v\in L_i$,}
%%%\begin{cases*}
%%%\|u-a_i\|\|v-b_i\|/l_i & if $\|u-a_i\|\le\|v-a_i\|$ \\
%%%\|v-a_i\|\|u-b_i\|/l_i & if $\|v-a_i\|\le\|u-a_i\|$
%%%\end{cases*}
%%%%\left[\min\{d_i(u),d_i(v)\}-d_i(u)d_i(v)\right]l_i,\quad u,v\in L_i.
%%\]
%%%as follows. 
%%%First, $Z_0(u)=B_0(u)$ if $u\in V$. Second, 
%%% $Z_0$ is defined on each line segment $L_i$ by linear interpolation so that 
%%%% if 
%%%%\[d_i(u)= 
%%%%\begin{cases*}
%%%%                    \|u-a_i\|/l_i & for  $u\in L_i\setminus\{a_i,b_i\}$  \\
%%%%                     0 & for $u\in\{a_i,b_i\}$
%%%%                 \end{cases*}
%%%%                 \]
%%%%is the distance of $u\in L_i$ as a proportion of the length
%%%%of $L_i$, then
%%%%\[Z_0(u)=\left(1-d_i(u)\right)Z_0(a_i)+d_i(u)Z_0(b_i).\]
%%%\[Z_0(u)=\frac{\|u-b_i\|}{l_i}Z_0(a_i)+\frac{\|u-a_i\|}{l_i}Z_0(b_i)\quad\mbox{if }u\in L_i.\]
%%and define
%%\[Z_i(u)=\begin{cases*}
%%                    B_i(u) & for  $u\in L_i$  \\
%%                     0 & for $u\in L\setminus L_i$.
%%                 \end{cases*}
%%                 \]
%%%Finally, % assuming $B_0,B_1,\ldots,B_m$ are independent, 
%%%define 
%%%\begin{equation*}
%%%    Z(u)=\sum_{i=0}^m\!Z_i(u),\quad u\in L.
%%%\end{equation*}
%%%Thus $Z$ is a sum of $m+1$ independent, zero mean, GRFs. 
%%The resistance metric on $L$ is defined by %is defined by
%% \[d_{\mathcal{R}}(u,v) %={\mathrm{var}}(Z(u)-Z(v))
%% =\sum_{i=0}^m \mathbb{V}{\mathrm{ar}}(Z_i(u)-Z_i(v))\quad\mbox{for } u,v\in L.\]
%%Clearly, $d_{\mathcal{R}}(u,v)=d_V(u,v)$ if $u,v\in V$. 
%% 
%%For $u\in L_j$ and $v\in L_i$,  a straightforward calculation gives
%%\begin{align}
%%d_{\mathcal{R}}(u,v) %{\mathrm{var}}(Z_0(u)-Z_0(v))
%%=\,&
%%\frac{\|u-b_j\|^2}{l_j^2}\Sigma(a_j,a_j)+\frac{\|u-a_j\|^2}{l_j^2}\Sigma(b_j,b_j)+2\frac{\|u-b_j\|\|u-a_j\|}{l_j^2}\Sigma(a_j,b_j)\nonumber\\
%%&+\frac{\|v-b_i\|^2}{l_i^2}\Sigma(a_i,a_i)+\frac{\|v-a_i\|^2}{l_i^2}\Sigma(b_i,b_i)+2\frac{\|v-b_i\|\|v-a_i\|}{l_i^2}\Sigma(a_i,b_i)\nonumber\\
%%&-2\frac{\|u-b_j\|\|v-b_i\|}{l_jl_i}\Sigma(a_j,a_i)
%%-2\frac{\|u-b_j\|\|v-a_i\|}{l_jl_i}\Sigma(a_j,b_i)\nonumber\\
%%&-2\frac{\|u-a_j\|\|v-b_i\|}{l_jl_i}\Sigma(b_j,a_i)
%%-2\frac{\|u-a_j\|\|v-a_i\|}{l_jl_i}\Sigma(b_j,b_i)\nonumber\\
%%&+\frac{\|v-a_i\|\|v-b_i\|}{l_i}+\frac{\|u-a_j\|\|u-b_j\|}{l_j}\nonumber\\
%%&-2\frac{1(i=j)}{l_i}\min\{\|u-a_i\|\|v-b_i\|,\|v-a_i\|\|u-b_i\|\}. \label{e:d-R}
%%\end{align}
%%% \citet[][Proposition 4]{AnderesEtAl} states that $d_{\mathcal{R}}\le d_{\mathcal{G}}$
%%%with equality if and only if $G$ is a tree. For example, $d_{\mathcal{R}}(u,v)=d_{\mathcal{G}}(u,v)-d_{\mathcal{G}}(u,v)^2/\sum_{i=1}^ml_i$ if $G$ is a cycle. 
%
%
%% Further, for $h=1,\ldots,m$, it follows immediately that
%%\begin{align*}{\mathrm{var}}&(Z_h(u)-Z_h(v))=\\
%%&\frac{1}{l_h}\times\begin{cases*}
%%\|u-a_h\|\|u-b_h\|+\|v-a_h\|\|v-b_h\|-2\min\{\|u-a_h\|\|v-b_h\|,\|v-a_h\|\|u-b_h\|\} & if $h=i=j$\\
%%\|v-a_h\|\|v-b_h\| & if $h=i\not=j$\\
%%\|u-a_h\|\|u-b_h\| & if $h=j\not=i$\\
%%0 & otherwise.
%%\end{cases*}
%%\end{align*}

%For the weight $w_\delta$ in \eqref{e:K3}, which also appears in non-parametric estimators of $K_L$, the calculation of $(\mathrm d/\mathrm dt) \delta(u,u_i(t))$
%%the Jacobian $J_\delta(u,v)=|(\mathrm d/\mathrm dt) \delta(u,u_i(t))|$ %|(\partial/\partial v)\delta(u,v)|$ 
%is needed %$(\partial/\partial v)\delta(u,v)$ is needed 
%for all $u\in L$ and all $v=u_i(t)\in L_i\setminus N$ with $i=1,\ldots,m$, where $N\subset L$ is finite. Let $\delta=d_{\mathcal{R}}$ and $N=V$, so  $v=u_i(t)=(1-t/l_i)a_i+(t/l_i)b_i$ with $0<t<l_i$. Then $\|v-a_i\|=t$, $\|v-b_i\|=l_i-t$, 
%\[ %\frac{\partial}{\partial v}\|v-b_i\|^2=
%\frac{\mathrm d}{\mathrm dt}\|u_i(t)-b_i\|^2=2(t-l_i)=-2\|v-b_i\|,\]
%and
%\[ %\frac{\partial}{\partial v}\|v-b_i\|=
%\frac{\mathrm d}{\mathrm dt}\|u_i(t)-b_i\|=-1. \]
%Therefore, 
%\begin{align}
%%\frac{\partial}{\partial v}d_{\mathcal{R}}(u,v)=&\,\frac{\mathrm d}{\mathrm d t}d_{\mathcal{R}}(u,u_i(t))\nonumber\\
%\frac{\mathrm d}{\mathrm d t}d_{\mathcal{R}}(u,u_i(t))=&\,\frac{\mathrm d}{\mathrm d t}\Bigg\{
%\frac{(l_i-t)^2}{l_i^2}\Sigma(a_i,a_i)+\frac{t^2}{l_i^2}\Sigma(b_i,b_i)+\frac{2t(l_i-t)}{l_i^2}\Sigma(a_i,b_i)\nonumber\\
%&-\frac{2\|u-b_j\|(l_i-t)}{l_jl_i}\Sigma(a_j,a_i)-\frac{2\|u-b_j\|t}{l_jl_i}\Sigma(a_j,b_i)\nonumber\\
%&-\frac{2\|u-a_j\|(l_i-t)}{l_jl_i}\Sigma(b_j,a_i)-\frac{2\|u-a_j\|t}{l_jl_i}\Sigma(b_j,b_i)\nonumber\\
%&+\frac{t(l_i-t)}{l_i}-2\frac{1(i=j)}{l_i}\min\{\|u-a_i\|(l_i-t),\|u-b_i\|t\}
%%\frac{1(i=j)}{l_i}\left[t(l_i-t)-2\min\{\|u-a_i\|(l_i-t),\|u-b_i\|t\}\right]+\frac{1(i\not=j)}{l_i}t(l_i-t)
%\Bigg\}\nonumber\\
%=&\,-2\frac{\|v-b_i\|}{l_i^2}\Sigma(a_i,a_i)+2\frac{\|v-a_i\|}{l_i^2}\Sigma(b_i,b_i)+2\frac{\|v-b_i\|-\|v-a_i\|}{l_i^2}\Sigma(a_i,b_i)\nonumber\\
%&+2\frac{\|u-b_j\|}{l_jl_i}\Sigma(a_j,a_i)-2\frac{\|u-b_j\|}{l_jl_i}\Sigma(a_j,b_i)
%+2\frac{\|u-a_j\|}{l_jl_i}\Sigma(b_j,a_i)-2\frac{\|u-a_j\|}{l_jl_i}\Sigma(b_j,b_i)\nonumber\\
%&+\frac{\|v-b_i\|-\|v-a_i\|}{l_i}+2\frac{\|u-a_i\|}{l_i}1(i=j,\,\|u-a_i\|\|\|v-b_i\|<\|u-b_i\|\|v-a_i\|)\nonumber\\
%&-2\frac{\|u-b_i\|}{l_i}1(i=j,\,\|u-a_i\|\|\|v-b_i\|>\|u-b_i\|\|v-a_i\|). \label{e:der-d-R} %-\frac{1(i\not=j)}{l_i}\|v-a_i\|.
%\end{align}
%
%Once the matrix $\Sigma$ has been calculated, the resistance metric and the Jacobian $J_\delta(u,v)=|(\mathrm d/\mathrm dt) \delta(u,u_i(t))|$ are quickly calculated from \eqref{e:d-R} and \eqref{e:der-d-R} \citep[][do not specify any such results]{RakshitEtAl}. 

%DOES IT FOLLOW THAT THIS GIVES THAT $J_{d_{\mathcal R}}$ IS AT MOST 1????

\bibliography{references}

\end{document} 







\begin{thebibliography}{49}
\expandafter\ifx\csname natexlab\endcsname\relax\def\natexlab#1{#1}\fi
\expandafter\ifx\csname url\endcsname\relax
  \def\url#1{\texttt{#1}}\fi
\expandafter\ifx\csname urlprefix\endcsname\relax\def\urlprefix{URL: }\fi

\bibitem[{Anderes et~al.(2020)Anderes, M{\o}ller and Rasmussen}]{AnderesEtAl}
Anderes, E., M{\o}ller, J. and Rasmussen, J.~G. (2020) Isotropic covariance
  functions on graphs and their edges.
\newblock \textit{Annals of Statistics}, \textbf{48}, 2478--2503.

\bibitem[{Ang et~al.(2012)Ang, Baddeley and Nair}]{AngEtAl}
Ang, Q., Baddeley, A. and Nair, G. (2012) Geometrically corrected second order
  analysis of events on a linear network, with applications to ecology and
  criminology.
\newblock \textit{Scandinavian Journal of Statistics}, \textbf{39}, 591--617.

\bibitem[{Baddeley et~al.(2022)Baddeley, Davies, Hazelton, Rakshit and
  Turner}]{BadEtAl2022}
Baddeley, A., Davies, T.~M., Hazelton, M.~L., Rakshit, S. and Turner, R. (2022)
  Fundamental problems in fitting spatial cluster process models.
\newblock \textit{Spatial Statistics}, \textbf{52}, 100709.

\bibitem[{Baddeley et~al.(2000)Baddeley, M{\o}ller and Waagepetersen}]{BMW}
Baddeley, A., M{\o}ller, J. and Waagepetersen, R. (2000) Non- and
  semi-parametric estimation of interaction in inhomogeneous point patterns.
\newblock \textit{Statistica Neerlandica}, \textbf{54}, 329--350.

\bibitem[{Baddeley et~al.(2017)Baddeley, Nair, Rakshit and
  McSwiggan}]{Stationary}
Baddeley, A., Nair, G., Rakshit, S. and McSwiggan, G. (2017) "{S}tationary"
  point processes are uncommon on linear networks.
\newblock \textit{Stat}, \textbf{6}, 68--78.

\bibitem[{Baddeley et~al.(2015)Baddeley, Rubak and Turner}]{spatstat}
Baddeley, A., Rubak, E. and Turner, R. (2015) \textit{Spatial Point Patterns:
  Methodology and Applications with {R}}.
\newblock Boca Raton: Chapman and Hall/CRC Press.

\bibitem[{Bolin et~al.(2022)Bolin, Simas and Wallin}]{BSW}
Bolin, D., Simas, A.~B. and Wallin, J. (2022) Gaussian {W}hittle-{M}at{\'e}rn
  fields on metric graphs.
\newblock Available at arXiv:2205.06163.

\bibitem[{Christensen and M{\o}ller(2020)}]{HeidiMe}
Christensen, H.~S. and M{\o}ller, J. (2020) Modelling spine locations on
  dendrite trees using inhomogeneous {C}ox point processes.
\newblock \textit{Spatial Statistics}, \textbf{39}, 100478.

\bibitem[{Coeurjolly et~al.(2017)Coeurjolly, M{\o}ller and
  Waagepetersen}]{CoeurjollyEtAl}
Coeurjolly, J.-F., M{\o}ller, J. and Waagepetersen, R. (2017) Palm
  distributions for log {G}aussian {C}ox processes.
\newblock \textit{Scandinavian Journal of Statistics}, \textbf{44}, 192--203.

\bibitem[{Cox(1955)}]{Cox55}
Cox, D.~R. (1955) Some statistical methods connected with series of events.
\newblock \textit{Journal of the Royal Statistical Society: Series B
  (Methodological)}, \textbf{17}, 129--157.

\bibitem[{Cronie et~al.(2020)Cronie, Moradi and Mateu}]{CronieEtAl}
Cronie, O., Moradi, M. and Mateu, J. (2020) Inhomogeneous higher-order summary
  statistics for point processes on linear networks.
\newblock \textit{Statistics and Computing}, \textbf{30}, 1221–1239.

\bibitem[{Cuevas-Pacheco and M{\o}ller(2018)}]{CM}
Cuevas-Pacheco, F. and M{\o}ller, J. (2018) Log {G}aussian {C}ox processes on
  the sphere.
\newblock \textit{Spatial Statistics}, \textbf{26}, 69--82.

\bibitem[{Dai and Xu(2013)}]{MR3060033}
Dai, F. and Xu, Y. (2013) \textit{Approximation Theory and Harmonic Analysis on
  Spheres and Balls}.
\newblock Springer Monographs in Mathematics. Springer, New York.

\bibitem[{Daley and Vere-Jones(2003)}]{DaleyVere-Jones}
Daley, D.~J. and Vere-Jones, D. (2003) \textit{An Introduction to the Theory of
  Point Processes. Volume I: Elementary Theory and Methods}.
\newblock second edn. Springer, New York.

\bibitem[{Diggle(2014)}]{Diggle}
Diggle, P. (2014) \textit{Statistical Analysis of Spatial and Spatio-Temporal
  Point Patterns. 3rd ed}.
\newblock Boca Raton, FL: CRC Press, Taylor and Francis Group.

\bibitem[{Geyer(1999)}]{Geyer}
Geyer, C.~J. (1999) Likelihood inference for spatial point processes.
\newblock In \textit{Stochastic Geometry: Likelihood and Computation} (eds.
  O.~E. Barndorff-Nielsen, W.~S. Kendall and M.~N.~M. {Van Lieshout}), 79--140.
  Chapman and Hall/CRC, Boca Raton, Florida.

\bibitem[{Gneiting(2001)}]{Gneiting01}
Gneiting, T. (2001) Criteria of {P}{\'o}lya type for radial positive definite
  functions.
\newblock \textit{Proceedings of American Mathematical Society}, \textbf{129},
  2309--2318.

\bibitem[{Gneiting(2013)}]{gneiting2013strictly}
--- (2013) Strictly and non-strictly positive definite functions on spheres.
\newblock \textit{Bernoulli}, \textbf{19}, 1327–1349.

\bibitem[{Guttorp and Thorarinsdottir(2012)}]{GT}
Guttorp, P. and Thorarinsdottir, T.~L. (2012) Bayesian inference for
  non-{M}arkovian point processes.
\newblock In \textit{Advances and Challenges in Space-time Modelling of Natural
  Events} (eds. E.~Porcu, J.~Montero and M.~Schlather), vol. 207 of
  \textit{Lecture Notes in Statistics}, 79--102. Springer: Berlin Heidelberg.

\bibitem[{Jorgensen and Pearse(2010)}]{JorgensenPearse}
Jorgensen, P. E.~T. and Pearse, E. P.~J. (2010) A {H}ilbert space approach to
  effective resistance metric.
\newblock \textit{Complex Analysis and Operator Theory}, \textbf{4},
  975–1013.

\bibitem[{Kigami(2003)}]{Kigami}
Kigami, J. (2003) Harmonic analysis for resistance forms.
\newblock \textit{Journal of Functional Analysis}, \textbf{204}, 399--444.

\bibitem[{Klein and M.~Randi{\'c}(1993)}]{KleinRandic}
Klein, D.~J. and M.~Randi{\'c}, M. (1993) Resistance distance.
\newblock \textit{Journal of Mathematical Chemistry}, \textbf{12}, 81--95.

\bibitem[{Lavancier and M{\o}ller(2016)}]{FredMe}
Lavancier, F. and M{\o}ller, J. (2016) Modelling aggregation on the large scale
  and regularity on the small scale in spatial point pattern datasets.
\newblock \textit{Scandinavian Journal of Statistics}, \textbf{43}, 587--609.

\bibitem[{Lavancier et~al.(2015)Lavancier, M{\o}ller and Rubak}]{FredEtAl}
Lavancier, F., M{\o}ller, J. and Rubak, E. (2015) Determinantal point process
  models and statistical inference.
\newblock \textit{Journal of the Royal Statistical Society: Series B
  (Methodological)}, \textbf{77}, 853--877.

\bibitem[{Lawrence et~al.(2016)Lawrence, Baddeley, Milne and
  Nair}]{LawrenceEtAl}
Lawrence, T., Baddeley, A., Milne, R.~K. and Nair, G. (2016) Point pattern
  analysis on a region of a sphere.
\newblock \textit{Stat}, \textbf{5}, 144–157.

\bibitem[{Macchi(1975)}]{Macchi:75}
Macchi, O. (1975) The coincidence approach to stochastic point processes.
\newblock \textit{Advances in Applied Probability}, \textbf{7}, 83--122.

\bibitem[{McCullagh and M{\o}ller(2006)}]{MacCM}
McCullagh, P. and M{\o}ller, J. (2006) The permanental process.
\newblock \textit{Advances in Applied Probability}, \textbf{38}, 873--888.

\bibitem[{Minc(1978)}]{Minc}
Minc, H. (1978) \textit{Permanents}.
\newblock Reading, MA: Addison-Wesley.

\bibitem[{M{\o}ller and Rubak(2016)}]{MollerRubak}
M{\o}ller, J. and Rubak, E. (2016) Functional summary statistics on the sphere
  with an application to determinantal point processes.
\newblock \textit{Spatial Statistics}, \textbf{18}, 4--23.

\bibitem[{M{\o}ller et~al.(1998)M{\o}ller, Syversveen and Waagepetersen}]{LGCP}
M{\o}ller, J., Syversveen, A.~R. and Waagepetersen, R.~P. (1998) Log {G}aussian
  {C}ox processes.
\newblock \textit{Scandinavian Journal of Statistics}, \textbf{25}, 451--482.

\bibitem[{M{\o}ller and Waagepetersen(2007)}]{MW07}
M{\o}ller, J. and Waagepetersen, R. (2007) Modern statistics for spatial point
  processes (with discussion).
\newblock \textit{Scandinavian Journal of Statistics}, \textbf{34}, 643--711.

\bibitem[{M{\o}ller and Waagepetersen(2017)}]{MW17}
--- (2017) Some recent developments in statistics for spatial point patterns.
\newblock \textit{Annual Review of Statistics and Its Applications},
  \textbf{4}, 317--342.

\bibitem[{M{\o}ller and Waagepetersen(2004)}]{MW04}
M{\o}ller, J. and Waagepetersen, R.~P. (2004) \textit{Statistical Inference and
  Simulation for Spatial Point Processes}.
\newblock Boca Raton: Chapman and Hall/CRC.

\bibitem[{Mrkvi{\v{c}}ka et~al.(2020)Mrkvi{\v{c}}ka, Myllym{\"a}ki, Jilik and
  Hahn}]{Mylly2}
Mrkvi{\v{c}}ka, T., Myllym{\"a}ki, M., Jilik, M. and Hahn, U. (2020) A one-way
  {ANOVA} test for functional data with graphical interpretation.
\newblock \textit{Kybernetika}, \textbf{56}, 432--458.

\bibitem[{Myllym{\"a}ki and Mrkvi{\v{c}}ka(2019)}]{Mylly3}
Myllym{\"a}ki, M. and Mrkvi{\v{c}}ka, T. (2019) {GET}: Global envelopes in {R}.
\newblock ArXiv preprint arXiv:1911.06583.

\bibitem[{Myllym{\"a}ki et~al.(2017)Myllym{\"a}ki, Mrkvi{\v{c}}ka, Grabarnik,
  Seijo and Hahn}]{Mylly1}
Myllym{\"a}ki, M., Mrkvi{\v{c}}ka, T., Grabarnik, P., Seijo, H. and Hahn, U.
  (2017) Global envelope tests for spatial processes.
\newblock \textit{Journal of the Royal Statistical Society: Series B
  (Statistical Methodology)}, \textbf{79}, 381–404.

\bibitem[{Pitt(1971)}]{Pitt71}
Pitt, L.~D. (1971) A {Markov} property for {G}aussian processes with a
  multidimensional parameter.
\newblock \textit{Archive for Rational Mechanics and Analysis}, \textbf{43},
  367--391.

\bibitem[{Rakshit et~al.(2019)Rakshit, Davies, Moradi, McSwiggan, Nair, Mateu
  and Baddeley}]{RakshitEtAl19}
Rakshit, S., Davies, T., Moradi, M., McSwiggan, G., Nair, G., Mateu, J. and
  Baddeley, A. (2019) Fast kernel smoothing of point patterns on a large
  network using two-dimensional convolution.
\newblock \textit{International Statistical Review}, \textbf{87}, 531–556.

\bibitem[{Rakshit et~al.(2017)Rakshit, Nair and Baddeley}]{RakshitEtAl}
Rakshit, S., Nair, G. and Baddeley, A. (2017) Second-order analysis of point
  patterns on a network using any distance metric.
\newblock \textit{Spatial Statistics}, \textbf{22}, 129--154.

\bibitem[{Rasmussen and Christensen(2021)}]{RasmussenChristensen}
Rasmussen, J.~G. and Christensen, H.~S. (2021) Point processes on directed
  linear networks.
\newblock \textit{Methodology and Computing in Applied Probability},
  \textbf{23}, 647--667.

\bibitem[{Rue and Held(2005)}]{rue:held}
Rue, H. and Held, L. (2005) \textit{Theory of Gaussian Markov Random Fields}.
\newblock Boca Raton: Chapman and Hall/CRC Press.

\bibitem[{Rue et~al.(2009)Rue, Martino and Chopin}]{RueEtAl}
Rue, H., Martino, S. and Chopin, N. (2009) Approximate {B}ayesian inference for
  latent {G}aussian models by using integrated nested {L}aplace approximations.
\newblock \textit{Journal of the Royal Statistical Society: Series B
  (Methodological)}, \textbf{71}, 319–392.

\bibitem[{Schoenberg(1938)}]{Schoenberg38}
Schoenberg, I. (1938) Metric spaces and completely monotone functions.
\newblock \textit{Annals of Mathematics}, \textbf{39}, 811–841.

\bibitem[{Schoenberg(1942)}]{Schoenberg}
--- (1942) Positive definite functions on spheres.
\newblock \textit{Duke Mathematical Journal}, \textbf{9}, 96–108.

\bibitem[{Shaw et~al.(2021)Shaw, M{\o}ller and Waagepetersen}]{ShawEtAl}
Shaw, T., M{\o}ller, J. and Waagepetersen, R.~P. (2021) Globally
  intensity-reweighted estimators for {K}- and pair correlation functions.
\newblock \textit{Australian and New Zealand Journal of Statistics},
  \textbf{63}, 93--118.

\bibitem[{Shirai and Takahashi(2003)}]{ShiraiTakahashi}
Shirai, T. and Takahashi, Y. (2003) Random point fields associated with certain
  {F}redholm determinants i: fermion, {P}oisson and boson point processes.
\newblock \textit{Journal of Functional Analysis}, \textbf{205}, 414–463.

\bibitem[{Stoyan(1979)}]{Stoyan}
Stoyan, D. (1979) Interrupted point processes.
\newblock \textit{Biometrical Journal}, \textbf{21}, 607--610.

\bibitem[{Valiant(1979)}]{Valiant}
Valiant, L.~G. (1979) The complexity of computing the permanent.
\newblock \textit{Theoretical Computer Science}, \textbf{8}, 189–201.

\bibitem[{{Van Lieshout}(2011)}]{Inhom-J-fct}
{Van Lieshout}, M. N.~M. (2011) A {J}-function for inhomogeneous point
  processes.
\newblock \textit{Statistica Neerlandica}, \textbf{65}, 183–201.

\end{thebibliography}
